% Two editions from one source. The Kindle edition is selected by defining
% \KindleEdition before the class is read, which the Makefile does on the
% command line:  pdflatex "\def\KindleEdition{}% Two editions from one source. The Kindle edition is selected by defining
% \KindleEdition before the class is read, which the Makefile does on the
% command line:  pdflatex "\def\KindleEdition{}% Two editions from one source. The Kindle edition is selected by defining
% \KindleEdition before the class is read, which the Makefile does on the
% command line:  pdflatex "\def\KindleEdition{}% Two editions from one source. The Kindle edition is selected by defining
% \KindleEdition before the class is read, which the Makefile does on the
% command line:  pdflatex "\def\KindleEdition{}\input{main.tex}"
\ifdefined\KindleEdition
  \documentclass[10pt,oneside]{book}
\else
  \documentclass[11pt,a4paper,twoside,openright]{book}
\fi
\input{preamble}
\renewcommand*{\thisvolume}{mixed}

\title{\Huge Mixed Problems\\[0.4em]
  \large Cross-Chapter Sets with Worked Solutions}
\author{}
\date{}

\begin{document}
\frontmatter
\maketitle

\chapter*{About this sheet}
\addcontentsline{toc}{chapter}{About this sheet}

Every other volume in this series drills one technique at a time. That is
the right way to learn a method and the wrong way to prepare for an
examination, because a chapter heading tells you what to do before you have
read the question. An exam does not. The single most common way to lose
marks is not botched arithmetic but reaching for the wrong tool: setting up
a Lagrangian where the constraint does not bind, differentiating where a
limit was wanted, solving a system by inverting a matrix that turns out to
be singular.

This sheet exists to remove that crutch. Each set draws on the whole
sequence --- algebra and functions, matrices and elimination, limits and
L'H\^opital, curve discussion, multivariable calculus, constrained
optimization, linear programming --- in no announced order, and no problem
says which chapter it came from. Deciding that \emph{is} the exercise.

Work a set in one sitting, with a time limit and without looking anything
up. Sixty to ninety minutes for twenty problems is a realistic target. Then
read the solutions, and for every problem you got wrong, note not the
arithmetic but the \emph{decision}: which signal in the question should have
pointed at which method? That note is worth more than redoing the
computation.

Solutions name the technique in their first words, so the sets can also be
used the other way round --- read only the first line of each solution as a
drill in classification.

The fourth set works the other way. It announces the methods, all of them,
and takes a single function of two variables through the whole chain:
domain, range, derivative, level curves, gradient, Hessian, critical
points, contour sets and a checklist of properties. There the difficulty is
not choosing a tool but keeping a dozen answers consistent with one
another.

\tableofcontents

\mainmatter
\include{chapters/ch01}
\include{chapters/ch02}
\include{chapters/ch03}
\include{chapters/ch04}

\end{document}
"
\ifdefined\KindleEdition
  \documentclass[10pt,oneside]{book}
\else
  \documentclass[11pt,a4paper,twoside,openright]{book}
\fi
% Shared preamble for the Mixed Problems training sheet.
\usepackage[T1]{fontenc}
\usepackage[utf8]{inputenc}
\usepackage{lmodern}
\usepackage[a4paper,margin=2.4cm]{geometry}
\usepackage{amsmath,amssymb,amsthm}
\usepackage{microtype}
\usepackage{enumitem}
\usepackage{xcolor}
\usepackage[most]{tcolorbox}
\usepackage{fancyhdr}
\usepackage[hidelinks]{hyperref}

\definecolor{accent}{HTML}{1F4E79}
\definecolor{soft}{HTML}{F4F7FA}

% --- function graphs ------------------------------------------------------
\usepackage{pgfplots}
\pgfplotsset{compat=1.18}
\usetikzlibrary{arrows.meta}

% One consistent look for every function graph in these books.
\pgfplotsset{
  fnaxis/.style={
    width=5.9cm, height=4.6cm,
    axis lines=middle,
    axis line style={-{Stealth[length=4pt]}, gray!65!black, line width=0.3pt},
    tick style={gray!65!black},
    tick label style={font=\tiny, gray!50!black},
    label style={font=\scriptsize},
    xlabel={$x$}, ylabel={$y$},
    every axis x label/.style={at={(ticklabel* cs:1.0)}, anchor=west},
    every axis y label/.style={at={(ticklabel* cs:1.0)}, anchor=south east,
      xshift=-1pt},
    every axis plot/.append style={accent, very thick, smooth, samples=201},
    clip=true, enlargelimits=false,
    title style={font=\scriptsize, yshift=4pt},
  },
  fnaxis/.append style={/pgfplots/legend style={font=\tiny, draw=none,
    fill=none, at={(0.98,0.03)}, anchor=south east}},
}

% A row of small graphs, evenly spaced with \hfill between them, plus an
% optional caption line: \begin{gallery} <pic>\hfill<pic> \gcap{...} \end{gallery}
\newenvironment{gallery}
  {\par\medskip\noindent\ignorespaces}
  {\par\medskip\noindent}
\newcommand{\gcap}[1]{\par\smallskip{\centering\footnotesize #1\par}}

% --- plane regions, implicit curves, surfaces -----------------------------
\usepgfplotslibrary{fillbetween}
\definecolor{regionfill}{HTML}{9DB4D8}
\pgfplotsset{colormap={house}{color=(soft) color=(accent!45) color=(accent)}}

\pgfplotsset{
  % Squared, gridded axis for plane regions and implicitly given curves.
  geoaxis/.style={
    fnaxis,
    width=5.9cm, height=5.9cm,
    grid=both,
    grid style={gray!22, line width=0.2pt},
    minor tick num=1,
  },
  % Level-curve maps: many thin curves, no single curve emphasised.
  contouraxis/.style={
    geoaxis,
    every axis plot/.append style={accent!75, line width=0.5pt, smooth,
      samples=201},
  },
  % Surfaces: \addplot3[surf, domain=..., y domain=...] {expr};
  surfaxis/.style={
    width=6.0cm, height=5.0cm,
    view={130}{22},
    colormap name=house,
    ticks=none,
    xlabel={$x$}, ylabel={$y$}, zlabel={$f$},
    label style={font=\tiny, gray!50!black},
    title style={font=\scriptsize},
    every axis plot/.append style={samples=25, shader=faceted interp,
      draw=accent!35, line width=0.1pt},
  },
}

% Shorthand for a shaded region between two named paths.
% \region{<path A>}{<path B>}{<soft clip options>}
\newcommand{\region}[3]{%
  \addplot[regionfill, opacity=0.55, forget plot, sharp plot]
    fill between[of=#1 and #2, #3];}

% --- rule boxes -----------------------------------------------------------
\newtcolorbox{rulebox}[1]{%
  breakable, enhanced, colback=soft, colframe=accent,
  boxrule=0.4pt, left=7pt, right=7pt, top=6pt, bottom=6pt,
  fonttitle=\bfseries, coltitle=white, colbacktitle=accent,
  title={#1}}

% --- worked examples ------------------------------------------------------
% One fully narrated solution per chapter, carrying the chapter's central
% technique and the mistakes that technique invites:
%   \begin{worked}{Task} ... \wstep{Step 1.} ... \wtrap{...} ... \end{worked}
\newtcolorbox{worked}[1]{%
  breakable, enhanced, colback=white, colframe=accent!65,
  boxrule=0.4pt, left=7pt, right=7pt, top=6pt, bottom=6pt,
  fonttitle=\bfseries, coltitle=accent, colbacktitle=soft,
  title={#1}}

% Run-in step heading inside a worked example.
\newcommand{\wstep}[1]{\par\smallskip\noindent\textbf{#1}\enspace\ignorespaces}

% A pitfall, set off from the running solution so it can be found again.
\newcommand{\wtrap}[1]{%
  \par\smallskip
  {\leftskip=1.3em\noindent
   \llap{\textcolor{accent}{\rule[0.22ex]{0.9em}{0.6pt}}\hspace{0.4em}}%
   \textbf{Trap.}\enspace\itshape #1\par}%
  \smallskip}

% --- problem / solution lists --------------------------------------------
\newlist{problems}{enumerate}{1}
\setlist[problems]{label=\textbf{\arabic*.}, leftmargin=2.1em,
  itemsep=0.55ex, parsep=0pt, topsep=0.8ex}
\newlist{solutions}{enumerate}{1}
\setlist[solutions]{label=\textbf{\arabic*.}, leftmargin=2.1em,
  itemsep=1.0ex, parsep=0pt, topsep=0.8ex}

% --- running heads --------------------------------------------------------
\pagestyle{fancy}
\fancyhf{}
\fancyhead[LE,RO]{\thepage}
\fancyhead[RE]{\small\itshape\nouppercase{\leftmark}}
\fancyhead[LO]{\small\itshape\nouppercase{\rightmark}}
\setlength{\headheight}{14pt}
\renewcommand{\headrulewidth}{0.3pt}
\renewcommand{\chaptermark}[1]{\markboth{\chaptername\ \thechapter.\ #1}{}}
\renewcommand{\sectionmark}[1]{\markright{\thesection\ #1}}

% --- shorthands used throughout -------------------------------------------
\newcommand{\dd}[2]{\frac{\mathrm{d}#1}{\mathrm{d}#2}}
\newcommand{\pd}[2]{\frac{\partial #1}{\partial #2}}
\newcommand{\pdd}[3]{\frac{\partial^{2} #1}{\partial #2\,\partial #3}}
\newcommand{\R}{\mathbb{R}}
\newcommand{\grad}{\nabla}
\DeclareMathOperator{\sgn}{sgn}
\DeclareMathOperator{\Hess}{H}

% German gloss, used on first mention of a technical term
\newcommand{\de}[1]{\emph{#1}}

\renewcommand*{\thisvolume}{mixed}

\title{\Huge Mixed Problems\\[0.4em]
  \large Cross-Chapter Sets with Worked Solutions}
\author{}
\date{}

\begin{document}
\frontmatter
\maketitle

\chapter*{About this sheet}
\addcontentsline{toc}{chapter}{About this sheet}

Every other volume in this series drills one technique at a time. That is
the right way to learn a method and the wrong way to prepare for an
examination, because a chapter heading tells you what to do before you have
read the question. An exam does not. The single most common way to lose
marks is not botched arithmetic but reaching for the wrong tool: setting up
a Lagrangian where the constraint does not bind, differentiating where a
limit was wanted, solving a system by inverting a matrix that turns out to
be singular.

This sheet exists to remove that crutch. Each set draws on the whole
sequence --- algebra and functions, matrices and elimination, limits and
L'H\^opital, curve discussion, multivariable calculus, constrained
optimization, linear programming --- in no announced order, and no problem
says which chapter it came from. Deciding that \emph{is} the exercise.

Work a set in one sitting, with a time limit and without looking anything
up. Sixty to ninety minutes for twenty problems is a realistic target. Then
read the solutions, and for every problem you got wrong, note not the
arithmetic but the \emph{decision}: which signal in the question should have
pointed at which method? That note is worth more than redoing the
computation.

Solutions name the technique in their first words, so the sets can also be
used the other way round --- read only the first line of each solution as a
drill in classification.

The fourth set works the other way. It announces the methods, all of them,
and takes a single function of two variables through the whole chain:
domain, range, derivative, level curves, gradient, Hessian, critical
points, contour sets and a checklist of properties. There the difficulty is
not choosing a tool but keeping a dozen answers consistent with one
another.

\tableofcontents

\mainmatter
\chapter{Mixed Set A}

Twenty problems, no order, no announced methods. Target time: 75 minutes.

\section{Worked Example}

The problems below announce no method. That is the point: on an
examination, choosing the technique is most of the work, and it is the one
step no drill chapter can train, because a drill chapter has already told
you the answer in its title.

The example here is \emph{not} one of the twenty. It is a typical example of
the kind, a single function examined from five directions, shown so
that you can watch the recognition step happen before you have to do it
yourself.

\begin{worked}{A five-part question on $f(x)=\dfrac{\ln x}{x}$: domain, both edges, extremum, inflection, and $\int_{1}^{\infty}f$}
\wstep{Step 1. Read the parts and name the tools before computing anything.}
\begin{center}\footnotesize
\begin{tabular}{lll}
  \hline
  part & tool & volume\\
  \hline
  domain & obstruction list & 2\\
  edges & limits, growth hierarchy, L'H\^opital & 4, 5\\
  extremum & $f'=0$, then $f''$ & 5\\
  inflection & sign change of $f''$ & 5\\
  convergence & improper integral, substitution & 6\\
  \hline
\end{tabular}
\end{center}
A minute spent on this table keeps the rest of the solution organised.

\wstep{Step 2. Domain.}
$\ln x$ needs $x>0$, and the denominator needs $x\neq0$. The first
condition subsumes the second:
\[
  D_{f}=(0,\infty).
\]

\wstep{Step 3. The left edge --- and it is \emph{not} indeterminate.}
As $x\to0^{+}$: $\ln x\to-\infty$ and $1/x\to+\infty$. The form is
$(-\infty)\cdot(+\infty)$, which is determinate:
\[
  \lim_{x\to0^{+}}\frac{\ln x}{x}=-\infty .
\]
\wtrap{Reaching for L'H\^opital here is the standard reflex and it is wrong
--- the rule needs $0/0$ or $\infty/\infty$, and this is neither. Always
name the form before choosing a tool. Here both factors push the same way,
so there is nothing to resolve.}

\wstep{Step 4. The right edge --- and here it \emph{is} indeterminate.}
As $x\to\infty$ the form is $\infty/\infty$, so either L'H\^opital or the
growth hierarchy applies:
\[
  \lim_{x\to\infty}\frac{\ln x}{x}
  \overset{\infty/\infty}{=}\lim_{x\to\infty}\frac{1/x}{1}=0 ,
\]
consistent with $\ln x\ll x$. So $y=0$ is a horizontal asymptote.

\wstep{Step 5. The extremum.}
\[
  f'(x)=\frac{\tfrac1x\cdot x-\ln x}{x^{2}}=\frac{1-\ln x}{x^{2}} .
\]
$f'=0$ when $\ln x=1$, i.e.\ $x=\mathrm{e}$, and
$f(\mathrm{e})=1/\mathrm{e}\approx0.368$.

The denominator $x^{2}$ is positive throughout the domain, so the sign of
$f'$ is the sign of $1-\ln x$: positive for $x<\mathrm{e}$, negative
after. A \textbf{maximum}, and since it is the only critical point and the
function rises to it from $-\infty$ and falls to $0$, it is the
\emph{global} maximum.
\wtrap{With a single critical point, the sign of $f'$ on either side settles
the classification faster than computing $f''$ --- and it settles the global
question too, which $f''$ never does.}

\wstep{Step 6. The inflection.}
\[
  f''(x)=\frac{2\ln x-3}{x^{3}} .
\]
This vanishes at $\ln x=\tfrac32$, i.e.\ $x=\mathrm{e}^{3/2}\approx4.482$,
and changes sign there (the denominator is positive), so it is a genuine
inflection at
\[
  \Bigl(\mathrm{e}^{3/2},\ \tfrac{3}{2}\mathrm{e}^{-3/2}\Bigr)\approx(4.48,\ 0.335).
\]
Check the ordering: $\mathrm{e}\approx2.72<4.48$, so the curve peaks and
\emph{then} changes curvature --- which is what a graph tending to zero from
above must do.

\wstep{Step 7. The integral.}
Substituting $u=\ln x$, $\mathrm{d}u=\mathrm{d}x/x$:
\[
  \int_{1}^{b}\frac{\ln x}{x}\,\mathrm{d}x
  =\int_{0}^{\ln b}u\,\mathrm{d}u
  =\frac{(\ln b)^{2}}{2}\xrightarrow[b\to\infty]{}\infty .
\]
\[
  \boxed{\int_{1}^{\infty}\frac{\ln x}{x}\,\mathrm{d}x\ \text{diverges.}}
\]
\wtrap{The integrand tends to $0$ --- Step~4 proved it --- and the integral
still diverges. ``The function goes to zero'' is never sufficient for
convergence, exactly as ``the terms go to zero'' was never sufficient for a
series. Here $f$ behaves like $1/x$ up to a logarithm, and $\int_{1}^{\infty}
\mathrm{d}x/x$ diverges.}

\wstep{Step 8. Assemble the picture, and check it against every part.}
Rising from $-\infty$ at the origin, crossing zero at $x=1$, peaking at
$(\mathrm{e},1/\mathrm{e})$, inflecting at $x\approx4.48$, and decaying to
$0$ --- with a tail thick enough to have infinite area. Every part of the
answer must be visible in that description, and if one is not, something is
wrong.
\begin{center}
\begin{tikzpicture}
\begin{axis}[fnaxis, width=8.0cm, height=4.4cm, domain=0.13:12,
             samples=201, xmin=-0.6, xmax=12.6, ymin=-2.2, ymax=0.9,
             xtick={1,2.718,4.482,8}, xticklabels={$1$,$\mathrm{e}$,$\mathrm{e}^{3/2}$,$8$},
             ytick={0.368}, yticklabels={$1/\mathrm{e}$}]
  \addplot {ln(x)/x};
  \addplot[only marks, mark=*, mark size=1.2pt, accent!70, forget plot]
    coordinates {(2.71828,0.36788) (4.48169,0.33470)};
\end{axis}
\end{tikzpicture}
\end{center}
\end{worked}

\section{Problems}

\begin{problems}
  \item Give the maximal domain of
        $f(x)=\dfrac{\sqrt{x^{2}-9}}{\ln(x-2)}$.
  \item Solve $3^{2x}-4\cdot3^{x}+3=0$.
  \item Find the inverse of $f(x)=\dfrac{2x-1}{x+3}$ and state its domain.
  \item Describe the set $\{(x,y): x^{2}+y^{2}\le25,\ y\ge x\}$ and decide
        whether it is convex.
  \item Give the level curves of $f(x,y)=x^{2}-4y$ at $c=0$ and $c=4$, and
        name the family.
  \item Evaluate $\displaystyle\lim_{x\to0}
        \frac{\mathrm{e}^{3x}-1-3x}{x^{2}}$.
  \item Evaluate $\displaystyle\lim_{x\to\infty}\frac{(\ln x)^{3}}{x}$.
  \item For $f(x)=\dfrac{x}{x^{2}-1}$: give the domain, the symmetry, all
        asymptotes, and the intervals of monotonicity.
  \item Solve $x+2y-z=3$, \;$2x-y+z=5$, \;$3x+y+2z=14$.
  \item Give the rank of
        $A=\begin{pmatrix}1&2&3\\2&4&6\\1&0&1\end{pmatrix}$ and say whether
        $A$ is invertible.
  \item Find the eigenvalues of
        $\begin{pmatrix}2&1\\1&2\end{pmatrix}$.
  \item Determine the definiteness of the Hessian of
        $f(x,y)=x^{2}+4xy+y^{2}$ and what it implies about the critical
        point at the origin.
  \item Find and classify all critical points of
        $f(x,y)=x^{3}-3xy+y^{3}$.
  \item Maximise $U=xy$ subject to $3x+4y=24$, and give the multiplier.
  \item Using the previous answer only, estimate the change in maximal
        utility when the income $24$ rises to $25$, then compute the exact
        change.
  \item Maximise $f(x,y)=x+y$ subject to $x^{2}+y^{2}\le2$, $x,y\ge0$.
  \item Maximise $3x+2y$ subject to $x+y\le4$, $x\le3$, $x,y\ge0$.
  \item Evaluate $\displaystyle\int x\,\mathrm{e}^{x}\,\mathrm{d}x$.
  \item Compute the price elasticity of demand for $q(p)=100/p^{2}$.
  \item At a continuously compounded rate of $5\%$, how long does a
        capital take to double?
\end{problems}

\section{Solutions}

\begin{solutions}
  \item \textbf{Maximal domain.} Three conditions: $x^{2}-9\ge0$ gives
        $|x|\ge3$; $\ln(x-2)$ needs $x>2$; and the denominator must not
        vanish, so $x-2\neq1$, i.e. $x\neq3$. Intersecting,
        $x\ge3$ and $x>2$ and $x\neq3$: \;$D=(3,\infty)$.
  \item \textbf{Substitution in an exponential equation.} With
        $u=3^{x}>0$: $u^{2}-4u+3=0$, so $u=1$ or $u=3$, giving
        $x=0$ and $x=1$.
  \item \textbf{Inverse function.} $y(x+3)=2x-1\Rightarrow x(y-2)=-1-3y$,
        so $x=\dfrac{3y+1}{2-y}$ and
        $f^{-1}(x)=\dfrac{3x+1}{2-x}$ with
        $D(f^{-1})=\R\setminus\{2\}$ --- which is exactly the range of
        $f$, as it must be.
  \item \textbf{Region in the plane.} The half of the disc of radius $5$
        lying on or above the line $y=x$. It is the intersection of two
        convex sets (a disc and a half-plane), hence \textbf{convex}.
  \item \textbf{Level curves.} $x^{2}-4y=c\Rightarrow y=\dfrac{x^{2}-c}{4}$:
        for $c=0$, $y=x^{2}/4$; for $c=4$, $y=\dfrac{x^{2}-4}{4}$. All
        level curves are vertical translates of the same parabola.
  \item \textbf{L'H\^opital, twice.} $\dfrac{3\mathrm{e}^{3x}-3}{2x}$ is
        still $0/0$, then
        $\dfrac{9\mathrm{e}^{3x}}{2}\to\dfrac92$.
  \item \textbf{Growth hierarchy.} $(\ln x)^{3}\ll x$, so the limit is
        $0$. (Or three applications of L'H\^opital, each reducing the power
        of the logarithm by one.)
  \item \textbf{Curve discussion.} Domain $\R\setminus\{-1,1\}$; the
        function is \emph{odd}; single root $x=0$; poles at $x=\pm1$;
        horizontal asymptote $y=0$. And
        $f'=\dfrac{(x^{2}-1)-x\cdot2x}{(x^{2}-1)^{2}}
        =\dfrac{-(x^{2}+1)}{(x^{2}-1)^{2}}<0$ everywhere, so $f$ is
        strictly decreasing on each of the three branches, with \emph{no}
        extrema.
  \item \textbf{Elimination.} Adding the first two equations gives
        $3x+y=8$; substituting into the third, $8+2z=14$ so $z=3$. Then
        $x+2y=6$ and $y=8-3x$ give $x=2$, $y=2$. Solution $(2,2,3)$.
        (Check the second: $4-2+3=5$ \checkmark.)
  \item \textbf{Rank by row reduction.} $R_{2}=2R_{1}$, so that row
        vanishes; $R_{3}-R_{1}=(0,-2,-2)$ survives.
        $\mathrm{rk}(A)=2<3$, so $A$ is \textbf{singular} --- equivalently
        $\det A=0$.
  \item \textbf{Characteristic polynomial.}
        $(2-\lambda)^{2}-1=0\Rightarrow\lambda-2=\pm1$, so
        $\lambda_{1}=3$ and $\lambda_{2}=1$. Both positive, so the matrix
        is positive definite.
  \item \textbf{Definiteness.}
        $\Hess f=\begin{pmatrix}2&4\\4&2\end{pmatrix}$ with
        $\det=4-16=-12<0$: \textbf{indefinite}. So the origin is a
        \textbf{saddle point}, not an extremum.
  \item \textbf{Critical points in two variables.}
        $f_{x}=3x^{2}-3y=0$ gives $y=x^{2}$; $f_{y}=-3x+3y^{2}=0$ gives
        $x=y^{2}$. Hence $x=x^{4}$, so $x(x^{3}-1)=0$ and $x=0$ or $x=1$:
        the points $(0,0)$ and $(1,1)$. With $f_{xx}=6x$, $f_{yy}=6y$,
        $f_{xy}=-3$: at $(0,0)$ the determinant is $0-9=-9<0$, a
        \textbf{saddle}; at $(1,1)$ it is $36-9=27>0$ with $f_{xx}=6>0$, a
        \textbf{local minimum} with $f=-1$.
  \item \textbf{Lagrange.} $\mathcal{L}=xy+\lambda(24-3x-4y)$ gives
        $y=3\lambda$ and $x=4\lambda$, so $4y=3x$. With the constraint,
        $3x+3x=24$, hence $x^{*}=4$, $y^{*}=3$, $U^{*}=12$ and
        $\lambda^{*}=1$.
  \item \textbf{Envelope theorem.} $\mathrm{d}U^{*}/\mathrm{d}M
        =\lambda^{*}=1$, so the estimate is $+1$. Exactly:
        $U^{*}(M)=M^{2}/48$, so
        $U^{*}(25)-U^{*}(24)=\dfrac{625-576}{48}=\dfrac{49}{48}
        \approx1.021$. The multiplier is the first-order term.
  \item \textbf{Kuhn--Tucker (or symmetry).} The objective increases in
        both variables, so the circle constraint binds; by symmetry
        $x=y$, and $2x^{2}=2$ gives $x=y=1$. The non-negativity
        constraints are slack, and $f^{*}=2$. (Formally
        $\grad f=(1,1)^{T}=\lambda(2x,2y)^{T}$ forces $x=y$ with
        $\lambda=\tfrac12>0$.)
  \item \textbf{Linear programming, graphically.} The feasible vertices
        are $(0,0)$, $(3,0)$, $(3,1)$ and $(0,4)$, where the objective
        takes the values $0$, $9$, $11$ and $8$. The maximum is
        $\mathbf{11}$ at $(3,1)$, where both $x\le3$ and $x+y\le4$ bind.
  \item \textbf{Integration by parts.} With $u=x$,
        $v'=\mathrm{e}^{x}$:
        $\int x\mathrm{e}^{x}\mathrm{d}x
        =x\mathrm{e}^{x}-\int\mathrm{e}^{x}\mathrm{d}x
        =(x-1)\mathrm{e}^{x}+C$.
  \item \textbf{Elasticity.}
        $\varepsilon=\dfrac{q'(p)\,p}{q(p)}
        =\dfrac{-200p^{-3}\cdot p}{100p^{-2}}=-2$: constant, independent of
        the price --- the signature of a power demand function.
  \item \textbf{Continuous compounding.}
        $\mathrm{e}^{0.05t}=2\Rightarrow t=\dfrac{\ln2}{0.05}
        \approx13.86$ years. (The ``rule of $70$'': $70/5=14$.)
\end{solutions}

\chapter{Mixed Set B}

Twenty problems, no order, no announced methods. Target time: 75 minutes.

\section{Worked Example}

Again the example is not one of the twenty. This one is a constrained
optimisation, one of the most frequently examined topics, worked
with every check that the marker expects to see.

\begin{worked}{Maximise $f(x,y)=\ln x+2\ln y$ subject to $x+y=12$, with $x,y>0$: solve, certify, and interpret the multiplier}
\wstep{Step 1. Recognise the type before writing anything.}
Two variables, one equality constraint, a differentiable objective: this is
Lagrange. The constraint is linear and could be substituted instead --- keep
that in reserve as the check.

\wstep{Step 2. Set up and differentiate.}
\[
  \mathcal L=\ln x+2\ln y+\lambda(12-x-y),
\]
\[
  \mathcal L_{x}=\frac1x-\lambda=0,
  \qquad
  \mathcal L_{y}=\frac2y-\lambda=0,
  \qquad
  x+y=12 .
\]

\wstep{Step 3. Eliminate $\lambda$.}
\[
  \frac1x=\frac2y
  \quad\Longrightarrow\quad
  y=2x .
\]
The optimal bundle puts twice as much into $y$ as into $x$ --- the ratio of
the coefficients $2$ and $1$ in the objective, exactly as Cobb--Douglas
shares behave, and $\ln$-linear objectives are Cobb--Douglas in disguise.

\wstep{Step 4. Impose the constraint.}
\[
  x+2x=12
  \quad\Longrightarrow\quad
  \boxed{x^{*}=4,\quad y^{*}=8,\quad \lambda^{*}=\tfrac14 }
\]
\[
  f^{*}=\ln4+2\ln8=\ln4+\ln64=\ln256\approx5.545 .
\]

\wstep{Step 5. Certify it with the bordered Hessian.}
\[
  \mathcal L_{xx}=-\frac{1}{x^{2}}=-\tfrac{1}{16},
  \qquad
  \mathcal L_{yy}=-\frac{2}{y^{2}}=-\tfrac{1}{32},
  \qquad
  \mathcal L_{xy}=0,
  \qquad
  g_{x}=g_{y}=1 .
\]
\[
  \overline H=\begin{pmatrix}0&1&1\\1&-\tfrac{1}{16}&0\\1&0&-\tfrac{1}{32}\end{pmatrix},
  \qquad
  \det\overline H=\tfrac{1}{32}+\tfrac{1}{16}=\tfrac{3}{32}>0 ,
\]
which for $n=2$, $m=1$ means a constrained \textbf{maximum}. \checkmark

\wstep{Step 6. Check again by substitution, which is available here.}
$y=12-x$ gives $h(x)=\ln x+2\ln(12-x)$, so
\[
  h'(x)=\frac1x-\frac{2}{12-x}=0
  \;\Longrightarrow\;
  12-x=2x
  \;\Longrightarrow\;
  x=4 ,
\]
\[
  h''(x)=-\frac{1}{x^{2}}-\frac{2}{(12-x)^{2}}<0\ \text{everywhere:\ concave, so the maximum is global.}
\]
\wtrap{Two independent routes to the same answer is the best use of spare
examination time --- far better than re-reading the first one. And here the
substitution route delivers something the bordered Hessian does not: a
\emph{global} conclusion, because $h''<0$ on the whole feasible interval,
not merely at the candidate point.}

\wstep{Step 7. Verify the multiplier's meaning.}
With a general budget $M$ in place of $12$: $x^{*}=M/3$, $y^{*}=2M/3$, so
\[
  V(M)=\ln\frac M3+2\ln\frac{2M}{3}=3\ln M+\text{constant},
  \qquad
  \dd{V}{M}=\frac3M .
\]
At $M=12$ this is $\tfrac14=\lambda^{*}$. \checkmark\ One more unit of budget
raises the maximised objective by about $0.25$.

\wstep{Step 8. The checks a marker looks for.}
\begin{center}\footnotesize
\begin{tabular}{ll}
  \hline
  check & here\\
  \hline
  constraint satisfied & $4+8=12$ \checkmark\\
  domain respected & $x,y>0$ \checkmark\\
  second-order condition & $\det\overline H>0$ \checkmark\\
  multiplier sign & $\lambda>0$: relaxing helps \checkmark\\
  interpretation stated & $\lambda=\partial V/\partial M$ \checkmark\\
  \hline
\end{tabular}
\end{center}
Every one of these is a line, and together they are usually worth more marks
than the algebra that produced the numbers.
\end{worked}

\section{Problems}

\begin{problems}
  \item Simplify $\dfrac{x^{2}-4}{x^{2}-x-6}$ and state where the
        simplification is valid.
  \item Solve $\log_{2}x+\log_{2}(x-2)=3$.
  \item Give the range of $f(x)=\dfrac{x^{2}+1}{x}$ on $x>0$.
  \item Write a system of inequalities for the part of the elliptical disc
        $\dfrac{x^{2}}{4}+\dfrac{y^{2}}{9}\le1$ lying in the second
        quadrant.
  \item Two functions have concentric circles as level curves. In one the
        circles are equally spaced, in the other they crowd outwards. Which
        is $\sqrt{x^{2}+y^{2}}$ and which is $x^{2}+y^{2}$?
  \item Evaluate $\displaystyle\lim_{x\to0^{+}}x^{\sin x}$.
  \item Evaluate $\displaystyle\lim_{x\to2}\frac{x^{3}-8}{x^{2}-4}$.
  \item For $f(x)=\dfrac{\mathrm{e}^{x}}{x}$: give the domain, the
        asymptotes, the monotonicity and any extremum.
  \item Solve $2x+y=7$, \;$4x+2y=14$, and describe the solution set.
  \item Compute
        $\det\begin{pmatrix}1&2&0\\0&3&1\\2&0&4\end{pmatrix}$.
  \item Solve $\begin{pmatrix}2&1\\1&3\end{pmatrix}
        \mathbf{x}=\binom{5}{10}$.
  \item Give $\grad f(1,2)$ for $f(x,y)=x^{2}y^{3}$.
  \item Give the directional derivative of the same $f$ at $(1,2)$ in the
        direction of $(3,4)^{T}$.
  \item Use the total differential to estimate $f(1.1,1.9)$ for the same
        $f$, and compare with the exact value.
  \item Find $\dd{y}{x}$ at $(1,2)$ on the curve $x^{2}+xy+y^{2}=7$.
  \item Minimise $2x+3y$ subject to $xy=54$, with $x,y>0$.
  \item Maximise $\ln x+y$ subject to $x+y\le4$, $x>0$, $y\ge0$.
  \item Maximise $4x+3y$ subject to $x+y\le10$, $2x+y\le16$, $x,y\ge0$.
  \item What is $1000$ received in five years worth today at $6\%$ annual
        interest?
  \item Sum the series $3+\tfrac32+\tfrac34+\tfrac38+\cdots$.
\end{problems}

\section{Solutions}

\begin{solutions}
  \item \textbf{Factor and cancel.}
        $\dfrac{(x-2)(x+2)}{(x-3)(x+2)}=\dfrac{x-2}{x-3}$, valid for
        $x\neq-2$ and $x\neq3$. The point $x=-2$ remains excluded from the
        domain even though the factor cancelled --- a hole, not a pole.
  \item \textbf{Logarithm law, then check the domain.}
        $\log_{2}\bigl(x(x-2)\bigr)=3\Rightarrow x^{2}-2x=8$, so
        $x^{2}-2x-8=0$ and $x=4$ or $x=-2$. The domain requires $x>2$, so
        $x=-2$ is rejected: \;$x=4$.
  \item \textbf{Range via AM--GM (or calculus).}
        $f(x)=x+\dfrac1x$, and for $x>0$, $x+1/x\ge2$ with equality at
        $x=1$. As $x\to0^{+}$ and $x\to\infty$ the function grows without
        bound, so $W=[2,\infty)$.
  \item \textbf{Region.} $\dfrac{x^{2}}{4}+\dfrac{y^{2}}{9}\le1$,
        \;$x\le0$, \;$y\ge0$.
  \item \textbf{Contour spacing.} Equally spaced circles belong to the
        \textbf{cone} $\sqrt{x^{2}+y^{2}}$, whose level $c$ has radius $c$.
        Circles crowding outwards belong to the \textbf{paraboloid}
        $x^{2}+y^{2}$, whose level $c$ has radius $\sqrt{c}$: the surface
        steepens, so equal steps in height need ever smaller steps in
        radius.
  \item \textbf{Exponential form, take logs.}
        $\ln y=\sin x\ln x$, and since $\sin x\sim x$ near $0$ this behaves
        like $x\ln x\to0$. Hence the limit is $\mathrm{e}^{0}=1$.
  \item \textbf{L'H\^opital (or factor).}
        $\dfrac{3x^{2}}{2x}\to\dfrac{12}{4}=3$. Factoring gives the same:
        $\dfrac{(x-2)(x^{2}+2x+4)}{(x-2)(x+2)}\to\dfrac{12}{4}$.
  \item \textbf{Curve discussion.} Domain $\R\setminus\{0\}$. As
        $x\to0^{+}$, $f\to+\infty$; as $x\to0^{-}$, $f\to-\infty$: a pole
        at $x=0$. As $x\to-\infty$, $f\to0^{-}$, so $y=0$ is a horizontal
        asymptote on the left only; as $x\to+\infty$, $f\to+\infty$. And
        $f'=\dfrac{\mathrm{e}^{x}(x-1)}{x^{2}}$, negative for $x<1$
        (excluding $0$) and positive for $x>1$: decreasing, then
        increasing, with a local (and global, on $x>0$) minimum at
        $(1,\mathrm{e})$.
  \item \textbf{Elimination.} The second equation is exactly twice the
        first, so it carries no information: one pivot, two unknowns.
        \textbf{Infinitely many} solutions, $y=7-2x$ with $x$ free ---
        the line itself.
  \item \textbf{Cofactor expansion along the first row.}
        $1\cdot(3\cdot4-1\cdot0)-2\cdot(0\cdot4-1\cdot2)+0
        =12+4=\mathbf{16}$.
  \item \textbf{Inverse or elimination.} $\det=6-1=5$, so
        $\mathbf{x}=\dfrac15\begin{pmatrix}3&-1\\-1&2\end{pmatrix}
        \binom{5}{10}=\dfrac15\binom{5}{15}=\binom{1}{3}$.
        (Check: $2+3=5$ \checkmark, $1+9=10$ \checkmark.)
  \item \textbf{Partial derivatives.} $f_{x}=2xy^{3}$ and
        $f_{y}=3x^{2}y^{2}$, so at $(1,2)$:
        $\grad f=(2\cdot8,\;3\cdot4)^{T}=(16,12)^{T}$.
  \item \textbf{Directional derivative --- normalise first.}
        $\|(3,4)\|=5$, so $\mathbf{u}=(0.6,0.8)^{T}$ and
        $D_{\mathbf{u}}f=(16,12)\cdot(0.6,0.8)=9.6+9.6=\mathbf{19.2}$.
  \item \textbf{Total differential.} $f(1,2)=1\cdot8=8$ and
        $\mathrm{d}f=16\,\mathrm{d}x+12\,\mathrm{d}y
        =16(0.1)+12(-0.1)=0.4$, so $f(1.1,1.9)\approx8.4$. Exactly
        $1.1^{2}\cdot1.9^{3}=1.21\cdot6.859=8.299$; the linear estimate is
        high by about $0.1$, the second-order error.
  \item \textbf{Implicit differentiation.} $2x+y+xy'+2yy'=0$, so
        $y'=-\dfrac{2x+y}{x+2y}$. At $(1,2)$ --- which does lie on the
        curve, since $1+2+4=7$ --- this is
        $-\dfrac{4}{5}$.
  \item \textbf{Lagrange (cost minimisation).}
        $2=\lambda y$ and $3=\lambda x$ give $\dfrac{x}{y}=\dfrac32$, so
        $x=1.5y$. Then $1.5y^{2}=54$, $y^{2}=36$, $y=6$ and $x=9$. Minimal
        cost $18+18=36$ --- the two expenditures are equal, as always for
        a product constraint with linear costs.
  \item \textbf{Kuhn--Tucker.} With
        $\mathcal{L}=\ln x+y+\lambda(4-x-y)$: $1/x=\lambda$ and
        $1=\lambda$, so $\lambda=1$ and $x=1$. The budget binds, so $y=3$,
        which satisfies $y\ge0$. Optimum $(1,3)$ with value
        $\ln1+3=3$.
  \item \textbf{Linear programming.} The two constraints meet where
        $x+y=10$ and $2x+y=16$, i.e. $(6,4)$. Vertices $(0,0)$, $(8,0)$,
        $(6,4)$, $(0,10)$ give objective values $0$, $32$, $36$, $30$: the
        maximum is $\mathbf{36}$ at $(6,4)$, where both constraints bind.
  \item \textbf{Discounting.}
        $\dfrac{1000}{1.06^{5}}=\dfrac{1000}{1.33823}\approx747.26$.
  \item \textbf{Geometric series.} First term $a=3$, ratio
        $r=\tfrac12$, so the sum is $\dfrac{3}{1-\tfrac12}=6$.
\end{solutions}

\chapter{Mixed Set C}

Twenty problems, no order, no announced methods. Target time: 75 minutes.

\section{Worked Example}

The last example is an applied one, of the kind that hides three techniques
inside a sentence of prose. As before it is not among the twenty.

\begin{worked}{A machine yields a net flow of $D(t)=8000\,\mathrm{e}^{-0.1t}$ euros per year, discounted continuously at $r=5\%$. Find (a) the present value over ten years, (b) over an infinite horizon, (c) the time at which half the infinite value has accrued}
\wstep{Step 1. Translate the sentence into an integral.}
``Continuous flow'' plus ``continuous discounting'' means
\[
  PV(T)=\int_{0}^{T}D(t)\,\mathrm{e}^{-rt}\,\mathrm{d}t .
\]
\wtrap{Two exponentials appear and they are different things: $\mathrm{e}^{-0.1t}$
is the machine \emph{decaying}, $\mathrm{e}^{-0.05t}$ is money being
\emph{discounted}. Both must be present, and using only one is the standard
misreading of the sentence.}

\wstep{Step 2. Combine the exponents --- this is what makes it easy.}
\[
  PV(T)=\int_{0}^{T}8000\,\mathrm{e}^{-0.1t}\mathrm{e}^{-0.05t}\,\mathrm{d}t
  =8000\int_{0}^{T}\mathrm{e}^{-0.15t}\,\mathrm{d}t .
\]
The two rates \emph{add}, because $\mathrm{e}^{a}\mathrm{e}^{b}
=\mathrm{e}^{a+b}$: the effective rate of decline is $15\%$.

\wstep{Step 3. Integrate.}
\[
  8000\Bigl[\frac{\mathrm{e}^{-0.15t}}{-0.15}\Bigr]_{0}^{T}
  =\frac{8000}{0.15}\bigl(1-\mathrm{e}^{-0.15T}\bigr)
  =53\,333.33\,\bigl(1-\mathrm{e}^{-0.15T}\bigr).
\]
\wtrap{The factor $1/(-0.15)$ turns $\mathrm{e}^{-0.15T}-1$ into
$1-\mathrm{e}^{-0.15T}$. Get this backwards and the present value comes out
negative, which signals an error: a positive flow cannot have a
negative value.}

\wstep{Step 4. Part (a): ten years.}
\[
  PV(10)=53\,333.33\,\bigl(1-\mathrm{e}^{-1.5}\bigr)
        =53\,333.33\,(1-0.22313)
        =\boxed{41\,433\ \text{euros}}
\]

\wstep{Step 5. Part (b): the infinite horizon.}
\[
  PV(\infty)=\lim_{T\to\infty}53\,333.33\bigl(1-\mathrm{e}^{-0.15T}\bigr)
  =\boxed{53\,333\ \text{euros}}
\]
The limit exists because $0.15>0$. Had the machine's output grown faster
than the discount rate --- $D(t)=8000\mathrm{e}^{+0.06t}$ against $r=0.05$
--- the integral would diverge, and the asset would have no finite value.
That convergence condition is the same $g<r$ that governs a growing
perpetuity in \chref{foundations}{2}.

\wstep{Step 6. Part (c): solve for the time, which is a logarithm problem.}
\[
  53\,333.33\bigl(1-\mathrm{e}^{-0.15T}\bigr)=\tfrac12\cdot53\,333.33
  \;\Longrightarrow\;
  1-\mathrm{e}^{-0.15T}=\tfrac12
  \;\Longrightarrow\;
  \mathrm{e}^{-0.15T}=\tfrac12 .
\]
\[
  -0.15T=\ln\tfrac12=-\ln2
  \quad\Longrightarrow\quad
  T=\frac{\ln2}{0.15}=\boxed{4.62\ \text{years}}
\]
The constant $53\,333.33$ cancels before any arithmetic is done --- notice
that before computing it.
\wtrap{This is a half-life, and it depends only on the \emph{combined} rate
$0.15$, not on the size of the flow. Recognising the shape
$\mathrm{e}^{-kT}=\tfrac12\Rightarrow T=\ln2/k$ saves the whole
manipulation.}

\wstep{Step 7. Sanity-check all three answers together.}
$41\,433<53\,333$ \checkmark: ten years cannot be worth more than the infinite horizon. The ten-year
figure is $77.7\%$ of the total, and $4.62$ years already delivers $50\%$
\checkmark --- consistent with a flow that decays. And $53\,333=8000/0.15$
\checkmark, the perpetuity formula with the combined rate.

\wstep{Step 8. What made this an applied problem rather than a calculus exercise.}
Three separate recognitions: that ``continuous flow'' means an integral,
that two exponential rates add, and that ``half the value'' is an equation
to be solved rather than a quantity to be computed. The calculus itself was
one substitution. That ratio --- most of the difficulty in the translation,
little in the technique --- is characteristic of the applied problems in
these sets, and it is why reading the sentence twice is time well spent.
\end{worked}

\section{Problems}

\begin{problems}
  \item Solve $|2x-5|\le3$.
  \item Give the maximal domain of $f(x)=\ln(x^{2}-x-6)$.
  \item Describe $y=-2\sqrt{x+1}+3$ as a chain of transformations, and give
        its domain and range.
  \item Describe the region $\{y\ge x^{2},\ y\le 2x+3\}$ and give the range
        of $x$ it occupies.
  \item Give the level curve of $f(x,y)=xy$ at $c=-4$ and say which
        quadrants it lies in.
  \item Evaluate $\displaystyle\lim_{x\to\infty}
        x\bigl(\sqrt{x^{2}+1}-x\bigr)$.
  \item Evaluate $\displaystyle\lim_{x\to1}\frac{x^{n}-1}{x-1}$ for a
        positive integer $n$.
  \item Carry out a complete discussion of $f(x)=x^{3}-12x$.
  \item For which $t$ does the system $x+y+tz=1$, \;$x+ty+z=1$,
        \;$tx+y+z=1$ fail to have a unique solution?
  \item Give the eigenvalues of
        $\begin{pmatrix}4&2\\2&1\end{pmatrix}$ and its definiteness.
  \item Determine the definiteness of the Hessian of
        $f(x,y)=x^{2}+y^{2}-2xy$ and describe the set of minimisers.
  \item At the point $(3,4)$ on $x^{2}+y^{2}=25$, may $y$ be written
        locally as a function of $x$? Give $\dd{y}{x}$ there.
  \item Show that $f(K,L)=K^{0.3}L^{0.7}$ is homogeneous of degree $1$ and
        verify Euler's theorem $Kf_{K}+Lf_{L}=f$.
  \item Maximise $-(x-2)^{2}-(y-3)^{2}$ subject to $x+y\le10$. Does the
        constraint bind?
  \item For $\max xy$ subject to $px+qy=M$, show $x^{*}=M/(2p)$ and give
        $\partial x^{*}/\partial p$.
  \item Minimise $2x+3y$ subject to $x+y\ge4$, $x\ge1$, $y\ge0$.
  \item Evaluate $\displaystyle\int_{0}^{1}x^{2}\mathrm{e}^{x^{3}}
        \,\mathrm{d}x$.
  \item Evaluate $\displaystyle\int_{1}^{\infty}\frac{1}{x^{2}}
        \,\mathrm{d}x$.
  \item Demand is $p=20-2q$. At the equilibrium quantity $q=5$, compute the
        consumer surplus.
  \item What is an annuity of $500$ per year for ten years worth today at
        $4\%$?
\end{problems}

\section{Solutions}

\begin{solutions}
  \item \textbf{Absolute-value inequality.}
        $-3\le2x-5\le3\Rightarrow2\le2x\le8$, so $1\le x\le4$.
  \item \textbf{Logarithm domain, hence a quadratic inequality.}
        $x^{2}-x-6=(x-3)(x+2)>0$ holds outside the roots:
        $D=(-\infty,-2)\cup(3,\infty)$.
  \item \textbf{Transformations.} From $\sqrt{x}$: shift $1$ left, stretch
        vertically by $2$, reflect in the $x$-axis, shift $3$ up. Domain
        $[-1,\infty)$; since $-2\sqrt{x+1}\in(-\infty,0]$, the range is
        $(-\infty,3]$, with the curve starting at $(-1,3)$ and falling.
  \item \textbf{Region between two curves.} They meet where
        $x^{2}=2x+3$, i.e. $x=-1$ and $x=3$, so the region is the lens
        above the parabola and below the line, occupying
        $-1\le x\le3$. Convex.
  \item \textbf{Level curve.} $xy=-4\Rightarrow y=-4/x$: a hyperbola whose
        branches lie in the \textbf{second and fourth} quadrants, since the
        product is negative.
  \item \textbf{Conjugate, then a limit at infinity.}
        $x\bigl(\sqrt{x^{2}+1}-x\bigr)
        =\dfrac{x}{\sqrt{x^{2}+1}+x}
        =\dfrac{1}{\sqrt{1+1/x^{2}}+1}\to\dfrac12$.
  \item \textbf{L'H\^opital (or the factorisation of $x^{n}-1$).}
        $\dfrac{nx^{n-1}}{1}\to n$.
  \item \textbf{Complete discussion.} Domain $\R$; \emph{odd}, so
        symmetric about the origin; $f=x(x^{2}-12)$ gives roots
        $0,\pm2\sqrt3$; ends $\mp\infty$. $f'=3x^{2}-12=3(x-2)(x+2)$:
        increasing on $(-\infty,-2)$, decreasing on $(-2,2)$, increasing on
        $(2,\infty)$, so a local maximum $(-2,16)$ and a local minimum
        $(2,-16)$. $f''=6x$: concave for $x<0$, convex for $x>0$, with an
        inflection at the origin $(0,0)$.
  \item \textbf{Determinant with a parameter.} Expanding,
        \[
          \det\begin{pmatrix}1&1&t\\1&t&1\\t&1&1\end{pmatrix}
          =(t-1)\bigl[2-t(1+t)\bigr]
          =-(t-1)^{2}(t+2).
        \]
        This vanishes for $t=1$ and $t=-2$. At $t=1$ all three equations
        become $x+y+z=1$: infinitely many solutions. At $t=-2$ the system
        is inconsistent (adding all three gives $0=3$).
  \item \textbf{Eigenvalues.} Trace $5$ and determinant $4-4=0$, so
        $\lambda^{2}-5\lambda=0$ and $\lambda=0$ or $\lambda=5$. Both
        $\ge0$ with one zero: \textbf{positive semidefinite}, not definite.
  \item \textbf{Hessian and the minimiser set.}
        $\Hess f=\begin{pmatrix}2&-2\\-2&2\end{pmatrix}$ with
        $\det=0$ and trace $4>0$: positive semidefinite, so $f$ is convex
        but not strictly so. Indeed $f=(x-y)^{2}\ge0$, and the minimum
        value $0$ is attained on the whole \textbf{line} $y=x$ --- the
        direction in which the Hessian is singular.
  \item \textbf{Implicit function theorem.} With
        $F=x^{2}+y^{2}-25$: $F_{y}=2y=8\neq0$ at $(3,4)$, so yes, a local
        $y(x)$ exists, and
        $\dd{y}{x}=-\dfrac{F_{x}}{F_{y}}=-\dfrac{6}{8}=-\dfrac34$.
  \item \textbf{Homogeneity and Euler.}
        $f(tK,tL)=(tK)^{0.3}(tL)^{0.7}=t^{0.3+0.7}K^{0.3}L^{0.7}=tf(K,L)$,
        so the degree is $1$. And
        $f_{K}=0.3K^{-0.7}L^{0.7}$, $f_{L}=0.7K^{0.3}L^{-0.3}$, so
        \[
          Kf_{K}+Lf_{L}=0.3K^{0.3}L^{0.7}+0.7K^{0.3}L^{0.7}
          =K^{0.3}L^{0.7}=f .
        \]
        Economically: with constant returns to scale, paying each factor
        its marginal product exactly exhausts the product.
  \item \textbf{Check the unconstrained optimum first.} The objective is
        maximal at $(2,3)$, where $x+y=5\le10$: the point is already
        feasible, so the constraint does \textbf{not} bind, its multiplier
        is $0$, and the maximum is $f=0$ at $(2,3)$. Setting up a
        Lagrangian with an equality would have produced the wrong answer.
  \item \textbf{Lagrange with general parameters.} $y=\lambda p$ and
        $x=\lambda q$ give $px=qy$, so the budget splits evenly:
        $px=M/2$, hence $x^{*}=\dfrac{M}{2p}$ and $y^{*}=\dfrac{M}{2q}$.
        Then $\dfrac{\partial x^{*}}{\partial p}=-\dfrac{M}{2p^{2}}<0$:
        the own-price effect is negative, and the cross-price effect
        $\partial x^{*}/\partial q=0$ --- Cobb--Douglas demands do not
        respond to the other price.
  \item \textbf{Linear programming, minimisation.} The vertices of the
        feasible set are $(4,0)$ (where $x+y=4$ meets $y=0$) and $(1,3)$
        (where $x=1$ meets $x+y=4$). The objective takes the values
        $8$ and $11$, so the minimum is $\mathbf{8}$ at $(4,0)$. The region
        is unbounded, but the objective increases in both variables, so a
        minimum exists.
  \item \textbf{Substitution.} With $u=x^{3}$, $\mathrm{d}u=3x^{2}
        \mathrm{d}x$:
        $\dfrac13\displaystyle\int_{0}^{1}\mathrm{e}^{u}\mathrm{d}u
        =\dfrac{\mathrm{e}-1}{3}\approx0.573$.
  \item \textbf{Improper integral.}
        $\displaystyle\int_{1}^{b}x^{-2}\mathrm{d}x
        =1-\dfrac1b\to1$ as $b\to\infty$: the integral \textbf{converges}
        to $1$.
  \item \textbf{Consumer surplus.} At $q=5$ the price is
        $p=20-10=10$, and
        \[
          CS=\int_{0}^{5}(20-2q)\,\mathrm{d}q-10\cdot5
            =\bigl[20q-q^{2}\bigr]_{0}^{5}-50
            =75-50=25 .
        \]
  \item \textbf{Annuity present value.}
        $500\cdot\dfrac{1-1.04^{-10}}{0.04}
        =500\cdot8.1109\approx4055.4$.
\end{solutions}

\chapter{Mixed Set D: Full Investigations}

The first three sets drill one skill per problem. This set does the
opposite. Each problem hands you a single function of two variables and
runs it through the whole chain of questions that a function of two
variables invites: where it lives, what values it takes, where it can be
differentiated, what its level curves look like, where its gradient
points, what its Hessian says, whether it has critical points, which of
its contour sets are convex, and finally a checklist of properties, each
to be decided \emph{with a reason}.

None of the steps is new. What is new is that they feed each other: the
domain decides where the derivative may be taken, the Hessian decides
concavity, concavity decides quasiconcavity, the domain and the range
decide whether a maximum exists. A slip early in the chain propagates,
and a good investigation catches it, because the later answers stop
fitting together. Target time: twenty to twenty-five minutes per
function. The solutions tag every step with the volume and chapter whose
rule it uses, so that a weak link can be traced back and repaired.

\section{The Investigation Chain}

\begin{rulebox}{The chain, in order}
\begin{enumerate}[label=\arabic*.,leftmargin=1.6em,itemsep=0.3ex]
  \item \textbf{Domain.} Collect the obstructions (roots need a
        nonnegative radicand, logarithms a positive argument,
        denominators must not vanish), sketch the set, and say whether it
        is closed, bounded and convex. (Vol.~2, Ch.~1 and Ch.~4;
        Vol.~5, Ch.~10; Vol.~7, Ch.~2.)
  \item \textbf{Range.} Find the smallest and largest values, or show that
        there are none, and use continuity on a connected domain to fill
        in everything between. (Vol.~2, Ch.~1; Vol.~5, Ch.~9.)
  \item \textbf{Derivative.} Compute $f_{x}$, $f_{y}$ and say where they
        exist and are continuous: on that open set $f$ is differentiable
        (continuous partials are sufficient, not necessary). Where a
        partial fails to exist, $f$ is not differentiable. Roots are the
        usual offenders on the boundary. (Vol.~5, Ch.~10.)
  \item \textbf{Level curves.} Solve $f(x,y)=c$ for one variable and
        sketch several values of $c$. (Vol.~5, Ch.~9.)
  \item \textbf{Gradient at a point.} Evaluate $\grad f$, draw it, check
        that it is orthogonal to the level curve through the point, and
        read it: it points in the direction of steepest ascent, and its
        norm $\lVert\grad f\rVert$ is the slope in that direction.
        (Vol.~5, Ch.~10 and Ch.~13.)
  \item \textbf{Hessian.} Second partials, Young's theorem, and the
        definiteness by Sylvester's criterion. (Vol.~5, Ch.~11.)
  \item \textbf{Critical points.} Solve $\grad f=\mathbf{0}$ \emph{inside
        the domain}, or explain why there is no solution; classify with the
        Hessian. (Vol.~7, Ch.~5.)
  \item \textbf{Contour sets.} Mark an upper set $\{f\ge c\}$ and a lower
        set $\{f\le c\}$ and decide whether each is convex.
        (Vol.~5, Ch.~9; Vol.~7, Ch.~2.)
  \item \textbf{Checklist.} Concave, convex, quasiconcave, quasiconvex,
        has a maximum, has a minimum, increasing, decreasing, epigraph
        convex: each answered with a proof or a counterexample.
        (Vol.~7, Ch.~1 and Ch.~3.)
\end{enumerate}
\end{rulebox}

\begin{rulebox}{Three definitions the checklist needs}
\begin{itemize}[leftmargin=1.6em,itemsep=0.3ex]
  \item \textbf{Monotone in several variables.} Order points componentwise:
        $(x,y)\le(x',y')$ means $x\le x'$ and $y\le y'$. Then $f$ is
        \emph{increasing} \de{monoton wachsend} if
        $(x,y)\le(x',y')\Rightarrow f(x,y)\le f(x',y')$, and
        \emph{decreasing} if the inequality between the values is reversed.
        On a convex domain with interior points, a continuous $f$ with
        $f_{x}\ge0$ and $f_{y}\ge0$ at every interior point is increasing.
        To refute it, exhibit two ordered points whose values are in the
        wrong order.
  \item \textbf{Epigraph and hypograph.} The \emph{epigraph}
        \de{Epigraph} of $f$ is the set of points on or above its graph,
        $\{(x,y,z): (x,y)\in D,\ z\ge f(x,y)\}$; the \emph{hypograph} is
        the set on or below it. On a convex domain, the epigraph is a convex
        set exactly when $f$ is convex, and the hypograph exactly when $f$
        is concave. So the epigraph question is the convexity question
        asked again, and the two answers must agree.
  \item \textbf{Existence of extrema.} A continuous function on a closed
        and bounded domain attains a maximum and a minimum (Weierstrass,
        Vol.~7, Ch.~1). On an unbounded or open domain, look at the range
        first: if $f$ grows without bound there is no maximum, and if the
        infimum is approached but never reached there is no minimum.
        Extrema on the boundary of the domain are not found by
        $\grad f=\mathbf{0}$, and are often points where $f$ is not even
        differentiable.
\end{itemize}
\end{rulebox}

\begin{rulebox}{Shortcuts that keep the checklist consistent}
\begin{itemize}[leftmargin=1.6em,itemsep=0.3ex]
  \item Concave $\Rightarrow$ quasiconcave; convex $\Rightarrow$
        quasiconvex. A function that is both concave and convex is affine.
  \item A strictly increasing transform $F$ ($\ln$ of a positive $f$,
        $\sqrt{\cdot}$ or a positive power of a nonnegative $f$) only
        relabels the contour sets, $\{F(f)\ge F(c)\}=\{f\ge c\}$, so it
        preserves quasiconcavity and quasiconvexity but not necessarily
        concavity.
  \item If $g$ is concave and $F$ is increasing and concave, $F\circ g$ is
        concave: this is how $\sqrt{\text{concave}}$ and
        $\ln(\text{concave})$ are handled without a Hessian.
  \item One straight line in the domain along which $f$ is strictly convex
        refutes concavity; one along which it is strictly concave refutes
        convexity. Two points in a contour set whose midpoint is not in it
        refute quasiconcavity or quasiconvexity.
  \item Cobb--Douglas $x^{a}y^{b}$ ($a,b>0$) on $\R^{2}_{+}$: always
        quasiconcave; concave iff $a+b\le1$; never convex.
        (Vol.~7, Ch.~3.)
\end{itemize}
\end{rulebox}

\section{Worked Example}

\begin{worked}{Investigate $f(x,y)=\sqrt{8-2x^{2}-y^{2}}$ completely: domain, range, derivative, level curves, gradient at $(1,2)$, Hessian, critical points, contour sets, and the full property checklist}
\wstep{Step 1. Domain (Vol.~2, Ch.~1; Vol.~2, Ch.~4).}
The only obstruction is the root: $8-2x^{2}-y^{2}\ge0$, that is
\[
  D=\{(x,y): 2x^{2}+y^{2}\le8\}
   =\Bigl\{\tfrac{x^{2}}{4}+\tfrac{y^{2}}{8}\le1\Bigr\},
\]
the filled ellipse with semi-axes $2$ along $x$ and $2\sqrt2$ along $y$,
boundary included. It is \textbf{closed} (the inequality is weak),
\textbf{bounded}, and \textbf{convex}: it is a lower contour set of the
convex function $2x^{2}+y^{2}$ (Vol.~7, Ch.~2).
\wtrap{Writing $2x^{2}+y^{2}<8$ loses the boundary. The root of zero is
zero, so the boundary ellipse belongs to the domain, and it will carry the
minimum.}

\wstep{Step 2. Range (Vol.~2, Ch.~1).}
The radicand is largest at the origin, where it equals $8$, and smallest
($0$) on the boundary. Since $\sqrt{\cdot}$ is increasing,
$0\le f\le\sqrt8=2\sqrt2$. $f$ is continuous on the connected set $D$, so
every value in between is taken:
\[
  f(D)=[\,0,\,2\sqrt2\,].
\]

\wstep{Step 3. Derivative, and where it exists (Vol.~5, Ch.~10).}
By the chain rule,
\[
  f_{x}=\frac{-4x}{2\sqrt{8-2x^{2}-y^{2}}}=-\frac{2x}{f},\qquad
  f_{y}=\frac{-2y}{2\sqrt{8-2x^{2}-y^{2}}}=-\frac{y}{f}.
\]
Both are continuous wherever $f>0$, so $f$ is differentiable on the
\textbf{interior} $2x^{2}+y^{2}<8$. On the boundary the formulas divide by
zero, and the difference quotient confirms the failure: at $(2,0)$,
$\bigl(f(2-h,0)-f(2,0)\bigr)/h=\sqrt{8h-2h^{2}}/h\to\infty$ as
$h\to0^{+}$. The surface meets the plane $z=0$ vertically, like the edge
of a dome.

\wstep{Step 4. Level curves (Vol.~5, Ch.~9).}
For $0\le c\le2\sqrt2$, $f=c$ means $2x^{2}+y^{2}=8-c^{2}$: concentric
ellipses with the same shape as the boundary, shrinking to the origin as
$c\to2\sqrt2$. The values $c=0,1,\sqrt2,2,\sqrt7$ give
$2x^{2}+y^{2}=8,7,6,4,1$ (left panel below). The level $c=0$ is the
boundary of the domain; the level $c=2\sqrt2$ is the single point
$(0,0)$.

\wstep{Step 5. Gradient at $(1,2)$ (Vol.~5, Ch.~10).}
Here $f(1,2)=\sqrt{8-2-4}=\sqrt2$, so
\[
  \grad f(1,2)=\begin{pmatrix}-2/\sqrt2\\-2/\sqrt2\end{pmatrix}
  =\begin{pmatrix}-\sqrt2\\-\sqrt2\end{pmatrix},
  \qquad
  \lVert\grad f(1,2)\rVert=\sqrt{2+2}=2 .
\]
\emph{Orthogonality.} The level curve through $(1,2)$ is
$2x^{2}+y^{2}=6$; implicit differentiation gives
$y'=-4x/(2y)=-1$ there, so a tangent direction is $(1,-1)$, and
$(-\sqrt2)(1)+(-\sqrt2)(-1)=0$ \checkmark.
\emph{Direction.} The gradient points into the ellipse, towards higher
levels: the direction of steepest ascent. \emph{Norm.} $2$ is the slope of
the surface in that direction, the largest directional derivative at
$(1,2)$ (Vol.~5, Ch.~13). In every direction along the level curve the
slope is $0$.

\begin{gallery}
\begin{tikzpicture}
  \begin{axis}[contouraxis, xmin=-3.2, xmax=3.2, ymin=-3.2, ymax=3.2,
      xtick={-2,-1,1,2}, ytick={-2,-1,1,2},
      title={domain, levels, $\grad f(1,2)$}]
    \addplot[name path=top, draw=none, domain=-2:2]
      {sqrt(max(8-2*x^2,0))};
    \addplot[name path=bot, draw=none, domain=-2:2]
      {-sqrt(max(8-2*x^2,0))};
    \region{top}{bot}{}
    \addplot[accent, line width=1pt, domain=0:360, samples=121]
      ({2*cos(x)},{sqrt(8)*sin(x)});
    \addplot[domain=0:360, samples=121]({sqrt(3.5)*cos(x)},{sqrt(7)*sin(x)});
    \addplot[domain=0:360, samples=121]({sqrt(3)*cos(x)},{sqrt(6)*sin(x)});
    \addplot[domain=0:360, samples=121]({sqrt(2)*cos(x)},{2*sin(x)});
    \addplot[domain=0:360, samples=121]({sqrt(0.5)*cos(x)},{sin(x)});
    \draw[dashed, gray!70!black, line width=0.5pt]
      (axis cs:0.1,2.9) -- (axis cs:1.9,1.1);
    \draw[-{Stealth[length=6pt,width=4pt]}, accent, line width=1pt]
      (axis cs:1,2) -- (axis cs:-0.414,0.586);
    \addplot[only marks, mark=*, mark size=1.2pt, black, forget plot]
      coordinates {(1,2) (0,0)};
  \end{axis}
\end{tikzpicture}\hfill
\begin{tikzpicture}
  \begin{axis}[contouraxis, xmin=-3.2, xmax=3.2, ymin=-3.2, ymax=3.2,
      xtick={-2,-1,1,2}, ytick={-2,-1,1,2},
      title={upper set $\{f\ge2\}$ (shaded)}]
    \addplot[name path=itop, draw=none, domain=-1.4142:1.4142]
      {sqrt(max(4-2*x^2,0))};
    \addplot[name path=ibot, draw=none, domain=-1.4142:1.4142]
      {-sqrt(max(4-2*x^2,0))};
    \region{itop}{ibot}{}
    \addplot[accent, line width=1pt, domain=0:360, samples=121]
      ({2*cos(x)},{sqrt(8)*sin(x)});
    \addplot[domain=0:360, samples=121]({sqrt(2)*cos(x)},{2*sin(x)});
    \draw[gray!70!black, line width=0.5pt] (axis cs:0,2) -- (axis cs:0,-2);
    \addplot[only marks, mark=*, mark size=1.2pt, black, forget plot]
      coordinates {(0,2) (0,-2)};
    \addplot[only marks, mark=o, mark size=1.6pt, black, forget plot]
      coordinates {(0,0)};
  \end{axis}
\end{tikzpicture}
\gcap{Left: the domain (shaded), its boundary $f=0$ (bold), the levels
$f=1,\sqrt2,2,\sqrt7$, and the gradient at $(1,2)$, orthogonal to the
dashed tangent of its level curve. Right: the upper set $\{f\ge2\}$ is a
filled ellipse; the lower set $\{f\le2\}$ is the ring between the two
curves, which contains $(0,\pm2)$ but not their midpoint.}
\end{gallery}

\wstep{Step 6. Hessian (Vol.~5, Ch.~11).}
Differentiate $f_{x}=-2xf^{-1}$ and $f_{y}=-yf^{-1}$ once more, using
$f_{x}$, $f_{y}$ from Step 3:
\[
  f_{xx}=-\frac{2}{f}+\frac{2x f_{x}}{f^{2}}
        =-\frac{2f^{2}+4x^{2}}{f^{3}}
        =-\frac{2(8-y^{2})}{f^{3}},\qquad
  f_{yy}=-\frac{1}{f}+\frac{y f_{y}}{f^{2}}
        =-\frac{8-2x^{2}}{f^{3}},
\]
\[
  f_{xy}=\frac{2x f_{y}}{f^{2}}=-\frac{2xy}{f^{3}}
  \quad(\text{equal to } f_{yx}\text{ by Young}).
\]
The determinant simplifies remarkably:
\[
  \det\Hess f=\frac{2(8-y^{2})(8-2x^{2})-4x^{2}y^{2}}{f^{6}}
  =\frac{16(8-2x^{2}-y^{2})}{f^{6}}=\frac{16}{f^{4}}>0 .
\]
In the interior $y^{2}<8$, so $f_{xx}<0$: the Hessian is
\textbf{negative definite} at every interior point. At $(1,2)$, with
$f^{3}=2\sqrt2$:
\[
  \Hess f(1,2)=\begin{pmatrix}-2\sqrt2&-\sqrt2\\[2pt]-\sqrt2&-\tfrac32\sqrt2
  \end{pmatrix},\qquad \det=6-2=4=\frac{16}{(\sqrt2)^{4}}\ \checkmark
\]
\wtrap{Do not expand $\det\Hess f$ with the roots still in it. Writing
every entry over the common denominator $f^{3}$ and replacing $f^{2}$ by
$8-2x^{2}-y^{2}$ is what makes the sign visible.}

\wstep{Step 7. Critical points (Vol.~7, Ch.~5).}
$f_{x}=0\iff x=0$ and $f_{y}=0\iff y=0$, so the only candidate is
$(0,0)$, which lies in the interior. There
$\Hess f(0,0)=\begin{pmatrix}-1/\sqrt2&0\\0&-1/(2\sqrt2)\end{pmatrix}$ is
negative definite: a strict local maximum, with value $2\sqrt2$. The
boundary points are not critical points; $f$ is not differentiable there.

\wstep{Step 8. Contour sets (Vol.~5, Ch.~9; Vol.~7, Ch.~2).}
Take $c=2$. The upper set $\{f\ge2\}=\{2x^{2}+y^{2}\le4\}$ is a filled
ellipse: \textbf{convex}. The lower set
$\{f\le2\}=\{4\le2x^{2}+y^{2}\le8\}$ is an elliptical ring: it contains
$(0,2)$ and $(0,-2)$ but not their midpoint $(0,0)$, so it is
\textbf{not convex} (right panel).

\wstep{Step 9. The checklist, each with its reason.}
\begin{itemize}[leftmargin=1.6em,itemsep=0.3ex]
  \item \emph{Concave: yes.} The radicand $8-2x^{2}-y^{2}$ is concave and
        $\sqrt{\cdot}$ is increasing and concave, so the composition is
        concave on $D$ (Vol.~7, Ch.~3). Step 6 confirms it, even strictly,
        on the interior.
  \item \emph{Convex: no.} The Hessian is negative definite; concretely,
        $f(0,0)=2\sqrt2$ lies above the chord value
        $\tfrac12\bigl(f(-2,0)+f(2,0)\bigr)=0$.
  \item \emph{Quasiconcave: yes}, because concave; every upper set is a
        filled ellipse (Step 8).
  \item \emph{Quasiconvex: no.} The lower set $\{f\le2\}$ is not convex
        (Step 8).
  \item \emph{Maximum: yes}, $2\sqrt2$ at $(0,0)$, unique; global because
        $f$ is concave and $(0,0)$ is a critical point (Vol.~7, Ch.~5).
  \item \emph{Minimum: yes}, $0$, attained at \emph{every} point of the
        boundary ellipse. Weierstrass guarantees it (continuous on a closed
        bounded set); the first-order condition could never find it, since
        $f$ is not differentiable there.
  \item \emph{Increasing: no.} $(0,0)\le(1,0)$ but
        $f(1,0)=\sqrt6<2\sqrt2=f(0,0)$.
  \item \emph{Decreasing: no.} $(-1,0)\le(0,0)$ but
        $f(-1,0)=\sqrt6<f(0,0)$.
  \item \emph{Epigraph convex: no}, since $f$ is not convex. Its
        hypograph, the solid dome, is convex, matching concavity.
\end{itemize}

\wstep{Step 10. Do the answers fit together?}
Concave with an interior critical point gives a global maximum
\checkmark. Not quasiconvex, yet concave: a dome has convex upper sets and
ring-shaped lower sets \checkmark. The range $[0,2\sqrt2]$ matches the
maximum and minimum found \checkmark. Neither monotone, as one expects of
a function that peaks in the middle of its domain \checkmark.
\end{worked}

\section{Problems}

Each problem is one investigation. Work every part, and give a reason for
every item of the checklist.

\begin{problems}
  \item Investigate $f(x,y)=\sqrt{6-x-y^{2}}$.
        \begin{enumerate}[label=(\alph*),leftmargin=1.8em,itemsep=0.2ex]
          \item Give and sketch the domain. Is it closed, bounded, convex?
                What values does $f$ take?
          \item Compute $f_{x}$ and $f_{y}$. Is $f$ differentiable on its
                whole domain?
          \item Describe and sketch the level curves.
          \item Compute $\grad f(1,1)$, draw it, show that it is orthogonal
                to the level curve through $(1,1)$, and state what its
                direction and its norm mean.
          \item Compute the Hessian and its definiteness.
          \item Does $f$ have critical points?
          \item Mark an upper and a lower contour set at level $2$. Which is
                convex?
          \item Decide, with reasons: concave, convex, quasiconcave,
                quasiconvex, has a maximum, has a minimum, increasing,
                decreasing, epigraph convex.
        \end{enumerate}
  \item Investigate the utility function
        $u:\R^{2}_{++}\to\R$, $u(x,y)=\ln\bigl(x^{1/3}y^{1/2}\bigr)$, where
        $\R^{2}_{++}=\{(x,y):x>0,\ y>0\}$.
        \begin{enumerate}[label=(\alph*),leftmargin=1.8em,itemsep=0.2ex]
          \item Rewrite $u$ as a sum. Describe the domain (closed? bounded?
                convex?) and the range.
          \item Describe and sketch the level curves. How do they compare
                with those of $x^{1/3}y^{1/2}$?
          \item Compute $\grad u(1,1)$, draw it, check orthogonality to the
                level curve, and interpret direction and norm.
          \item Compute the Hessian. Show that both marginal utilities are
                diminishing, and give the definiteness.
          \item Does $u$ have critical points?
          \item Mark an upper and a lower contour set at level $0$. Which is
                convex?
          \item Decide, with reasons: concave, convex, quasiconcave,
                quasiconvex, has a maximum, has a minimum, increasing,
                decreasing, epigraph convex.
        \end{enumerate}
  \item Investigate the production function $f:\R^{2}_{+}\to\R$,
        $f(x,y)=x^{3/4}y^{1/2}$.
        \begin{enumerate}[label=(\alph*),leftmargin=1.8em,itemsep=0.2ex]
          \item Name this type of function and its returns to scale. Give
                the domain (closed? convex?) and the range.
          \item Compute the derivative and say where it exists. Show that
                both marginal products are diminishing.
          \item Compute the Hessian and its definiteness. Does diminishing
                marginal productivity make $f$ concave?
          \item Sketch the level curves, and draw $\grad f(1,1)$; check
                orthogonality and interpret direction and norm.
          \item Does $f$ have critical points? Where does its minimum come
                from?
          \item Mark an upper and a lower contour set at level $1$. Which is
                convex?
          \item Decide, with reasons: concave, convex, quasiconcave,
                quasiconvex, has a maximum, has a minimum, increasing,
                decreasing, epigraph convex.
        \end{enumerate}
  \item Investigate $f:\R^{2}_{+}\to\R$, $f(x,y)=x^{1/3}y^{2/3}$.
        \begin{enumerate}[label=(\alph*),leftmargin=1.8em,itemsep=0.2ex]
          \item Name the function and its returns to scale; give domain and
                range.
          \item Compute the derivative and the Hessian, and the
                definiteness. Which direction does the Hessian annihilate?
          \item Sketch the level curves and draw $\grad f(1,8)$; check
                orthogonality, interpret direction and norm, and verify
                Euler's theorem at $(1,8)$.
          \item Does $f$ have critical points?
          \item Mark an upper and a lower contour set at level $4$. Which is
                convex?
          \item Decide, with reasons: concave, strictly concave, convex,
                quasiconcave, quasiconvex, has a maximum, has a minimum,
                increasing, decreasing, epigraph convex.
        \end{enumerate}
  \item Investigate $f(x,y)=\ln\bigl(4+x^{2}+y^{2}\bigr)$.
        \begin{enumerate}[label=(\alph*),leftmargin=1.8em,itemsep=0.2ex]
          \item Give the domain and the range.
          \item Compute the gradient; where is $f$ differentiable?
          \item Describe the level curves, compute $\grad f(1,1)$, and check
                that it is orthogonal to the level curve through $(1,1)$.
          \item Compute the Hessian, its determinant in general, and its
                definiteness at $(0,0)$, $(1,1)$ and $(2,2)$.
          \item Find and classify the critical points.
          \item Mark the upper and lower contour sets at level $\ln8$.
                Which is convex?
          \item Decide, with reasons: concave, convex, quasiconcave,
                quasiconvex, has a maximum, has a minimum, increasing,
                decreasing, epigraph convex. On which disk centred at the
                origin is $f$ convex?
        \end{enumerate}
\end{problems}

\section{Solutions}

\begin{solutions}
  \item \textbf{Root of a concave expression on an unbounded closed domain.}
        \begin{enumerate}[label=(\alph*),leftmargin=1.8em,itemsep=0.4ex]
          \item \emph{Domain (Vol.~2, Ch.~1 and Ch.~4).} Radicand
                $\ge0$: $D=\{x\le6-y^{2}\}$, the region on and to the left of
                the parabola $x=6-y^{2}$ (vertex $(6,0)$, opening to the
                left). Closed (weak inequality), unbounded (it contains the
                whole ray $y=0$, $x\le6$), and convex: it is
                the upper set $\{g\ge0\}$ of the concave function
                $g=6-x-y^{2}$ (Vol.~7, Ch.~2). \emph{Range
                (Vol.~2, Ch.~1):} $f=0$ on the parabola and
                $f(x,0)=\sqrt{6-x}\to\infty$ as $x\to-\infty$; by
                continuity $f(D)=[0,\infty)$.
          \item \emph{Derivative (Vol.~5, Ch.~10).}
                $f_{x}=\dfrac{-1}{2\sqrt{6-x-y^{2}}}=-\dfrac{1}{2f}$,
                \;$f_{y}=\dfrac{-2y}{2f}=-\dfrac{y}{f}$, continuous where
                $f>0$. So $f$ is differentiable on the interior
                $x<6-y^{2}$ only. On the parabola the partials blow up:
                at $(6,0)$, $\bigl(f(6-h,0)-0\bigr)/h=\sqrt h/h\to\infty$.
          \item \emph{Level curves (Vol.~5, Ch.~9).} $f=c\ge0$ gives
                $x=6-c^{2}-y^{2}$: copies of the boundary parabola, shifted
                left by $c^{2}$. For $c=0,1,2,3$ the vertices sit at
                $x=6,5,2,-3$; higher levels lie further left.
          \item \emph{Gradient (Vol.~5, Ch.~10).} $f(1,1)=\sqrt4=2$, so
                \[
                  \grad f(1,1)=\begin{pmatrix}-1/4\\-1/2\end{pmatrix},\qquad
                  \lVert\grad f(1,1)\rVert=\sqrt{\tfrac1{16}+\tfrac14}
                  =\frac{\sqrt5}{4}\approx0.56 .
                \]
                The level curve through $(1,1)$ is $x=2-y^{2}$; there
                $\mathrm{d}x/\mathrm{d}y=-2y=-2$, so
                $\mathrm{d}y/\mathrm{d}x=-\tfrac12$ and $(2,-1)$ is a tangent
                direction. $(-\tfrac14)(2)+(-\tfrac12)(-1)=0$ \checkmark.
                The gradient points left and down, towards the inside of the
                parabolas where the levels rise: steepest ascent, with slope
                $\sqrt5/4$ in that direction (Vol.~5, Ch.~13).
          \item \emph{Hessian (Vol.~5, Ch.~11).} From $f_{x}=-\tfrac12f^{-1}$
                and $f_{y}=-yf^{-1}$:
                \[
                  f_{xx}=\tfrac12f^{-2}f_{x}=-\frac{1}{4f^{3}},\qquad
                  f_{xy}=\tfrac12f^{-2}f_{y}=-\frac{y}{2f^{3}},
                \]
                \[
                  f_{yy}=-\frac1f+\frac{yf_{y}}{f^{2}}
                        =-\frac{f^{2}+y^{2}}{f^{3}}=-\frac{6-x}{f^{3}} .
                \]
                $\det\Hess f=\dfrac{(6-x)-y^{2}}{4f^{6}}
                =\dfrac{f^{2}}{4f^{6}}=\dfrac{1}{4f^{4}}>0$ and
                $f_{xx}<0$: \textbf{negative definite} on the interior. At
                $(1,1)$: $\Hess f=\begin{pmatrix}-1/32&-1/16\\-1/16&-5/8
                \end{pmatrix}$, $\det=\tfrac{5}{256}-\tfrac{1}{256}
                =\tfrac1{64}=\tfrac{1}{4\cdot2^{4}}$ \checkmark.
          \item \emph{Critical points (Vol.~7, Ch.~5).} None:
                $f_{x}=-1/(2f)$ is never $0$. Moving left always raises $f$.
          \item \emph{Contour sets (Vol.~5, Ch.~9; Vol.~7, Ch.~2).}
                $\{f\ge2\}=\{x\le2-y^{2}\}$, the inside of a parabola:
                \textbf{convex}. $\{f\le2\}=\{2-y^{2}\le x\le6-y^{2}\}$, a
                curved band between two parabolas: it contains $(-2,2)$ and
                $(-2,-2)$ (check: $-2\le-2\le2$) but not their midpoint
                $(-2,0)$ (since $2\le-2$ fails). \textbf{Not convex.}
          \item \emph{Checklist (Vol.~7, Ch.~1 and Ch.~3).}
                \emph{Concave: yes}, $\sqrt{\cdot}$ increasing and concave
                applied to the concave $6-x-y^{2}$ (and (e) confirms it).
                \emph{Convex: no}, negative definite Hessian.
                \emph{Quasiconcave: yes}, because concave.
                \emph{Quasiconvex: no}, by the band in (g).
                \emph{Maximum: no}, $f$ is unbounded above.
                \emph{Minimum: yes}, $0$ on the whole parabola $x=6-y^{2}$,
                boundary points where $f$ is not differentiable.
                \emph{Increasing: no}, $f$ falls in $x$: $(0,0)\le(1,0)$
                with $f(1,0)=\sqrt5<\sqrt6=f(0,0)$.
                \emph{Decreasing: no}, $(0,-1)\le(0,0)$ with
                $f(0,-1)=\sqrt5<\sqrt6=f(0,0)$.
                \emph{Epigraph convex: no}, since $f$ is not convex (the
                hypograph is).
        \end{enumerate}
  \item \textbf{Log-Cobb--Douglas.}
        \begin{enumerate}[label=(\alph*),leftmargin=1.8em,itemsep=0.4ex]
          \item \emph{Logarithm rules (Vol.~1, Ch.~2).}
                $u=\tfrac13\ln x+\tfrac12\ln y$. \emph{Domain
                (Vol.~5, Ch.~10):} the open quadrant $\R^{2}_{++}$, open
                (not closed), unbounded, convex. \emph{Range:}
                $u(t,1)=\tfrac13\ln t$ runs through all of $\R$ as $t$
                runs through $(0,\infty)$, so $u(\R^{2}_{++})=\R$.
          \item \emph{Level curves (Vol.~5, Ch.~9).}
                $\tfrac13\ln x+\tfrac12\ln y=c\iff\ln y=2c-\tfrac23\ln x
                \iff y=\mathrm{e}^{2c}x^{-2/3}$: decreasing, convex curves
                asymptotic to both axes, moving outwards as $c$ grows. They
                are exactly the level curves of $x^{1/3}y^{1/2}$, relabelled
                ($u=\ln v$ with $v$ the Cobb--Douglas): an increasing
                transform changes the labels, not the curves.
          \item \emph{Gradient (Vol.~5, Ch.~10).} $u_{x}=\dfrac{1}{3x}$,
                $u_{y}=\dfrac{1}{2y}$, so $\grad u(1,1)=(\tfrac13,\tfrac12)$
                with norm $\sqrt{\tfrac19+\tfrac14}=\tfrac{\sqrt{13}}{6}
                \approx0.60$. The level curve through $(1,1)$ is $u=0$, i.e.
                $y=x^{-2/3}$, with slope $-\tfrac23$ at $x=1$ and tangent
                direction $(3,-2)$; $\tfrac13\cdot3-\tfrac12\cdot2=0$
                \checkmark. The gradient points up and to the right, the
                direction in which utility rises fastest, at rate
                $\sqrt{13}/6$.
          \item \emph{Hessian (Vol.~5, Ch.~11).}
                $\Hess u=\begin{pmatrix}-\frac{1}{3x^{2}}&0\\0&-\frac{1}{2y^{2}}
                \end{pmatrix}$. $u_{xx}<0$ and $u_{yy}<0$: each marginal
                utility $u_{x}=1/(3x)$, $u_{y}=1/(2y)$ falls as its own good
                increases. Diagonal with negative entries: \textbf{negative
                definite} everywhere.
          \item \emph{Critical points (Vol.~7, Ch.~5).} None: $u_{x}=1/(3x)>0$.
          \item \emph{Contour sets.} $u\le0\iff x^{1/3}y^{1/2}\le1\iff
                x^{2}y^{3}\le1$. The upper set $\{u\ge0\}=\{y\ge x^{-2/3}\}$
                lies above a convex curve: \textbf{convex}. The lower set is
                not: $(8,\tfrac14)$ and $(\tfrac18,4)$ both satisfy
                $x^{2}y^{3}=1$, but at the midpoint
                $(\tfrac{65}{16},\tfrac{17}{8})$ the product is far above
                $1$ (already $x^{2}>16$ and $y^{3}>8$).
          \item \emph{Checklist (Vol.~7, Ch.~3).}
                \emph{Concave: yes}, strictly, by the negative definite
                Hessian (or as a positive sum of the concave $\ln x$,
                $\ln y$).
                \emph{Convex: no}, the Hessian is negative definite.
                \emph{Quasiconcave: yes}, because concave.
                \emph{Quasiconvex: no}, by (f).
                \emph{Maximum: no}, $u(t,1)\to\infty$ as $t\to\infty$.
                \emph{Minimum: no}, $u(t,1)\to-\infty$ as $t\to0^{+}$; the
                open domain has no boundary to stop at.
                \emph{Increasing: yes}, strictly: $u_{x}>0$, $u_{y}>0$ on a
                convex domain.
                \emph{Decreasing: no}, as it is strictly increasing.
                \emph{Epigraph convex: no}, $u$ is not convex.
        \end{enumerate}
  \item \textbf{Cobb--Douglas with exponent sum above $1$.}
        \begin{enumerate}[label=(\alph*),leftmargin=1.8em,itemsep=0.4ex]
          \item \emph{Name (Vol.~5, Ch.~14).} A Cobb--Douglas function,
                homogeneous of degree $\tfrac34+\tfrac12=\tfrac54>1$:
                $f(tx,ty)=t^{5/4}f(x,y)$, \textbf{increasing returns to
                scale}. Domain $\R^{2}_{+}$: closed, unbounded, convex.
                Range $[0,\infty)$: $f=0$ on the axes and $f(t,t)=t^{5/4}$
                is unbounded.
          \item \emph{Derivative (Vol.~5, Ch.~10).}
                \[
                  f_{x}=\tfrac34x^{-1/4}y^{1/2},\qquad
                  f_{y}=\tfrac12x^{3/4}y^{-1/2},
                \]
                continuous on the open quadrant $x,y>0$, so $f$ is
                differentiable there. On the axis $x=0$ (with $y>0$) the
                factor $x^{-1/4}$ is infinite, and likewise $y^{-1/2}$ on the
                axis $y=0$ (with $x>0$): the partials do not exist there.
                The origin is the exception: $f$ vanishes on both axes, so
                $f_{x}(0,0)=f_{y}(0,0)=0$, and
                $0\le f(h,k)\le\lVert(h,k)\rVert^{5/4}$ shows that $f$ is
                differentiable there (from inside the quadrant) with zero
                derivative. Second partials:
                $f_{xx}=-\tfrac{3}{16}x^{-5/4}y^{1/2}<0$,
                $f_{yy}=-\tfrac14x^{3/4}y^{-3/2}<0$: both marginal products
                \textbf{diminish}.
          \item \emph{Hessian (Vol.~5, Ch.~11).}
                $f_{xy}=\tfrac38x^{-1/4}y^{-1/2}$, and
                \[
                  \det\Hess f=\tfrac{3}{64}x^{-1/2}y^{-1}
                  -\tfrac{9}{64}x^{-1/2}y^{-1}
                  =-\tfrac{3}{32}\,x^{-1/2}y^{-1}<0 :
                \]
                \textbf{indefinite} at every interior point, so $f$ is
                neither concave nor convex. Diminishing marginal products
                (negative diagonal) are \emph{not} enough for concavity;
                the cross effect $f_{xy}$ is too strong. Along the diagonal
                the function is $f(t,t)=t^{5/4}$, strictly convex in $t$,
                which refutes concavity directly.
          \item \emph{Level curves and gradient (Vol.~5, Ch.~9, Ch.~10).}
                $x^{3/4}y^{1/2}=c$ gives $y=c^{2}x^{-3/2}$: decreasing convex
                isoquants. $\grad f(1,1)=(\tfrac34,\tfrac12)$, norm
                $\sqrt{\tfrac{9}{16}+\tfrac{4}{16}}=\tfrac{\sqrt{13}}{4}
                \approx0.90$. The isoquant $y=x^{-3/2}$ through $(1,1)$ has
                slope $-\tfrac32$, tangent $(2,-3)$, and
                $\tfrac34\cdot2-\tfrac12\cdot3=0$ \checkmark. The gradient
                points to the upper right, towards higher output, and
                $\sqrt{13}/4$ is the output gain per unit step in that best
                direction.
          \item \emph{Critical points (Vol.~7, Ch.~5).} None in the
                interior, since $f_{x}>0$. The minimum value $0$ is attained
                on the two axes, boundary points of the domain, which the
                first-order condition does not see. (At the corner $(0,0)$
                the gradient from (b) does vanish, but the other minimizers
                on the axes are points where the partials do not even
                exist.)
          \item \emph{Contour sets.} $\{f\ge1\}=\{y\ge x^{-3/2}\}$, the
                region above a convex curve: \textbf{convex}. $\{f\le1\}$ is
                not: $(4,\tfrac18)$ and $(\tfrac14,8)$ lie on $f=1$, but the
                midpoint $(\tfrac{17}{8},\tfrac{65}{16})$ has
                $f>1$ (both coordinates exceed $1$).
          \item \emph{Checklist (Vol.~7, Ch.~3).}
                \emph{Concave: no}, $f(t,t)=t^{5/4}$ is strictly convex.
                \emph{Convex: no}, $f(t,1)=t^{3/4}$ is strictly concave.
                \emph{Quasiconcave: yes}, all upper sets lie above convex
                isoquants (Cobb--Douglas always is).
                \emph{Quasiconvex: no}, by (f).
                \emph{Maximum: no}, unbounded.
                \emph{Minimum: yes}, $0$ on the axes.
                \emph{Increasing: yes}: if $x\le x'$ and $y\le y'$ then
                $x^{3/4}\le x'^{3/4}$ and $y^{1/2}\le y'^{1/2}$, and the
                product of nonnegative factors grows.
                \emph{Decreasing: no}, $f(0,0)=0<1=f(1,1)$.
                \emph{Epigraph convex: no}, $f$ is not convex.
        \end{enumerate}
  \item \textbf{Cobb--Douglas with constant returns.}
        \begin{enumerate}[label=(\alph*),leftmargin=1.8em,itemsep=0.4ex]
          \item \emph{Name (Vol.~5, Ch.~14).} Cobb--Douglas, homogeneous of
                degree $\tfrac13+\tfrac23=1$: \textbf{constant returns to
                scale}. Domain $\R^{2}_{+}$ (closed, convex, unbounded),
                range $[0,\infty)$, since $f(t,t)=t$.
          \item \emph{Derivative and Hessian (Vol.~5, Ch.~10, Ch.~11).}
                On $x,y>0$:
                $f_{x}=\tfrac13x^{-2/3}y^{2/3}$,
                $f_{y}=\tfrac23x^{1/3}y^{-1/3}$, and
                \[
                  \Hess f=\frac29\begin{pmatrix}
                    -x^{-5/3}y^{2/3}& x^{-2/3}y^{-1/3}\\
                    x^{-2/3}y^{-1/3}& -x^{1/3}y^{-4/3}\end{pmatrix},\qquad
                  \det\Hess f=\tfrac{4}{81}x^{-4/3}y^{-2/3}
                  -\tfrac{4}{81}x^{-4/3}y^{-2/3}=0 .
                \]
                Diagonal negative, determinant zero: \textbf{negative
                semidefinite}, not definite. The null direction is the ray
                itself: the first row times $(x,y)$ is
                $\tfrac29\bigl(-x^{-2/3}y^{2/3}+x^{-2/3}y^{2/3}\bigr)=0$,
                and likewise the second. Along a ray $f(tx,ty)=tf(x,y)$ is
                linear, so it has no curvature there.
          \item \emph{Level curves, gradient, Euler (Vol.~5, Ch.~9,
                Ch.~10, Ch.~14).} $x^{1/3}y^{2/3}=c\iff y=c^{3/2}x^{-1/2}$.
                $f(1,8)=1\cdot4=4$, so the isoquant is $y=8x^{-1/2}$, slope
                $-4x^{-3/2}=-4$ at $x=1$. $\grad f(1,8)=
                (\tfrac13\cdot4,\ \tfrac23\cdot\tfrac12)=(\tfrac43,\tfrac13)$,
                norm $\tfrac{\sqrt{17}}{3}\approx1.37$; with tangent
                $(1,-4)$: $\tfrac43-\tfrac43=0$ \checkmark. Steepest ascent
                points mostly in the $x$ direction, because $x$ is the
                scarce factor at $(1,8)$. Euler: $1\cdot\tfrac43+8\cdot
                \tfrac13=4=f(1,8)$ \checkmark.
          \item \emph{Critical points.} None in the interior ($f_{x}>0$);
                the minimum $0$ sits on the axes.
          \item \emph{Contour sets.} $\{f\ge4\}=\{y\ge8x^{-1/2}\}$:
                \textbf{convex}. $\{f\le4\}$: $(16,2)$ and $(\tfrac14,16)$
                lie on $f=4$, the midpoint $(\tfrac{65}{8},9)$ gives
                $f=(\tfrac{65}{8})^{1/3}\cdot9^{2/3}\approx8.7>4$.
                \textbf{Not convex.}
          \item \emph{Checklist (Vol.~7, Ch.~3).}
                \emph{Concave: yes}, the Hessian is negative semidefinite
                on the interior, and concavity carries over to the axes by
                continuity (Cobb--Douglas with $a+b\le1$).
                \emph{Strictly concave: no}, $f(t,t)=t$ is linear along the
                diagonal.
                \emph{Convex: no}, $f(t,1)=t^{1/3}$ is strictly concave.
                \emph{Quasiconcave: yes}, because concave.
                \emph{Quasiconvex: no}, by (e).
                \emph{Maximum: no}. \emph{Minimum: yes}, $0$ on the axes.
                \emph{Increasing: yes}, as in Problem~3.
                \emph{Decreasing: no}.
                \emph{Epigraph convex: no}; the hypograph is convex.
        \end{enumerate}
  \item \textbf{Logarithm of a positive quadratic.}
        \begin{enumerate}[label=(\alph*),leftmargin=1.8em,itemsep=0.4ex]
          \item \emph{Domain and range (Vol.~2, Ch.~1).}
                $q=4+x^{2}+y^{2}\ge4>0$ always, so $D=\R^{2}$. $q$ takes
                every value in $[4,\infty)$ and $\ln$ is increasing, so
                $f(\R^{2})=[\ln4,\infty)$.
          \item \emph{Gradient (Vol.~5, Ch.~10).}
                $\grad f=\Bigl(\dfrac{2x}{q},\dfrac{2y}{q}\Bigr)$,
                continuous everywhere since $q\ne0$: differentiable on all
                of $\R^{2}$.
          \item \emph{Level curves and gradient (Vol.~5, Ch.~9).}
                $f=c\iff x^{2}+y^{2}=\mathrm{e}^{c}-4$: circles about the
                origin for $c>\ln4$, the origin itself for $c=\ln4$, empty
                below. Through $(1,1)$ runs the circle $x^{2}+y^{2}=2$
                (level $\ln6$). $\grad f(1,1)=(\tfrac26,\tfrac26)
                =(\tfrac13,\tfrac13)$, norm $\tfrac{\sqrt2}{3}$; it is a
                multiple of the radius vector $(1,1)$, and a radius is
                orthogonal to a circle: tangent $(1,-1)$,
                $\tfrac13-\tfrac13=0$ \checkmark. It points outwards,
                towards larger circles and higher values.
          \item \emph{Hessian (Vol.~5, Ch.~11).} Quotient rule:
                \[
                  f_{xx}=\frac{2(4-x^{2}+y^{2})}{q^{2}},\quad
                  f_{yy}=\frac{2(4+x^{2}-y^{2})}{q^{2}},\quad
                  f_{xy}=-\frac{4xy}{q^{2}},
                \]
                \[
                  \det\Hess f=\frac{4\bigl[16-(x^{2}-y^{2})^{2}\bigr]
                  -16x^{2}y^{2}}{q^{4}}
                  =\frac{4\bigl[16-(x^{2}+y^{2})^{2}\bigr]}{q^{4}} .
                \]
                At $(0,0)$: $\Hess f=\mathrm{diag}(\tfrac12,\tfrac12)$,
                positive definite. At $(1,1)$ ($q=6$):
                $\begin{pmatrix}2/9&-1/9\\-1/9&2/9\end{pmatrix}$,
                $\det=\tfrac{1}{27}>0$, $f_{xx}>0$: positive definite. At
                $(2,2)$ ($q=12$):
                $\begin{pmatrix}1/18&-1/9\\-1/9&1/18\end{pmatrix}$,
                $\det=-\tfrac{1}{108}<0$: indefinite.
          \item \emph{Critical points (Vol.~7, Ch.~5).} $2x/q=0$ and
                $2y/q=0$ give only $(0,0)$; positive definite Hessian: strict
                local minimum $\ln4$, and global, since $f\ge\ln4$ by (a).
          \item \emph{Contour sets.} $\{f\le\ln8\}=\{x^{2}+y^{2}\le4\}$, a
                disk: \textbf{convex}. $\{f\ge\ln8\}=\{x^{2}+y^{2}\ge4\}$,
                the outside of the disk: it contains $(2,0)$ and $(-2,0)$
                but not their midpoint $(0,0)$. \textbf{Not convex.}
          \item \emph{Checklist (Vol.~7, Ch.~3).}
                \emph{Concave: no}, the Hessian is positive definite at the
                origin.
                \emph{Convex: no}, it is indefinite at $(2,2)$; along the
                $x$-axis $g(t)=\ln(4+t^{2})$ has
                $g''(t)=\dfrac{2(4-t^{2})}{(4+t^{2})^{2}}<0$ for $|t|>2$.
                \emph{Quasiconvex: yes}, every lower set is a disk (or
                empty, or a point).
                \emph{Quasiconcave: no}, by (f).
                \emph{Maximum: no}, unbounded above.
                \emph{Minimum: yes}, $\ln4$ at $(0,0)$.
                \emph{Increasing: no}, $(-1,0)\le(0,0)$ but
                $f(-1,0)=\ln5>\ln4=f(0,0)$.
                \emph{Decreasing: no}, $(0,0)\le(1,0)$ but
                $f(1,0)=\ln5>f(0,0)$.
                \emph{Epigraph convex: no}, $f$ is not convex.
                \emph{Convex on a disk:} the determinant is positive
                exactly when $x^{2}+y^{2}<2^{2}$, and there
                $f_{xx}=2(4-x^{2}+y^{2})/q^{2}>0$ because $x^{2}<4$. So $f$
                is convex on the disk of radius $2$ (closed, by continuity),
                and on no larger disk, since the determinant turns negative
                just outside it. This is the typical shape of a
                quasiconvex but not convex function: a bowl at the bottom
                that flattens into a trumpet.
        \end{enumerate}
\end{solutions}


\end{document}
"
\ifdefined\KindleEdition
  \documentclass[10pt,oneside]{book}
\else
  \documentclass[11pt,a4paper,twoside,openright]{book}
\fi
% Shared preamble for the Mixed Problems training sheet.
\usepackage[T1]{fontenc}
\usepackage[utf8]{inputenc}
\usepackage{lmodern}
\usepackage[a4paper,margin=2.4cm]{geometry}
\usepackage{amsmath,amssymb,amsthm}
\usepackage{microtype}
\usepackage{enumitem}
\usepackage{xcolor}
\usepackage[most]{tcolorbox}
\usepackage{fancyhdr}
\usepackage[hidelinks]{hyperref}

\definecolor{accent}{HTML}{1F4E79}
\definecolor{soft}{HTML}{F4F7FA}

% --- function graphs ------------------------------------------------------
\usepackage{pgfplots}
\pgfplotsset{compat=1.18}
\usetikzlibrary{arrows.meta}

% One consistent look for every function graph in these books.
\pgfplotsset{
  fnaxis/.style={
    width=5.9cm, height=4.6cm,
    axis lines=middle,
    axis line style={-{Stealth[length=4pt]}, gray!65!black, line width=0.3pt},
    tick style={gray!65!black},
    tick label style={font=\tiny, gray!50!black},
    label style={font=\scriptsize},
    xlabel={$x$}, ylabel={$y$},
    every axis x label/.style={at={(ticklabel* cs:1.0)}, anchor=west},
    every axis y label/.style={at={(ticklabel* cs:1.0)}, anchor=south east,
      xshift=-1pt},
    every axis plot/.append style={accent, very thick, smooth, samples=201},
    clip=true, enlargelimits=false,
    title style={font=\scriptsize, yshift=4pt},
  },
  fnaxis/.append style={/pgfplots/legend style={font=\tiny, draw=none,
    fill=none, at={(0.98,0.03)}, anchor=south east}},
}

% A row of small graphs, evenly spaced with \hfill between them, plus an
% optional caption line: \begin{gallery} <pic>\hfill<pic> \gcap{...} \end{gallery}
\newenvironment{gallery}
  {\par\medskip\noindent\ignorespaces}
  {\par\medskip\noindent}
\newcommand{\gcap}[1]{\par\smallskip{\centering\footnotesize #1\par}}

% --- plane regions, implicit curves, surfaces -----------------------------
\usepgfplotslibrary{fillbetween}
\definecolor{regionfill}{HTML}{9DB4D8}
\pgfplotsset{colormap={house}{color=(soft) color=(accent!45) color=(accent)}}

\pgfplotsset{
  % Squared, gridded axis for plane regions and implicitly given curves.
  geoaxis/.style={
    fnaxis,
    width=5.9cm, height=5.9cm,
    grid=both,
    grid style={gray!22, line width=0.2pt},
    minor tick num=1,
  },
  % Level-curve maps: many thin curves, no single curve emphasised.
  contouraxis/.style={
    geoaxis,
    every axis plot/.append style={accent!75, line width=0.5pt, smooth,
      samples=201},
  },
  % Surfaces: \addplot3[surf, domain=..., y domain=...] {expr};
  surfaxis/.style={
    width=6.0cm, height=5.0cm,
    view={130}{22},
    colormap name=house,
    ticks=none,
    xlabel={$x$}, ylabel={$y$}, zlabel={$f$},
    label style={font=\tiny, gray!50!black},
    title style={font=\scriptsize},
    every axis plot/.append style={samples=25, shader=faceted interp,
      draw=accent!35, line width=0.1pt},
  },
}

% Shorthand for a shaded region between two named paths.
% \region{<path A>}{<path B>}{<soft clip options>}
\newcommand{\region}[3]{%
  \addplot[regionfill, opacity=0.55, forget plot, sharp plot]
    fill between[of=#1 and #2, #3];}

% --- rule boxes -----------------------------------------------------------
\newtcolorbox{rulebox}[1]{%
  breakable, enhanced, colback=soft, colframe=accent,
  boxrule=0.4pt, left=7pt, right=7pt, top=6pt, bottom=6pt,
  fonttitle=\bfseries, coltitle=white, colbacktitle=accent,
  title={#1}}

% --- worked examples ------------------------------------------------------
% One fully narrated solution per chapter, carrying the chapter's central
% technique and the mistakes that technique invites:
%   \begin{worked}{Task} ... \wstep{Step 1.} ... \wtrap{...} ... \end{worked}
\newtcolorbox{worked}[1]{%
  breakable, enhanced, colback=white, colframe=accent!65,
  boxrule=0.4pt, left=7pt, right=7pt, top=6pt, bottom=6pt,
  fonttitle=\bfseries, coltitle=accent, colbacktitle=soft,
  title={#1}}

% Run-in step heading inside a worked example.
\newcommand{\wstep}[1]{\par\smallskip\noindent\textbf{#1}\enspace\ignorespaces}

% A pitfall, set off from the running solution so it can be found again.
\newcommand{\wtrap}[1]{%
  \par\smallskip
  {\leftskip=1.3em\noindent
   \llap{\textcolor{accent}{\rule[0.22ex]{0.9em}{0.6pt}}\hspace{0.4em}}%
   \textbf{Trap.}\enspace\itshape #1\par}%
  \smallskip}

% --- problem / solution lists --------------------------------------------
\newlist{problems}{enumerate}{1}
\setlist[problems]{label=\textbf{\arabic*.}, leftmargin=2.1em,
  itemsep=0.55ex, parsep=0pt, topsep=0.8ex}
\newlist{solutions}{enumerate}{1}
\setlist[solutions]{label=\textbf{\arabic*.}, leftmargin=2.1em,
  itemsep=1.0ex, parsep=0pt, topsep=0.8ex}

% --- running heads --------------------------------------------------------
\pagestyle{fancy}
\fancyhf{}
\fancyhead[LE,RO]{\thepage}
\fancyhead[RE]{\small\itshape\nouppercase{\leftmark}}
\fancyhead[LO]{\small\itshape\nouppercase{\rightmark}}
\setlength{\headheight}{14pt}
\renewcommand{\headrulewidth}{0.3pt}
\renewcommand{\chaptermark}[1]{\markboth{\chaptername\ \thechapter.\ #1}{}}
\renewcommand{\sectionmark}[1]{\markright{\thesection\ #1}}

% --- shorthands used throughout -------------------------------------------
\newcommand{\dd}[2]{\frac{\mathrm{d}#1}{\mathrm{d}#2}}
\newcommand{\pd}[2]{\frac{\partial #1}{\partial #2}}
\newcommand{\pdd}[3]{\frac{\partial^{2} #1}{\partial #2\,\partial #3}}
\newcommand{\R}{\mathbb{R}}
\newcommand{\grad}{\nabla}
\DeclareMathOperator{\sgn}{sgn}
\DeclareMathOperator{\Hess}{H}

% German gloss, used on first mention of a technical term
\newcommand{\de}[1]{\emph{#1}}

\renewcommand*{\thisvolume}{mixed}

\title{\Huge Mixed Problems\\[0.4em]
  \large Cross-Chapter Sets with Worked Solutions}
\author{}
\date{}

\begin{document}
\frontmatter
\maketitle

\chapter*{About this sheet}
\addcontentsline{toc}{chapter}{About this sheet}

Every other volume in this series drills one technique at a time. That is
the right way to learn a method and the wrong way to prepare for an
examination, because a chapter heading tells you what to do before you have
read the question. An exam does not. The single most common way to lose
marks is not botched arithmetic but reaching for the wrong tool: setting up
a Lagrangian where the constraint does not bind, differentiating where a
limit was wanted, solving a system by inverting a matrix that turns out to
be singular.

This sheet exists to remove that crutch. Each set draws on the whole
sequence --- algebra and functions, matrices and elimination, limits and
L'H\^opital, curve discussion, multivariable calculus, constrained
optimization, linear programming --- in no announced order, and no problem
says which chapter it came from. Deciding that \emph{is} the exercise.

Work a set in one sitting, with a time limit and without looking anything
up. Sixty to ninety minutes for twenty problems is a realistic target. Then
read the solutions, and for every problem you got wrong, note not the
arithmetic but the \emph{decision}: which signal in the question should have
pointed at which method? That note is worth more than redoing the
computation.

Solutions name the technique in their first words, so the sets can also be
used the other way round --- read only the first line of each solution as a
drill in classification.

The fourth set works the other way. It announces the methods, all of them,
and takes a single function of two variables through the whole chain:
domain, range, derivative, level curves, gradient, Hessian, critical
points, contour sets and a checklist of properties. There the difficulty is
not choosing a tool but keeping a dozen answers consistent with one
another.

\tableofcontents

\mainmatter
\chapter{Mixed Set A}

Twenty problems, no order, no announced methods. Target time: 75 minutes.

\section{Worked Example}

The problems below announce no method. That is the point: on an
examination, choosing the technique is most of the work, and it is the one
step no drill chapter can train, because a drill chapter has already told
you the answer in its title.

The example here is \emph{not} one of the twenty. It is a typical example of
the kind, a single function examined from five directions, shown so
that you can watch the recognition step happen before you have to do it
yourself.

\begin{worked}{A five-part question on $f(x)=\dfrac{\ln x}{x}$: domain, both edges, extremum, inflection, and $\int_{1}^{\infty}f$}
\wstep{Step 1. Read the parts and name the tools before computing anything.}
\begin{center}\footnotesize
\begin{tabular}{lll}
  \hline
  part & tool & volume\\
  \hline
  domain & obstruction list & 2\\
  edges & limits, growth hierarchy, L'H\^opital & 4, 5\\
  extremum & $f'=0$, then $f''$ & 5\\
  inflection & sign change of $f''$ & 5\\
  convergence & improper integral, substitution & 6\\
  \hline
\end{tabular}
\end{center}
A minute spent on this table keeps the rest of the solution organised.

\wstep{Step 2. Domain.}
$\ln x$ needs $x>0$, and the denominator needs $x\neq0$. The first
condition subsumes the second:
\[
  D_{f}=(0,\infty).
\]

\wstep{Step 3. The left edge --- and it is \emph{not} indeterminate.}
As $x\to0^{+}$: $\ln x\to-\infty$ and $1/x\to+\infty$. The form is
$(-\infty)\cdot(+\infty)$, which is determinate:
\[
  \lim_{x\to0^{+}}\frac{\ln x}{x}=-\infty .
\]
\wtrap{Reaching for L'H\^opital here is the standard reflex and it is wrong
--- the rule needs $0/0$ or $\infty/\infty$, and this is neither. Always
name the form before choosing a tool. Here both factors push the same way,
so there is nothing to resolve.}

\wstep{Step 4. The right edge --- and here it \emph{is} indeterminate.}
As $x\to\infty$ the form is $\infty/\infty$, so either L'H\^opital or the
growth hierarchy applies:
\[
  \lim_{x\to\infty}\frac{\ln x}{x}
  \overset{\infty/\infty}{=}\lim_{x\to\infty}\frac{1/x}{1}=0 ,
\]
consistent with $\ln x\ll x$. So $y=0$ is a horizontal asymptote.

\wstep{Step 5. The extremum.}
\[
  f'(x)=\frac{\tfrac1x\cdot x-\ln x}{x^{2}}=\frac{1-\ln x}{x^{2}} .
\]
$f'=0$ when $\ln x=1$, i.e.\ $x=\mathrm{e}$, and
$f(\mathrm{e})=1/\mathrm{e}\approx0.368$.

The denominator $x^{2}$ is positive throughout the domain, so the sign of
$f'$ is the sign of $1-\ln x$: positive for $x<\mathrm{e}$, negative
after. A \textbf{maximum}, and since it is the only critical point and the
function rises to it from $-\infty$ and falls to $0$, it is the
\emph{global} maximum.
\wtrap{With a single critical point, the sign of $f'$ on either side settles
the classification faster than computing $f''$ --- and it settles the global
question too, which $f''$ never does.}

\wstep{Step 6. The inflection.}
\[
  f''(x)=\frac{2\ln x-3}{x^{3}} .
\]
This vanishes at $\ln x=\tfrac32$, i.e.\ $x=\mathrm{e}^{3/2}\approx4.482$,
and changes sign there (the denominator is positive), so it is a genuine
inflection at
\[
  \Bigl(\mathrm{e}^{3/2},\ \tfrac{3}{2}\mathrm{e}^{-3/2}\Bigr)\approx(4.48,\ 0.335).
\]
Check the ordering: $\mathrm{e}\approx2.72<4.48$, so the curve peaks and
\emph{then} changes curvature --- which is what a graph tending to zero from
above must do.

\wstep{Step 7. The integral.}
Substituting $u=\ln x$, $\mathrm{d}u=\mathrm{d}x/x$:
\[
  \int_{1}^{b}\frac{\ln x}{x}\,\mathrm{d}x
  =\int_{0}^{\ln b}u\,\mathrm{d}u
  =\frac{(\ln b)^{2}}{2}\xrightarrow[b\to\infty]{}\infty .
\]
\[
  \boxed{\int_{1}^{\infty}\frac{\ln x}{x}\,\mathrm{d}x\ \text{diverges.}}
\]
\wtrap{The integrand tends to $0$ --- Step~4 proved it --- and the integral
still diverges. ``The function goes to zero'' is never sufficient for
convergence, exactly as ``the terms go to zero'' was never sufficient for a
series. Here $f$ behaves like $1/x$ up to a logarithm, and $\int_{1}^{\infty}
\mathrm{d}x/x$ diverges.}

\wstep{Step 8. Assemble the picture, and check it against every part.}
Rising from $-\infty$ at the origin, crossing zero at $x=1$, peaking at
$(\mathrm{e},1/\mathrm{e})$, inflecting at $x\approx4.48$, and decaying to
$0$ --- with a tail thick enough to have infinite area. Every part of the
answer must be visible in that description, and if one is not, something is
wrong.
\begin{center}
\begin{tikzpicture}
\begin{axis}[fnaxis, width=8.0cm, height=4.4cm, domain=0.13:12,
             samples=201, xmin=-0.6, xmax=12.6, ymin=-2.2, ymax=0.9,
             xtick={1,2.718,4.482,8}, xticklabels={$1$,$\mathrm{e}$,$\mathrm{e}^{3/2}$,$8$},
             ytick={0.368}, yticklabels={$1/\mathrm{e}$}]
  \addplot {ln(x)/x};
  \addplot[only marks, mark=*, mark size=1.2pt, accent!70, forget plot]
    coordinates {(2.71828,0.36788) (4.48169,0.33470)};
\end{axis}
\end{tikzpicture}
\end{center}
\end{worked}

\section{Problems}

\begin{problems}
  \item Give the maximal domain of
        $f(x)=\dfrac{\sqrt{x^{2}-9}}{\ln(x-2)}$.
  \item Solve $3^{2x}-4\cdot3^{x}+3=0$.
  \item Find the inverse of $f(x)=\dfrac{2x-1}{x+3}$ and state its domain.
  \item Describe the set $\{(x,y): x^{2}+y^{2}\le25,\ y\ge x\}$ and decide
        whether it is convex.
  \item Give the level curves of $f(x,y)=x^{2}-4y$ at $c=0$ and $c=4$, and
        name the family.
  \item Evaluate $\displaystyle\lim_{x\to0}
        \frac{\mathrm{e}^{3x}-1-3x}{x^{2}}$.
  \item Evaluate $\displaystyle\lim_{x\to\infty}\frac{(\ln x)^{3}}{x}$.
  \item For $f(x)=\dfrac{x}{x^{2}-1}$: give the domain, the symmetry, all
        asymptotes, and the intervals of monotonicity.
  \item Solve $x+2y-z=3$, \;$2x-y+z=5$, \;$3x+y+2z=14$.
  \item Give the rank of
        $A=\begin{pmatrix}1&2&3\\2&4&6\\1&0&1\end{pmatrix}$ and say whether
        $A$ is invertible.
  \item Find the eigenvalues of
        $\begin{pmatrix}2&1\\1&2\end{pmatrix}$.
  \item Determine the definiteness of the Hessian of
        $f(x,y)=x^{2}+4xy+y^{2}$ and what it implies about the critical
        point at the origin.
  \item Find and classify all critical points of
        $f(x,y)=x^{3}-3xy+y^{3}$.
  \item Maximise $U=xy$ subject to $3x+4y=24$, and give the multiplier.
  \item Using the previous answer only, estimate the change in maximal
        utility when the income $24$ rises to $25$, then compute the exact
        change.
  \item Maximise $f(x,y)=x+y$ subject to $x^{2}+y^{2}\le2$, $x,y\ge0$.
  \item Maximise $3x+2y$ subject to $x+y\le4$, $x\le3$, $x,y\ge0$.
  \item Evaluate $\displaystyle\int x\,\mathrm{e}^{x}\,\mathrm{d}x$.
  \item Compute the price elasticity of demand for $q(p)=100/p^{2}$.
  \item At a continuously compounded rate of $5\%$, how long does a
        capital take to double?
\end{problems}

\section{Solutions}

\begin{solutions}
  \item \textbf{Maximal domain.} Three conditions: $x^{2}-9\ge0$ gives
        $|x|\ge3$; $\ln(x-2)$ needs $x>2$; and the denominator must not
        vanish, so $x-2\neq1$, i.e. $x\neq3$. Intersecting,
        $x\ge3$ and $x>2$ and $x\neq3$: \;$D=(3,\infty)$.
  \item \textbf{Substitution in an exponential equation.} With
        $u=3^{x}>0$: $u^{2}-4u+3=0$, so $u=1$ or $u=3$, giving
        $x=0$ and $x=1$.
  \item \textbf{Inverse function.} $y(x+3)=2x-1\Rightarrow x(y-2)=-1-3y$,
        so $x=\dfrac{3y+1}{2-y}$ and
        $f^{-1}(x)=\dfrac{3x+1}{2-x}$ with
        $D(f^{-1})=\R\setminus\{2\}$ --- which is exactly the range of
        $f$, as it must be.
  \item \textbf{Region in the plane.} The half of the disc of radius $5$
        lying on or above the line $y=x$. It is the intersection of two
        convex sets (a disc and a half-plane), hence \textbf{convex}.
  \item \textbf{Level curves.} $x^{2}-4y=c\Rightarrow y=\dfrac{x^{2}-c}{4}$:
        for $c=0$, $y=x^{2}/4$; for $c=4$, $y=\dfrac{x^{2}-4}{4}$. All
        level curves are vertical translates of the same parabola.
  \item \textbf{L'H\^opital, twice.} $\dfrac{3\mathrm{e}^{3x}-3}{2x}$ is
        still $0/0$, then
        $\dfrac{9\mathrm{e}^{3x}}{2}\to\dfrac92$.
  \item \textbf{Growth hierarchy.} $(\ln x)^{3}\ll x$, so the limit is
        $0$. (Or three applications of L'H\^opital, each reducing the power
        of the logarithm by one.)
  \item \textbf{Curve discussion.} Domain $\R\setminus\{-1,1\}$; the
        function is \emph{odd}; single root $x=0$; poles at $x=\pm1$;
        horizontal asymptote $y=0$. And
        $f'=\dfrac{(x^{2}-1)-x\cdot2x}{(x^{2}-1)^{2}}
        =\dfrac{-(x^{2}+1)}{(x^{2}-1)^{2}}<0$ everywhere, so $f$ is
        strictly decreasing on each of the three branches, with \emph{no}
        extrema.
  \item \textbf{Elimination.} Adding the first two equations gives
        $3x+y=8$; substituting into the third, $8+2z=14$ so $z=3$. Then
        $x+2y=6$ and $y=8-3x$ give $x=2$, $y=2$. Solution $(2,2,3)$.
        (Check the second: $4-2+3=5$ \checkmark.)
  \item \textbf{Rank by row reduction.} $R_{2}=2R_{1}$, so that row
        vanishes; $R_{3}-R_{1}=(0,-2,-2)$ survives.
        $\mathrm{rk}(A)=2<3$, so $A$ is \textbf{singular} --- equivalently
        $\det A=0$.
  \item \textbf{Characteristic polynomial.}
        $(2-\lambda)^{2}-1=0\Rightarrow\lambda-2=\pm1$, so
        $\lambda_{1}=3$ and $\lambda_{2}=1$. Both positive, so the matrix
        is positive definite.
  \item \textbf{Definiteness.}
        $\Hess f=\begin{pmatrix}2&4\\4&2\end{pmatrix}$ with
        $\det=4-16=-12<0$: \textbf{indefinite}. So the origin is a
        \textbf{saddle point}, not an extremum.
  \item \textbf{Critical points in two variables.}
        $f_{x}=3x^{2}-3y=0$ gives $y=x^{2}$; $f_{y}=-3x+3y^{2}=0$ gives
        $x=y^{2}$. Hence $x=x^{4}$, so $x(x^{3}-1)=0$ and $x=0$ or $x=1$:
        the points $(0,0)$ and $(1,1)$. With $f_{xx}=6x$, $f_{yy}=6y$,
        $f_{xy}=-3$: at $(0,0)$ the determinant is $0-9=-9<0$, a
        \textbf{saddle}; at $(1,1)$ it is $36-9=27>0$ with $f_{xx}=6>0$, a
        \textbf{local minimum} with $f=-1$.
  \item \textbf{Lagrange.} $\mathcal{L}=xy+\lambda(24-3x-4y)$ gives
        $y=3\lambda$ and $x=4\lambda$, so $4y=3x$. With the constraint,
        $3x+3x=24$, hence $x^{*}=4$, $y^{*}=3$, $U^{*}=12$ and
        $\lambda^{*}=1$.
  \item \textbf{Envelope theorem.} $\mathrm{d}U^{*}/\mathrm{d}M
        =\lambda^{*}=1$, so the estimate is $+1$. Exactly:
        $U^{*}(M)=M^{2}/48$, so
        $U^{*}(25)-U^{*}(24)=\dfrac{625-576}{48}=\dfrac{49}{48}
        \approx1.021$. The multiplier is the first-order term.
  \item \textbf{Kuhn--Tucker (or symmetry).} The objective increases in
        both variables, so the circle constraint binds; by symmetry
        $x=y$, and $2x^{2}=2$ gives $x=y=1$. The non-negativity
        constraints are slack, and $f^{*}=2$. (Formally
        $\grad f=(1,1)^{T}=\lambda(2x,2y)^{T}$ forces $x=y$ with
        $\lambda=\tfrac12>0$.)
  \item \textbf{Linear programming, graphically.} The feasible vertices
        are $(0,0)$, $(3,0)$, $(3,1)$ and $(0,4)$, where the objective
        takes the values $0$, $9$, $11$ and $8$. The maximum is
        $\mathbf{11}$ at $(3,1)$, where both $x\le3$ and $x+y\le4$ bind.
  \item \textbf{Integration by parts.} With $u=x$,
        $v'=\mathrm{e}^{x}$:
        $\int x\mathrm{e}^{x}\mathrm{d}x
        =x\mathrm{e}^{x}-\int\mathrm{e}^{x}\mathrm{d}x
        =(x-1)\mathrm{e}^{x}+C$.
  \item \textbf{Elasticity.}
        $\varepsilon=\dfrac{q'(p)\,p}{q(p)}
        =\dfrac{-200p^{-3}\cdot p}{100p^{-2}}=-2$: constant, independent of
        the price --- the signature of a power demand function.
  \item \textbf{Continuous compounding.}
        $\mathrm{e}^{0.05t}=2\Rightarrow t=\dfrac{\ln2}{0.05}
        \approx13.86$ years. (The ``rule of $70$'': $70/5=14$.)
\end{solutions}

\chapter{Mixed Set B}

Twenty problems, no order, no announced methods. Target time: 75 minutes.

\section{Worked Example}

Again the example is not one of the twenty. This one is a constrained
optimisation, one of the most frequently examined topics, worked
with every check that the marker expects to see.

\begin{worked}{Maximise $f(x,y)=\ln x+2\ln y$ subject to $x+y=12$, with $x,y>0$: solve, certify, and interpret the multiplier}
\wstep{Step 1. Recognise the type before writing anything.}
Two variables, one equality constraint, a differentiable objective: this is
Lagrange. The constraint is linear and could be substituted instead --- keep
that in reserve as the check.

\wstep{Step 2. Set up and differentiate.}
\[
  \mathcal L=\ln x+2\ln y+\lambda(12-x-y),
\]
\[
  \mathcal L_{x}=\frac1x-\lambda=0,
  \qquad
  \mathcal L_{y}=\frac2y-\lambda=0,
  \qquad
  x+y=12 .
\]

\wstep{Step 3. Eliminate $\lambda$.}
\[
  \frac1x=\frac2y
  \quad\Longrightarrow\quad
  y=2x .
\]
The optimal bundle puts twice as much into $y$ as into $x$ --- the ratio of
the coefficients $2$ and $1$ in the objective, exactly as Cobb--Douglas
shares behave, and $\ln$-linear objectives are Cobb--Douglas in disguise.

\wstep{Step 4. Impose the constraint.}
\[
  x+2x=12
  \quad\Longrightarrow\quad
  \boxed{x^{*}=4,\quad y^{*}=8,\quad \lambda^{*}=\tfrac14 }
\]
\[
  f^{*}=\ln4+2\ln8=\ln4+\ln64=\ln256\approx5.545 .
\]

\wstep{Step 5. Certify it with the bordered Hessian.}
\[
  \mathcal L_{xx}=-\frac{1}{x^{2}}=-\tfrac{1}{16},
  \qquad
  \mathcal L_{yy}=-\frac{2}{y^{2}}=-\tfrac{1}{32},
  \qquad
  \mathcal L_{xy}=0,
  \qquad
  g_{x}=g_{y}=1 .
\]
\[
  \overline H=\begin{pmatrix}0&1&1\\1&-\tfrac{1}{16}&0\\1&0&-\tfrac{1}{32}\end{pmatrix},
  \qquad
  \det\overline H=\tfrac{1}{32}+\tfrac{1}{16}=\tfrac{3}{32}>0 ,
\]
which for $n=2$, $m=1$ means a constrained \textbf{maximum}. \checkmark

\wstep{Step 6. Check again by substitution, which is available here.}
$y=12-x$ gives $h(x)=\ln x+2\ln(12-x)$, so
\[
  h'(x)=\frac1x-\frac{2}{12-x}=0
  \;\Longrightarrow\;
  12-x=2x
  \;\Longrightarrow\;
  x=4 ,
\]
\[
  h''(x)=-\frac{1}{x^{2}}-\frac{2}{(12-x)^{2}}<0\ \text{everywhere:\ concave, so the maximum is global.}
\]
\wtrap{Two independent routes to the same answer is the best use of spare
examination time --- far better than re-reading the first one. And here the
substitution route delivers something the bordered Hessian does not: a
\emph{global} conclusion, because $h''<0$ on the whole feasible interval,
not merely at the candidate point.}

\wstep{Step 7. Verify the multiplier's meaning.}
With a general budget $M$ in place of $12$: $x^{*}=M/3$, $y^{*}=2M/3$, so
\[
  V(M)=\ln\frac M3+2\ln\frac{2M}{3}=3\ln M+\text{constant},
  \qquad
  \dd{V}{M}=\frac3M .
\]
At $M=12$ this is $\tfrac14=\lambda^{*}$. \checkmark\ One more unit of budget
raises the maximised objective by about $0.25$.

\wstep{Step 8. The checks a marker looks for.}
\begin{center}\footnotesize
\begin{tabular}{ll}
  \hline
  check & here\\
  \hline
  constraint satisfied & $4+8=12$ \checkmark\\
  domain respected & $x,y>0$ \checkmark\\
  second-order condition & $\det\overline H>0$ \checkmark\\
  multiplier sign & $\lambda>0$: relaxing helps \checkmark\\
  interpretation stated & $\lambda=\partial V/\partial M$ \checkmark\\
  \hline
\end{tabular}
\end{center}
Every one of these is a line, and together they are usually worth more marks
than the algebra that produced the numbers.
\end{worked}

\section{Problems}

\begin{problems}
  \item Simplify $\dfrac{x^{2}-4}{x^{2}-x-6}$ and state where the
        simplification is valid.
  \item Solve $\log_{2}x+\log_{2}(x-2)=3$.
  \item Give the range of $f(x)=\dfrac{x^{2}+1}{x}$ on $x>0$.
  \item Write a system of inequalities for the part of the elliptical disc
        $\dfrac{x^{2}}{4}+\dfrac{y^{2}}{9}\le1$ lying in the second
        quadrant.
  \item Two functions have concentric circles as level curves. In one the
        circles are equally spaced, in the other they crowd outwards. Which
        is $\sqrt{x^{2}+y^{2}}$ and which is $x^{2}+y^{2}$?
  \item Evaluate $\displaystyle\lim_{x\to0^{+}}x^{\sin x}$.
  \item Evaluate $\displaystyle\lim_{x\to2}\frac{x^{3}-8}{x^{2}-4}$.
  \item For $f(x)=\dfrac{\mathrm{e}^{x}}{x}$: give the domain, the
        asymptotes, the monotonicity and any extremum.
  \item Solve $2x+y=7$, \;$4x+2y=14$, and describe the solution set.
  \item Compute
        $\det\begin{pmatrix}1&2&0\\0&3&1\\2&0&4\end{pmatrix}$.
  \item Solve $\begin{pmatrix}2&1\\1&3\end{pmatrix}
        \mathbf{x}=\binom{5}{10}$.
  \item Give $\grad f(1,2)$ for $f(x,y)=x^{2}y^{3}$.
  \item Give the directional derivative of the same $f$ at $(1,2)$ in the
        direction of $(3,4)^{T}$.
  \item Use the total differential to estimate $f(1.1,1.9)$ for the same
        $f$, and compare with the exact value.
  \item Find $\dd{y}{x}$ at $(1,2)$ on the curve $x^{2}+xy+y^{2}=7$.
  \item Minimise $2x+3y$ subject to $xy=54$, with $x,y>0$.
  \item Maximise $\ln x+y$ subject to $x+y\le4$, $x>0$, $y\ge0$.
  \item Maximise $4x+3y$ subject to $x+y\le10$, $2x+y\le16$, $x,y\ge0$.
  \item What is $1000$ received in five years worth today at $6\%$ annual
        interest?
  \item Sum the series $3+\tfrac32+\tfrac34+\tfrac38+\cdots$.
\end{problems}

\section{Solutions}

\begin{solutions}
  \item \textbf{Factor and cancel.}
        $\dfrac{(x-2)(x+2)}{(x-3)(x+2)}=\dfrac{x-2}{x-3}$, valid for
        $x\neq-2$ and $x\neq3$. The point $x=-2$ remains excluded from the
        domain even though the factor cancelled --- a hole, not a pole.
  \item \textbf{Logarithm law, then check the domain.}
        $\log_{2}\bigl(x(x-2)\bigr)=3\Rightarrow x^{2}-2x=8$, so
        $x^{2}-2x-8=0$ and $x=4$ or $x=-2$. The domain requires $x>2$, so
        $x=-2$ is rejected: \;$x=4$.
  \item \textbf{Range via AM--GM (or calculus).}
        $f(x)=x+\dfrac1x$, and for $x>0$, $x+1/x\ge2$ with equality at
        $x=1$. As $x\to0^{+}$ and $x\to\infty$ the function grows without
        bound, so $W=[2,\infty)$.
  \item \textbf{Region.} $\dfrac{x^{2}}{4}+\dfrac{y^{2}}{9}\le1$,
        \;$x\le0$, \;$y\ge0$.
  \item \textbf{Contour spacing.} Equally spaced circles belong to the
        \textbf{cone} $\sqrt{x^{2}+y^{2}}$, whose level $c$ has radius $c$.
        Circles crowding outwards belong to the \textbf{paraboloid}
        $x^{2}+y^{2}$, whose level $c$ has radius $\sqrt{c}$: the surface
        steepens, so equal steps in height need ever smaller steps in
        radius.
  \item \textbf{Exponential form, take logs.}
        $\ln y=\sin x\ln x$, and since $\sin x\sim x$ near $0$ this behaves
        like $x\ln x\to0$. Hence the limit is $\mathrm{e}^{0}=1$.
  \item \textbf{L'H\^opital (or factor).}
        $\dfrac{3x^{2}}{2x}\to\dfrac{12}{4}=3$. Factoring gives the same:
        $\dfrac{(x-2)(x^{2}+2x+4)}{(x-2)(x+2)}\to\dfrac{12}{4}$.
  \item \textbf{Curve discussion.} Domain $\R\setminus\{0\}$. As
        $x\to0^{+}$, $f\to+\infty$; as $x\to0^{-}$, $f\to-\infty$: a pole
        at $x=0$. As $x\to-\infty$, $f\to0^{-}$, so $y=0$ is a horizontal
        asymptote on the left only; as $x\to+\infty$, $f\to+\infty$. And
        $f'=\dfrac{\mathrm{e}^{x}(x-1)}{x^{2}}$, negative for $x<1$
        (excluding $0$) and positive for $x>1$: decreasing, then
        increasing, with a local (and global, on $x>0$) minimum at
        $(1,\mathrm{e})$.
  \item \textbf{Elimination.} The second equation is exactly twice the
        first, so it carries no information: one pivot, two unknowns.
        \textbf{Infinitely many} solutions, $y=7-2x$ with $x$ free ---
        the line itself.
  \item \textbf{Cofactor expansion along the first row.}
        $1\cdot(3\cdot4-1\cdot0)-2\cdot(0\cdot4-1\cdot2)+0
        =12+4=\mathbf{16}$.
  \item \textbf{Inverse or elimination.} $\det=6-1=5$, so
        $\mathbf{x}=\dfrac15\begin{pmatrix}3&-1\\-1&2\end{pmatrix}
        \binom{5}{10}=\dfrac15\binom{5}{15}=\binom{1}{3}$.
        (Check: $2+3=5$ \checkmark, $1+9=10$ \checkmark.)
  \item \textbf{Partial derivatives.} $f_{x}=2xy^{3}$ and
        $f_{y}=3x^{2}y^{2}$, so at $(1,2)$:
        $\grad f=(2\cdot8,\;3\cdot4)^{T}=(16,12)^{T}$.
  \item \textbf{Directional derivative --- normalise first.}
        $\|(3,4)\|=5$, so $\mathbf{u}=(0.6,0.8)^{T}$ and
        $D_{\mathbf{u}}f=(16,12)\cdot(0.6,0.8)=9.6+9.6=\mathbf{19.2}$.
  \item \textbf{Total differential.} $f(1,2)=1\cdot8=8$ and
        $\mathrm{d}f=16\,\mathrm{d}x+12\,\mathrm{d}y
        =16(0.1)+12(-0.1)=0.4$, so $f(1.1,1.9)\approx8.4$. Exactly
        $1.1^{2}\cdot1.9^{3}=1.21\cdot6.859=8.299$; the linear estimate is
        high by about $0.1$, the second-order error.
  \item \textbf{Implicit differentiation.} $2x+y+xy'+2yy'=0$, so
        $y'=-\dfrac{2x+y}{x+2y}$. At $(1,2)$ --- which does lie on the
        curve, since $1+2+4=7$ --- this is
        $-\dfrac{4}{5}$.
  \item \textbf{Lagrange (cost minimisation).}
        $2=\lambda y$ and $3=\lambda x$ give $\dfrac{x}{y}=\dfrac32$, so
        $x=1.5y$. Then $1.5y^{2}=54$, $y^{2}=36$, $y=6$ and $x=9$. Minimal
        cost $18+18=36$ --- the two expenditures are equal, as always for
        a product constraint with linear costs.
  \item \textbf{Kuhn--Tucker.} With
        $\mathcal{L}=\ln x+y+\lambda(4-x-y)$: $1/x=\lambda$ and
        $1=\lambda$, so $\lambda=1$ and $x=1$. The budget binds, so $y=3$,
        which satisfies $y\ge0$. Optimum $(1,3)$ with value
        $\ln1+3=3$.
  \item \textbf{Linear programming.} The two constraints meet where
        $x+y=10$ and $2x+y=16$, i.e. $(6,4)$. Vertices $(0,0)$, $(8,0)$,
        $(6,4)$, $(0,10)$ give objective values $0$, $32$, $36$, $30$: the
        maximum is $\mathbf{36}$ at $(6,4)$, where both constraints bind.
  \item \textbf{Discounting.}
        $\dfrac{1000}{1.06^{5}}=\dfrac{1000}{1.33823}\approx747.26$.
  \item \textbf{Geometric series.} First term $a=3$, ratio
        $r=\tfrac12$, so the sum is $\dfrac{3}{1-\tfrac12}=6$.
\end{solutions}

\chapter{Mixed Set C}

Twenty problems, no order, no announced methods. Target time: 75 minutes.

\section{Worked Example}

The last example is an applied one, of the kind that hides three techniques
inside a sentence of prose. As before it is not among the twenty.

\begin{worked}{A machine yields a net flow of $D(t)=8000\,\mathrm{e}^{-0.1t}$ euros per year, discounted continuously at $r=5\%$. Find (a) the present value over ten years, (b) over an infinite horizon, (c) the time at which half the infinite value has accrued}
\wstep{Step 1. Translate the sentence into an integral.}
``Continuous flow'' plus ``continuous discounting'' means
\[
  PV(T)=\int_{0}^{T}D(t)\,\mathrm{e}^{-rt}\,\mathrm{d}t .
\]
\wtrap{Two exponentials appear and they are different things: $\mathrm{e}^{-0.1t}$
is the machine \emph{decaying}, $\mathrm{e}^{-0.05t}$ is money being
\emph{discounted}. Both must be present, and using only one is the standard
misreading of the sentence.}

\wstep{Step 2. Combine the exponents --- this is what makes it easy.}
\[
  PV(T)=\int_{0}^{T}8000\,\mathrm{e}^{-0.1t}\mathrm{e}^{-0.05t}\,\mathrm{d}t
  =8000\int_{0}^{T}\mathrm{e}^{-0.15t}\,\mathrm{d}t .
\]
The two rates \emph{add}, because $\mathrm{e}^{a}\mathrm{e}^{b}
=\mathrm{e}^{a+b}$: the effective rate of decline is $15\%$.

\wstep{Step 3. Integrate.}
\[
  8000\Bigl[\frac{\mathrm{e}^{-0.15t}}{-0.15}\Bigr]_{0}^{T}
  =\frac{8000}{0.15}\bigl(1-\mathrm{e}^{-0.15T}\bigr)
  =53\,333.33\,\bigl(1-\mathrm{e}^{-0.15T}\bigr).
\]
\wtrap{The factor $1/(-0.15)$ turns $\mathrm{e}^{-0.15T}-1$ into
$1-\mathrm{e}^{-0.15T}$. Get this backwards and the present value comes out
negative, which signals an error: a positive flow cannot have a
negative value.}

\wstep{Step 4. Part (a): ten years.}
\[
  PV(10)=53\,333.33\,\bigl(1-\mathrm{e}^{-1.5}\bigr)
        =53\,333.33\,(1-0.22313)
        =\boxed{41\,433\ \text{euros}}
\]

\wstep{Step 5. Part (b): the infinite horizon.}
\[
  PV(\infty)=\lim_{T\to\infty}53\,333.33\bigl(1-\mathrm{e}^{-0.15T}\bigr)
  =\boxed{53\,333\ \text{euros}}
\]
The limit exists because $0.15>0$. Had the machine's output grown faster
than the discount rate --- $D(t)=8000\mathrm{e}^{+0.06t}$ against $r=0.05$
--- the integral would diverge, and the asset would have no finite value.
That convergence condition is the same $g<r$ that governs a growing
perpetuity in \chref{foundations}{2}.

\wstep{Step 6. Part (c): solve for the time, which is a logarithm problem.}
\[
  53\,333.33\bigl(1-\mathrm{e}^{-0.15T}\bigr)=\tfrac12\cdot53\,333.33
  \;\Longrightarrow\;
  1-\mathrm{e}^{-0.15T}=\tfrac12
  \;\Longrightarrow\;
  \mathrm{e}^{-0.15T}=\tfrac12 .
\]
\[
  -0.15T=\ln\tfrac12=-\ln2
  \quad\Longrightarrow\quad
  T=\frac{\ln2}{0.15}=\boxed{4.62\ \text{years}}
\]
The constant $53\,333.33$ cancels before any arithmetic is done --- notice
that before computing it.
\wtrap{This is a half-life, and it depends only on the \emph{combined} rate
$0.15$, not on the size of the flow. Recognising the shape
$\mathrm{e}^{-kT}=\tfrac12\Rightarrow T=\ln2/k$ saves the whole
manipulation.}

\wstep{Step 7. Sanity-check all three answers together.}
$41\,433<53\,333$ \checkmark: ten years cannot be worth more than the infinite horizon. The ten-year
figure is $77.7\%$ of the total, and $4.62$ years already delivers $50\%$
\checkmark --- consistent with a flow that decays. And $53\,333=8000/0.15$
\checkmark, the perpetuity formula with the combined rate.

\wstep{Step 8. What made this an applied problem rather than a calculus exercise.}
Three separate recognitions: that ``continuous flow'' means an integral,
that two exponential rates add, and that ``half the value'' is an equation
to be solved rather than a quantity to be computed. The calculus itself was
one substitution. That ratio --- most of the difficulty in the translation,
little in the technique --- is characteristic of the applied problems in
these sets, and it is why reading the sentence twice is time well spent.
\end{worked}

\section{Problems}

\begin{problems}
  \item Solve $|2x-5|\le3$.
  \item Give the maximal domain of $f(x)=\ln(x^{2}-x-6)$.
  \item Describe $y=-2\sqrt{x+1}+3$ as a chain of transformations, and give
        its domain and range.
  \item Describe the region $\{y\ge x^{2},\ y\le 2x+3\}$ and give the range
        of $x$ it occupies.
  \item Give the level curve of $f(x,y)=xy$ at $c=-4$ and say which
        quadrants it lies in.
  \item Evaluate $\displaystyle\lim_{x\to\infty}
        x\bigl(\sqrt{x^{2}+1}-x\bigr)$.
  \item Evaluate $\displaystyle\lim_{x\to1}\frac{x^{n}-1}{x-1}$ for a
        positive integer $n$.
  \item Carry out a complete discussion of $f(x)=x^{3}-12x$.
  \item For which $t$ does the system $x+y+tz=1$, \;$x+ty+z=1$,
        \;$tx+y+z=1$ fail to have a unique solution?
  \item Give the eigenvalues of
        $\begin{pmatrix}4&2\\2&1\end{pmatrix}$ and its definiteness.
  \item Determine the definiteness of the Hessian of
        $f(x,y)=x^{2}+y^{2}-2xy$ and describe the set of minimisers.
  \item At the point $(3,4)$ on $x^{2}+y^{2}=25$, may $y$ be written
        locally as a function of $x$? Give $\dd{y}{x}$ there.
  \item Show that $f(K,L)=K^{0.3}L^{0.7}$ is homogeneous of degree $1$ and
        verify Euler's theorem $Kf_{K}+Lf_{L}=f$.
  \item Maximise $-(x-2)^{2}-(y-3)^{2}$ subject to $x+y\le10$. Does the
        constraint bind?
  \item For $\max xy$ subject to $px+qy=M$, show $x^{*}=M/(2p)$ and give
        $\partial x^{*}/\partial p$.
  \item Minimise $2x+3y$ subject to $x+y\ge4$, $x\ge1$, $y\ge0$.
  \item Evaluate $\displaystyle\int_{0}^{1}x^{2}\mathrm{e}^{x^{3}}
        \,\mathrm{d}x$.
  \item Evaluate $\displaystyle\int_{1}^{\infty}\frac{1}{x^{2}}
        \,\mathrm{d}x$.
  \item Demand is $p=20-2q$. At the equilibrium quantity $q=5$, compute the
        consumer surplus.
  \item What is an annuity of $500$ per year for ten years worth today at
        $4\%$?
\end{problems}

\section{Solutions}

\begin{solutions}
  \item \textbf{Absolute-value inequality.}
        $-3\le2x-5\le3\Rightarrow2\le2x\le8$, so $1\le x\le4$.
  \item \textbf{Logarithm domain, hence a quadratic inequality.}
        $x^{2}-x-6=(x-3)(x+2)>0$ holds outside the roots:
        $D=(-\infty,-2)\cup(3,\infty)$.
  \item \textbf{Transformations.} From $\sqrt{x}$: shift $1$ left, stretch
        vertically by $2$, reflect in the $x$-axis, shift $3$ up. Domain
        $[-1,\infty)$; since $-2\sqrt{x+1}\in(-\infty,0]$, the range is
        $(-\infty,3]$, with the curve starting at $(-1,3)$ and falling.
  \item \textbf{Region between two curves.} They meet where
        $x^{2}=2x+3$, i.e. $x=-1$ and $x=3$, so the region is the lens
        above the parabola and below the line, occupying
        $-1\le x\le3$. Convex.
  \item \textbf{Level curve.} $xy=-4\Rightarrow y=-4/x$: a hyperbola whose
        branches lie in the \textbf{second and fourth} quadrants, since the
        product is negative.
  \item \textbf{Conjugate, then a limit at infinity.}
        $x\bigl(\sqrt{x^{2}+1}-x\bigr)
        =\dfrac{x}{\sqrt{x^{2}+1}+x}
        =\dfrac{1}{\sqrt{1+1/x^{2}}+1}\to\dfrac12$.
  \item \textbf{L'H\^opital (or the factorisation of $x^{n}-1$).}
        $\dfrac{nx^{n-1}}{1}\to n$.
  \item \textbf{Complete discussion.} Domain $\R$; \emph{odd}, so
        symmetric about the origin; $f=x(x^{2}-12)$ gives roots
        $0,\pm2\sqrt3$; ends $\mp\infty$. $f'=3x^{2}-12=3(x-2)(x+2)$:
        increasing on $(-\infty,-2)$, decreasing on $(-2,2)$, increasing on
        $(2,\infty)$, so a local maximum $(-2,16)$ and a local minimum
        $(2,-16)$. $f''=6x$: concave for $x<0$, convex for $x>0$, with an
        inflection at the origin $(0,0)$.
  \item \textbf{Determinant with a parameter.} Expanding,
        \[
          \det\begin{pmatrix}1&1&t\\1&t&1\\t&1&1\end{pmatrix}
          =(t-1)\bigl[2-t(1+t)\bigr]
          =-(t-1)^{2}(t+2).
        \]
        This vanishes for $t=1$ and $t=-2$. At $t=1$ all three equations
        become $x+y+z=1$: infinitely many solutions. At $t=-2$ the system
        is inconsistent (adding all three gives $0=3$).
  \item \textbf{Eigenvalues.} Trace $5$ and determinant $4-4=0$, so
        $\lambda^{2}-5\lambda=0$ and $\lambda=0$ or $\lambda=5$. Both
        $\ge0$ with one zero: \textbf{positive semidefinite}, not definite.
  \item \textbf{Hessian and the minimiser set.}
        $\Hess f=\begin{pmatrix}2&-2\\-2&2\end{pmatrix}$ with
        $\det=0$ and trace $4>0$: positive semidefinite, so $f$ is convex
        but not strictly so. Indeed $f=(x-y)^{2}\ge0$, and the minimum
        value $0$ is attained on the whole \textbf{line} $y=x$ --- the
        direction in which the Hessian is singular.
  \item \textbf{Implicit function theorem.} With
        $F=x^{2}+y^{2}-25$: $F_{y}=2y=8\neq0$ at $(3,4)$, so yes, a local
        $y(x)$ exists, and
        $\dd{y}{x}=-\dfrac{F_{x}}{F_{y}}=-\dfrac{6}{8}=-\dfrac34$.
  \item \textbf{Homogeneity and Euler.}
        $f(tK,tL)=(tK)^{0.3}(tL)^{0.7}=t^{0.3+0.7}K^{0.3}L^{0.7}=tf(K,L)$,
        so the degree is $1$. And
        $f_{K}=0.3K^{-0.7}L^{0.7}$, $f_{L}=0.7K^{0.3}L^{-0.3}$, so
        \[
          Kf_{K}+Lf_{L}=0.3K^{0.3}L^{0.7}+0.7K^{0.3}L^{0.7}
          =K^{0.3}L^{0.7}=f .
        \]
        Economically: with constant returns to scale, paying each factor
        its marginal product exactly exhausts the product.
  \item \textbf{Check the unconstrained optimum first.} The objective is
        maximal at $(2,3)$, where $x+y=5\le10$: the point is already
        feasible, so the constraint does \textbf{not} bind, its multiplier
        is $0$, and the maximum is $f=0$ at $(2,3)$. Setting up a
        Lagrangian with an equality would have produced the wrong answer.
  \item \textbf{Lagrange with general parameters.} $y=\lambda p$ and
        $x=\lambda q$ give $px=qy$, so the budget splits evenly:
        $px=M/2$, hence $x^{*}=\dfrac{M}{2p}$ and $y^{*}=\dfrac{M}{2q}$.
        Then $\dfrac{\partial x^{*}}{\partial p}=-\dfrac{M}{2p^{2}}<0$:
        the own-price effect is negative, and the cross-price effect
        $\partial x^{*}/\partial q=0$ --- Cobb--Douglas demands do not
        respond to the other price.
  \item \textbf{Linear programming, minimisation.} The vertices of the
        feasible set are $(4,0)$ (where $x+y=4$ meets $y=0$) and $(1,3)$
        (where $x=1$ meets $x+y=4$). The objective takes the values
        $8$ and $11$, so the minimum is $\mathbf{8}$ at $(4,0)$. The region
        is unbounded, but the objective increases in both variables, so a
        minimum exists.
  \item \textbf{Substitution.} With $u=x^{3}$, $\mathrm{d}u=3x^{2}
        \mathrm{d}x$:
        $\dfrac13\displaystyle\int_{0}^{1}\mathrm{e}^{u}\mathrm{d}u
        =\dfrac{\mathrm{e}-1}{3}\approx0.573$.
  \item \textbf{Improper integral.}
        $\displaystyle\int_{1}^{b}x^{-2}\mathrm{d}x
        =1-\dfrac1b\to1$ as $b\to\infty$: the integral \textbf{converges}
        to $1$.
  \item \textbf{Consumer surplus.} At $q=5$ the price is
        $p=20-10=10$, and
        \[
          CS=\int_{0}^{5}(20-2q)\,\mathrm{d}q-10\cdot5
            =\bigl[20q-q^{2}\bigr]_{0}^{5}-50
            =75-50=25 .
        \]
  \item \textbf{Annuity present value.}
        $500\cdot\dfrac{1-1.04^{-10}}{0.04}
        =500\cdot8.1109\approx4055.4$.
\end{solutions}

\chapter{Mixed Set D: Full Investigations}

The first three sets drill one skill per problem. This set does the
opposite. Each problem hands you a single function of two variables and
runs it through the whole chain of questions that a function of two
variables invites: where it lives, what values it takes, where it can be
differentiated, what its level curves look like, where its gradient
points, what its Hessian says, whether it has critical points, which of
its contour sets are convex, and finally a checklist of properties, each
to be decided \emph{with a reason}.

None of the steps is new. What is new is that they feed each other: the
domain decides where the derivative may be taken, the Hessian decides
concavity, concavity decides quasiconcavity, the domain and the range
decide whether a maximum exists. A slip early in the chain propagates,
and a good investigation catches it, because the later answers stop
fitting together. Target time: twenty to twenty-five minutes per
function. The solutions tag every step with the volume and chapter whose
rule it uses, so that a weak link can be traced back and repaired.

\section{The Investigation Chain}

\begin{rulebox}{The chain, in order}
\begin{enumerate}[label=\arabic*.,leftmargin=1.6em,itemsep=0.3ex]
  \item \textbf{Domain.} Collect the obstructions (roots need a
        nonnegative radicand, logarithms a positive argument,
        denominators must not vanish), sketch the set, and say whether it
        is closed, bounded and convex. (Vol.~2, Ch.~1 and Ch.~4;
        Vol.~5, Ch.~10; Vol.~7, Ch.~2.)
  \item \textbf{Range.} Find the smallest and largest values, or show that
        there are none, and use continuity on a connected domain to fill
        in everything between. (Vol.~2, Ch.~1; Vol.~5, Ch.~9.)
  \item \textbf{Derivative.} Compute $f_{x}$, $f_{y}$ and say where they
        exist and are continuous: on that open set $f$ is differentiable
        (continuous partials are sufficient, not necessary). Where a
        partial fails to exist, $f$ is not differentiable. Roots are the
        usual offenders on the boundary. (Vol.~5, Ch.~10.)
  \item \textbf{Level curves.} Solve $f(x,y)=c$ for one variable and
        sketch several values of $c$. (Vol.~5, Ch.~9.)
  \item \textbf{Gradient at a point.} Evaluate $\grad f$, draw it, check
        that it is orthogonal to the level curve through the point, and
        read it: it points in the direction of steepest ascent, and its
        norm $\lVert\grad f\rVert$ is the slope in that direction.
        (Vol.~5, Ch.~10 and Ch.~13.)
  \item \textbf{Hessian.} Second partials, Young's theorem, and the
        definiteness by Sylvester's criterion. (Vol.~5, Ch.~11.)
  \item \textbf{Critical points.} Solve $\grad f=\mathbf{0}$ \emph{inside
        the domain}, or explain why there is no solution; classify with the
        Hessian. (Vol.~7, Ch.~5.)
  \item \textbf{Contour sets.} Mark an upper set $\{f\ge c\}$ and a lower
        set $\{f\le c\}$ and decide whether each is convex.
        (Vol.~5, Ch.~9; Vol.~7, Ch.~2.)
  \item \textbf{Checklist.} Concave, convex, quasiconcave, quasiconvex,
        has a maximum, has a minimum, increasing, decreasing, epigraph
        convex: each answered with a proof or a counterexample.
        (Vol.~7, Ch.~1 and Ch.~3.)
\end{enumerate}
\end{rulebox}

\begin{rulebox}{Three definitions the checklist needs}
\begin{itemize}[leftmargin=1.6em,itemsep=0.3ex]
  \item \textbf{Monotone in several variables.} Order points componentwise:
        $(x,y)\le(x',y')$ means $x\le x'$ and $y\le y'$. Then $f$ is
        \emph{increasing} \de{monoton wachsend} if
        $(x,y)\le(x',y')\Rightarrow f(x,y)\le f(x',y')$, and
        \emph{decreasing} if the inequality between the values is reversed.
        On a convex domain with interior points, a continuous $f$ with
        $f_{x}\ge0$ and $f_{y}\ge0$ at every interior point is increasing.
        To refute it, exhibit two ordered points whose values are in the
        wrong order.
  \item \textbf{Epigraph and hypograph.} The \emph{epigraph}
        \de{Epigraph} of $f$ is the set of points on or above its graph,
        $\{(x,y,z): (x,y)\in D,\ z\ge f(x,y)\}$; the \emph{hypograph} is
        the set on or below it. On a convex domain, the epigraph is a convex
        set exactly when $f$ is convex, and the hypograph exactly when $f$
        is concave. So the epigraph question is the convexity question
        asked again, and the two answers must agree.
  \item \textbf{Existence of extrema.} A continuous function on a closed
        and bounded domain attains a maximum and a minimum (Weierstrass,
        Vol.~7, Ch.~1). On an unbounded or open domain, look at the range
        first: if $f$ grows without bound there is no maximum, and if the
        infimum is approached but never reached there is no minimum.
        Extrema on the boundary of the domain are not found by
        $\grad f=\mathbf{0}$, and are often points where $f$ is not even
        differentiable.
\end{itemize}
\end{rulebox}

\begin{rulebox}{Shortcuts that keep the checklist consistent}
\begin{itemize}[leftmargin=1.6em,itemsep=0.3ex]
  \item Concave $\Rightarrow$ quasiconcave; convex $\Rightarrow$
        quasiconvex. A function that is both concave and convex is affine.
  \item A strictly increasing transform $F$ ($\ln$ of a positive $f$,
        $\sqrt{\cdot}$ or a positive power of a nonnegative $f$) only
        relabels the contour sets, $\{F(f)\ge F(c)\}=\{f\ge c\}$, so it
        preserves quasiconcavity and quasiconvexity but not necessarily
        concavity.
  \item If $g$ is concave and $F$ is increasing and concave, $F\circ g$ is
        concave: this is how $\sqrt{\text{concave}}$ and
        $\ln(\text{concave})$ are handled without a Hessian.
  \item One straight line in the domain along which $f$ is strictly convex
        refutes concavity; one along which it is strictly concave refutes
        convexity. Two points in a contour set whose midpoint is not in it
        refute quasiconcavity or quasiconvexity.
  \item Cobb--Douglas $x^{a}y^{b}$ ($a,b>0$) on $\R^{2}_{+}$: always
        quasiconcave; concave iff $a+b\le1$; never convex.
        (Vol.~7, Ch.~3.)
\end{itemize}
\end{rulebox}

\section{Worked Example}

\begin{worked}{Investigate $f(x,y)=\sqrt{8-2x^{2}-y^{2}}$ completely: domain, range, derivative, level curves, gradient at $(1,2)$, Hessian, critical points, contour sets, and the full property checklist}
\wstep{Step 1. Domain (Vol.~2, Ch.~1; Vol.~2, Ch.~4).}
The only obstruction is the root: $8-2x^{2}-y^{2}\ge0$, that is
\[
  D=\{(x,y): 2x^{2}+y^{2}\le8\}
   =\Bigl\{\tfrac{x^{2}}{4}+\tfrac{y^{2}}{8}\le1\Bigr\},
\]
the filled ellipse with semi-axes $2$ along $x$ and $2\sqrt2$ along $y$,
boundary included. It is \textbf{closed} (the inequality is weak),
\textbf{bounded}, and \textbf{convex}: it is a lower contour set of the
convex function $2x^{2}+y^{2}$ (Vol.~7, Ch.~2).
\wtrap{Writing $2x^{2}+y^{2}<8$ loses the boundary. The root of zero is
zero, so the boundary ellipse belongs to the domain, and it will carry the
minimum.}

\wstep{Step 2. Range (Vol.~2, Ch.~1).}
The radicand is largest at the origin, where it equals $8$, and smallest
($0$) on the boundary. Since $\sqrt{\cdot}$ is increasing,
$0\le f\le\sqrt8=2\sqrt2$. $f$ is continuous on the connected set $D$, so
every value in between is taken:
\[
  f(D)=[\,0,\,2\sqrt2\,].
\]

\wstep{Step 3. Derivative, and where it exists (Vol.~5, Ch.~10).}
By the chain rule,
\[
  f_{x}=\frac{-4x}{2\sqrt{8-2x^{2}-y^{2}}}=-\frac{2x}{f},\qquad
  f_{y}=\frac{-2y}{2\sqrt{8-2x^{2}-y^{2}}}=-\frac{y}{f}.
\]
Both are continuous wherever $f>0$, so $f$ is differentiable on the
\textbf{interior} $2x^{2}+y^{2}<8$. On the boundary the formulas divide by
zero, and the difference quotient confirms the failure: at $(2,0)$,
$\bigl(f(2-h,0)-f(2,0)\bigr)/h=\sqrt{8h-2h^{2}}/h\to\infty$ as
$h\to0^{+}$. The surface meets the plane $z=0$ vertically, like the edge
of a dome.

\wstep{Step 4. Level curves (Vol.~5, Ch.~9).}
For $0\le c\le2\sqrt2$, $f=c$ means $2x^{2}+y^{2}=8-c^{2}$: concentric
ellipses with the same shape as the boundary, shrinking to the origin as
$c\to2\sqrt2$. The values $c=0,1,\sqrt2,2,\sqrt7$ give
$2x^{2}+y^{2}=8,7,6,4,1$ (left panel below). The level $c=0$ is the
boundary of the domain; the level $c=2\sqrt2$ is the single point
$(0,0)$.

\wstep{Step 5. Gradient at $(1,2)$ (Vol.~5, Ch.~10).}
Here $f(1,2)=\sqrt{8-2-4}=\sqrt2$, so
\[
  \grad f(1,2)=\begin{pmatrix}-2/\sqrt2\\-2/\sqrt2\end{pmatrix}
  =\begin{pmatrix}-\sqrt2\\-\sqrt2\end{pmatrix},
  \qquad
  \lVert\grad f(1,2)\rVert=\sqrt{2+2}=2 .
\]
\emph{Orthogonality.} The level curve through $(1,2)$ is
$2x^{2}+y^{2}=6$; implicit differentiation gives
$y'=-4x/(2y)=-1$ there, so a tangent direction is $(1,-1)$, and
$(-\sqrt2)(1)+(-\sqrt2)(-1)=0$ \checkmark.
\emph{Direction.} The gradient points into the ellipse, towards higher
levels: the direction of steepest ascent. \emph{Norm.} $2$ is the slope of
the surface in that direction, the largest directional derivative at
$(1,2)$ (Vol.~5, Ch.~13). In every direction along the level curve the
slope is $0$.

\begin{gallery}
\begin{tikzpicture}
  \begin{axis}[contouraxis, xmin=-3.2, xmax=3.2, ymin=-3.2, ymax=3.2,
      xtick={-2,-1,1,2}, ytick={-2,-1,1,2},
      title={domain, levels, $\grad f(1,2)$}]
    \addplot[name path=top, draw=none, domain=-2:2]
      {sqrt(max(8-2*x^2,0))};
    \addplot[name path=bot, draw=none, domain=-2:2]
      {-sqrt(max(8-2*x^2,0))};
    \region{top}{bot}{}
    \addplot[accent, line width=1pt, domain=0:360, samples=121]
      ({2*cos(x)},{sqrt(8)*sin(x)});
    \addplot[domain=0:360, samples=121]({sqrt(3.5)*cos(x)},{sqrt(7)*sin(x)});
    \addplot[domain=0:360, samples=121]({sqrt(3)*cos(x)},{sqrt(6)*sin(x)});
    \addplot[domain=0:360, samples=121]({sqrt(2)*cos(x)},{2*sin(x)});
    \addplot[domain=0:360, samples=121]({sqrt(0.5)*cos(x)},{sin(x)});
    \draw[dashed, gray!70!black, line width=0.5pt]
      (axis cs:0.1,2.9) -- (axis cs:1.9,1.1);
    \draw[-{Stealth[length=6pt,width=4pt]}, accent, line width=1pt]
      (axis cs:1,2) -- (axis cs:-0.414,0.586);
    \addplot[only marks, mark=*, mark size=1.2pt, black, forget plot]
      coordinates {(1,2) (0,0)};
  \end{axis}
\end{tikzpicture}\hfill
\begin{tikzpicture}
  \begin{axis}[contouraxis, xmin=-3.2, xmax=3.2, ymin=-3.2, ymax=3.2,
      xtick={-2,-1,1,2}, ytick={-2,-1,1,2},
      title={upper set $\{f\ge2\}$ (shaded)}]
    \addplot[name path=itop, draw=none, domain=-1.4142:1.4142]
      {sqrt(max(4-2*x^2,0))};
    \addplot[name path=ibot, draw=none, domain=-1.4142:1.4142]
      {-sqrt(max(4-2*x^2,0))};
    \region{itop}{ibot}{}
    \addplot[accent, line width=1pt, domain=0:360, samples=121]
      ({2*cos(x)},{sqrt(8)*sin(x)});
    \addplot[domain=0:360, samples=121]({sqrt(2)*cos(x)},{2*sin(x)});
    \draw[gray!70!black, line width=0.5pt] (axis cs:0,2) -- (axis cs:0,-2);
    \addplot[only marks, mark=*, mark size=1.2pt, black, forget plot]
      coordinates {(0,2) (0,-2)};
    \addplot[only marks, mark=o, mark size=1.6pt, black, forget plot]
      coordinates {(0,0)};
  \end{axis}
\end{tikzpicture}
\gcap{Left: the domain (shaded), its boundary $f=0$ (bold), the levels
$f=1,\sqrt2,2,\sqrt7$, and the gradient at $(1,2)$, orthogonal to the
dashed tangent of its level curve. Right: the upper set $\{f\ge2\}$ is a
filled ellipse; the lower set $\{f\le2\}$ is the ring between the two
curves, which contains $(0,\pm2)$ but not their midpoint.}
\end{gallery}

\wstep{Step 6. Hessian (Vol.~5, Ch.~11).}
Differentiate $f_{x}=-2xf^{-1}$ and $f_{y}=-yf^{-1}$ once more, using
$f_{x}$, $f_{y}$ from Step 3:
\[
  f_{xx}=-\frac{2}{f}+\frac{2x f_{x}}{f^{2}}
        =-\frac{2f^{2}+4x^{2}}{f^{3}}
        =-\frac{2(8-y^{2})}{f^{3}},\qquad
  f_{yy}=-\frac{1}{f}+\frac{y f_{y}}{f^{2}}
        =-\frac{8-2x^{2}}{f^{3}},
\]
\[
  f_{xy}=\frac{2x f_{y}}{f^{2}}=-\frac{2xy}{f^{3}}
  \quad(\text{equal to } f_{yx}\text{ by Young}).
\]
The determinant simplifies remarkably:
\[
  \det\Hess f=\frac{2(8-y^{2})(8-2x^{2})-4x^{2}y^{2}}{f^{6}}
  =\frac{16(8-2x^{2}-y^{2})}{f^{6}}=\frac{16}{f^{4}}>0 .
\]
In the interior $y^{2}<8$, so $f_{xx}<0$: the Hessian is
\textbf{negative definite} at every interior point. At $(1,2)$, with
$f^{3}=2\sqrt2$:
\[
  \Hess f(1,2)=\begin{pmatrix}-2\sqrt2&-\sqrt2\\[2pt]-\sqrt2&-\tfrac32\sqrt2
  \end{pmatrix},\qquad \det=6-2=4=\frac{16}{(\sqrt2)^{4}}\ \checkmark
\]
\wtrap{Do not expand $\det\Hess f$ with the roots still in it. Writing
every entry over the common denominator $f^{3}$ and replacing $f^{2}$ by
$8-2x^{2}-y^{2}$ is what makes the sign visible.}

\wstep{Step 7. Critical points (Vol.~7, Ch.~5).}
$f_{x}=0\iff x=0$ and $f_{y}=0\iff y=0$, so the only candidate is
$(0,0)$, which lies in the interior. There
$\Hess f(0,0)=\begin{pmatrix}-1/\sqrt2&0\\0&-1/(2\sqrt2)\end{pmatrix}$ is
negative definite: a strict local maximum, with value $2\sqrt2$. The
boundary points are not critical points; $f$ is not differentiable there.

\wstep{Step 8. Contour sets (Vol.~5, Ch.~9; Vol.~7, Ch.~2).}
Take $c=2$. The upper set $\{f\ge2\}=\{2x^{2}+y^{2}\le4\}$ is a filled
ellipse: \textbf{convex}. The lower set
$\{f\le2\}=\{4\le2x^{2}+y^{2}\le8\}$ is an elliptical ring: it contains
$(0,2)$ and $(0,-2)$ but not their midpoint $(0,0)$, so it is
\textbf{not convex} (right panel).

\wstep{Step 9. The checklist, each with its reason.}
\begin{itemize}[leftmargin=1.6em,itemsep=0.3ex]
  \item \emph{Concave: yes.} The radicand $8-2x^{2}-y^{2}$ is concave and
        $\sqrt{\cdot}$ is increasing and concave, so the composition is
        concave on $D$ (Vol.~7, Ch.~3). Step 6 confirms it, even strictly,
        on the interior.
  \item \emph{Convex: no.} The Hessian is negative definite; concretely,
        $f(0,0)=2\sqrt2$ lies above the chord value
        $\tfrac12\bigl(f(-2,0)+f(2,0)\bigr)=0$.
  \item \emph{Quasiconcave: yes}, because concave; every upper set is a
        filled ellipse (Step 8).
  \item \emph{Quasiconvex: no.} The lower set $\{f\le2\}$ is not convex
        (Step 8).
  \item \emph{Maximum: yes}, $2\sqrt2$ at $(0,0)$, unique; global because
        $f$ is concave and $(0,0)$ is a critical point (Vol.~7, Ch.~5).
  \item \emph{Minimum: yes}, $0$, attained at \emph{every} point of the
        boundary ellipse. Weierstrass guarantees it (continuous on a closed
        bounded set); the first-order condition could never find it, since
        $f$ is not differentiable there.
  \item \emph{Increasing: no.} $(0,0)\le(1,0)$ but
        $f(1,0)=\sqrt6<2\sqrt2=f(0,0)$.
  \item \emph{Decreasing: no.} $(-1,0)\le(0,0)$ but
        $f(-1,0)=\sqrt6<f(0,0)$.
  \item \emph{Epigraph convex: no}, since $f$ is not convex. Its
        hypograph, the solid dome, is convex, matching concavity.
\end{itemize}

\wstep{Step 10. Do the answers fit together?}
Concave with an interior critical point gives a global maximum
\checkmark. Not quasiconvex, yet concave: a dome has convex upper sets and
ring-shaped lower sets \checkmark. The range $[0,2\sqrt2]$ matches the
maximum and minimum found \checkmark. Neither monotone, as one expects of
a function that peaks in the middle of its domain \checkmark.
\end{worked}

\section{Problems}

Each problem is one investigation. Work every part, and give a reason for
every item of the checklist.

\begin{problems}
  \item Investigate $f(x,y)=\sqrt{6-x-y^{2}}$.
        \begin{enumerate}[label=(\alph*),leftmargin=1.8em,itemsep=0.2ex]
          \item Give and sketch the domain. Is it closed, bounded, convex?
                What values does $f$ take?
          \item Compute $f_{x}$ and $f_{y}$. Is $f$ differentiable on its
                whole domain?
          \item Describe and sketch the level curves.
          \item Compute $\grad f(1,1)$, draw it, show that it is orthogonal
                to the level curve through $(1,1)$, and state what its
                direction and its norm mean.
          \item Compute the Hessian and its definiteness.
          \item Does $f$ have critical points?
          \item Mark an upper and a lower contour set at level $2$. Which is
                convex?
          \item Decide, with reasons: concave, convex, quasiconcave,
                quasiconvex, has a maximum, has a minimum, increasing,
                decreasing, epigraph convex.
        \end{enumerate}
  \item Investigate the utility function
        $u:\R^{2}_{++}\to\R$, $u(x,y)=\ln\bigl(x^{1/3}y^{1/2}\bigr)$, where
        $\R^{2}_{++}=\{(x,y):x>0,\ y>0\}$.
        \begin{enumerate}[label=(\alph*),leftmargin=1.8em,itemsep=0.2ex]
          \item Rewrite $u$ as a sum. Describe the domain (closed? bounded?
                convex?) and the range.
          \item Describe and sketch the level curves. How do they compare
                with those of $x^{1/3}y^{1/2}$?
          \item Compute $\grad u(1,1)$, draw it, check orthogonality to the
                level curve, and interpret direction and norm.
          \item Compute the Hessian. Show that both marginal utilities are
                diminishing, and give the definiteness.
          \item Does $u$ have critical points?
          \item Mark an upper and a lower contour set at level $0$. Which is
                convex?
          \item Decide, with reasons: concave, convex, quasiconcave,
                quasiconvex, has a maximum, has a minimum, increasing,
                decreasing, epigraph convex.
        \end{enumerate}
  \item Investigate the production function $f:\R^{2}_{+}\to\R$,
        $f(x,y)=x^{3/4}y^{1/2}$.
        \begin{enumerate}[label=(\alph*),leftmargin=1.8em,itemsep=0.2ex]
          \item Name this type of function and its returns to scale. Give
                the domain (closed? convex?) and the range.
          \item Compute the derivative and say where it exists. Show that
                both marginal products are diminishing.
          \item Compute the Hessian and its definiteness. Does diminishing
                marginal productivity make $f$ concave?
          \item Sketch the level curves, and draw $\grad f(1,1)$; check
                orthogonality and interpret direction and norm.
          \item Does $f$ have critical points? Where does its minimum come
                from?
          \item Mark an upper and a lower contour set at level $1$. Which is
                convex?
          \item Decide, with reasons: concave, convex, quasiconcave,
                quasiconvex, has a maximum, has a minimum, increasing,
                decreasing, epigraph convex.
        \end{enumerate}
  \item Investigate $f:\R^{2}_{+}\to\R$, $f(x,y)=x^{1/3}y^{2/3}$.
        \begin{enumerate}[label=(\alph*),leftmargin=1.8em,itemsep=0.2ex]
          \item Name the function and its returns to scale; give domain and
                range.
          \item Compute the derivative and the Hessian, and the
                definiteness. Which direction does the Hessian annihilate?
          \item Sketch the level curves and draw $\grad f(1,8)$; check
                orthogonality, interpret direction and norm, and verify
                Euler's theorem at $(1,8)$.
          \item Does $f$ have critical points?
          \item Mark an upper and a lower contour set at level $4$. Which is
                convex?
          \item Decide, with reasons: concave, strictly concave, convex,
                quasiconcave, quasiconvex, has a maximum, has a minimum,
                increasing, decreasing, epigraph convex.
        \end{enumerate}
  \item Investigate $f(x,y)=\ln\bigl(4+x^{2}+y^{2}\bigr)$.
        \begin{enumerate}[label=(\alph*),leftmargin=1.8em,itemsep=0.2ex]
          \item Give the domain and the range.
          \item Compute the gradient; where is $f$ differentiable?
          \item Describe the level curves, compute $\grad f(1,1)$, and check
                that it is orthogonal to the level curve through $(1,1)$.
          \item Compute the Hessian, its determinant in general, and its
                definiteness at $(0,0)$, $(1,1)$ and $(2,2)$.
          \item Find and classify the critical points.
          \item Mark the upper and lower contour sets at level $\ln8$.
                Which is convex?
          \item Decide, with reasons: concave, convex, quasiconcave,
                quasiconvex, has a maximum, has a minimum, increasing,
                decreasing, epigraph convex. On which disk centred at the
                origin is $f$ convex?
        \end{enumerate}
\end{problems}

\section{Solutions}

\begin{solutions}
  \item \textbf{Root of a concave expression on an unbounded closed domain.}
        \begin{enumerate}[label=(\alph*),leftmargin=1.8em,itemsep=0.4ex]
          \item \emph{Domain (Vol.~2, Ch.~1 and Ch.~4).} Radicand
                $\ge0$: $D=\{x\le6-y^{2}\}$, the region on and to the left of
                the parabola $x=6-y^{2}$ (vertex $(6,0)$, opening to the
                left). Closed (weak inequality), unbounded (it contains the
                whole ray $y=0$, $x\le6$), and convex: it is
                the upper set $\{g\ge0\}$ of the concave function
                $g=6-x-y^{2}$ (Vol.~7, Ch.~2). \emph{Range
                (Vol.~2, Ch.~1):} $f=0$ on the parabola and
                $f(x,0)=\sqrt{6-x}\to\infty$ as $x\to-\infty$; by
                continuity $f(D)=[0,\infty)$.
          \item \emph{Derivative (Vol.~5, Ch.~10).}
                $f_{x}=\dfrac{-1}{2\sqrt{6-x-y^{2}}}=-\dfrac{1}{2f}$,
                \;$f_{y}=\dfrac{-2y}{2f}=-\dfrac{y}{f}$, continuous where
                $f>0$. So $f$ is differentiable on the interior
                $x<6-y^{2}$ only. On the parabola the partials blow up:
                at $(6,0)$, $\bigl(f(6-h,0)-0\bigr)/h=\sqrt h/h\to\infty$.
          \item \emph{Level curves (Vol.~5, Ch.~9).} $f=c\ge0$ gives
                $x=6-c^{2}-y^{2}$: copies of the boundary parabola, shifted
                left by $c^{2}$. For $c=0,1,2,3$ the vertices sit at
                $x=6,5,2,-3$; higher levels lie further left.
          \item \emph{Gradient (Vol.~5, Ch.~10).} $f(1,1)=\sqrt4=2$, so
                \[
                  \grad f(1,1)=\begin{pmatrix}-1/4\\-1/2\end{pmatrix},\qquad
                  \lVert\grad f(1,1)\rVert=\sqrt{\tfrac1{16}+\tfrac14}
                  =\frac{\sqrt5}{4}\approx0.56 .
                \]
                The level curve through $(1,1)$ is $x=2-y^{2}$; there
                $\mathrm{d}x/\mathrm{d}y=-2y=-2$, so
                $\mathrm{d}y/\mathrm{d}x=-\tfrac12$ and $(2,-1)$ is a tangent
                direction. $(-\tfrac14)(2)+(-\tfrac12)(-1)=0$ \checkmark.
                The gradient points left and down, towards the inside of the
                parabolas where the levels rise: steepest ascent, with slope
                $\sqrt5/4$ in that direction (Vol.~5, Ch.~13).
          \item \emph{Hessian (Vol.~5, Ch.~11).} From $f_{x}=-\tfrac12f^{-1}$
                and $f_{y}=-yf^{-1}$:
                \[
                  f_{xx}=\tfrac12f^{-2}f_{x}=-\frac{1}{4f^{3}},\qquad
                  f_{xy}=\tfrac12f^{-2}f_{y}=-\frac{y}{2f^{3}},
                \]
                \[
                  f_{yy}=-\frac1f+\frac{yf_{y}}{f^{2}}
                        =-\frac{f^{2}+y^{2}}{f^{3}}=-\frac{6-x}{f^{3}} .
                \]
                $\det\Hess f=\dfrac{(6-x)-y^{2}}{4f^{6}}
                =\dfrac{f^{2}}{4f^{6}}=\dfrac{1}{4f^{4}}>0$ and
                $f_{xx}<0$: \textbf{negative definite} on the interior. At
                $(1,1)$: $\Hess f=\begin{pmatrix}-1/32&-1/16\\-1/16&-5/8
                \end{pmatrix}$, $\det=\tfrac{5}{256}-\tfrac{1}{256}
                =\tfrac1{64}=\tfrac{1}{4\cdot2^{4}}$ \checkmark.
          \item \emph{Critical points (Vol.~7, Ch.~5).} None:
                $f_{x}=-1/(2f)$ is never $0$. Moving left always raises $f$.
          \item \emph{Contour sets (Vol.~5, Ch.~9; Vol.~7, Ch.~2).}
                $\{f\ge2\}=\{x\le2-y^{2}\}$, the inside of a parabola:
                \textbf{convex}. $\{f\le2\}=\{2-y^{2}\le x\le6-y^{2}\}$, a
                curved band between two parabolas: it contains $(-2,2)$ and
                $(-2,-2)$ (check: $-2\le-2\le2$) but not their midpoint
                $(-2,0)$ (since $2\le-2$ fails). \textbf{Not convex.}
          \item \emph{Checklist (Vol.~7, Ch.~1 and Ch.~3).}
                \emph{Concave: yes}, $\sqrt{\cdot}$ increasing and concave
                applied to the concave $6-x-y^{2}$ (and (e) confirms it).
                \emph{Convex: no}, negative definite Hessian.
                \emph{Quasiconcave: yes}, because concave.
                \emph{Quasiconvex: no}, by the band in (g).
                \emph{Maximum: no}, $f$ is unbounded above.
                \emph{Minimum: yes}, $0$ on the whole parabola $x=6-y^{2}$,
                boundary points where $f$ is not differentiable.
                \emph{Increasing: no}, $f$ falls in $x$: $(0,0)\le(1,0)$
                with $f(1,0)=\sqrt5<\sqrt6=f(0,0)$.
                \emph{Decreasing: no}, $(0,-1)\le(0,0)$ with
                $f(0,-1)=\sqrt5<\sqrt6=f(0,0)$.
                \emph{Epigraph convex: no}, since $f$ is not convex (the
                hypograph is).
        \end{enumerate}
  \item \textbf{Log-Cobb--Douglas.}
        \begin{enumerate}[label=(\alph*),leftmargin=1.8em,itemsep=0.4ex]
          \item \emph{Logarithm rules (Vol.~1, Ch.~2).}
                $u=\tfrac13\ln x+\tfrac12\ln y$. \emph{Domain
                (Vol.~5, Ch.~10):} the open quadrant $\R^{2}_{++}$, open
                (not closed), unbounded, convex. \emph{Range:}
                $u(t,1)=\tfrac13\ln t$ runs through all of $\R$ as $t$
                runs through $(0,\infty)$, so $u(\R^{2}_{++})=\R$.
          \item \emph{Level curves (Vol.~5, Ch.~9).}
                $\tfrac13\ln x+\tfrac12\ln y=c\iff\ln y=2c-\tfrac23\ln x
                \iff y=\mathrm{e}^{2c}x^{-2/3}$: decreasing, convex curves
                asymptotic to both axes, moving outwards as $c$ grows. They
                are exactly the level curves of $x^{1/3}y^{1/2}$, relabelled
                ($u=\ln v$ with $v$ the Cobb--Douglas): an increasing
                transform changes the labels, not the curves.
          \item \emph{Gradient (Vol.~5, Ch.~10).} $u_{x}=\dfrac{1}{3x}$,
                $u_{y}=\dfrac{1}{2y}$, so $\grad u(1,1)=(\tfrac13,\tfrac12)$
                with norm $\sqrt{\tfrac19+\tfrac14}=\tfrac{\sqrt{13}}{6}
                \approx0.60$. The level curve through $(1,1)$ is $u=0$, i.e.
                $y=x^{-2/3}$, with slope $-\tfrac23$ at $x=1$ and tangent
                direction $(3,-2)$; $\tfrac13\cdot3-\tfrac12\cdot2=0$
                \checkmark. The gradient points up and to the right, the
                direction in which utility rises fastest, at rate
                $\sqrt{13}/6$.
          \item \emph{Hessian (Vol.~5, Ch.~11).}
                $\Hess u=\begin{pmatrix}-\frac{1}{3x^{2}}&0\\0&-\frac{1}{2y^{2}}
                \end{pmatrix}$. $u_{xx}<0$ and $u_{yy}<0$: each marginal
                utility $u_{x}=1/(3x)$, $u_{y}=1/(2y)$ falls as its own good
                increases. Diagonal with negative entries: \textbf{negative
                definite} everywhere.
          \item \emph{Critical points (Vol.~7, Ch.~5).} None: $u_{x}=1/(3x)>0$.
          \item \emph{Contour sets.} $u\le0\iff x^{1/3}y^{1/2}\le1\iff
                x^{2}y^{3}\le1$. The upper set $\{u\ge0\}=\{y\ge x^{-2/3}\}$
                lies above a convex curve: \textbf{convex}. The lower set is
                not: $(8,\tfrac14)$ and $(\tfrac18,4)$ both satisfy
                $x^{2}y^{3}=1$, but at the midpoint
                $(\tfrac{65}{16},\tfrac{17}{8})$ the product is far above
                $1$ (already $x^{2}>16$ and $y^{3}>8$).
          \item \emph{Checklist (Vol.~7, Ch.~3).}
                \emph{Concave: yes}, strictly, by the negative definite
                Hessian (or as a positive sum of the concave $\ln x$,
                $\ln y$).
                \emph{Convex: no}, the Hessian is negative definite.
                \emph{Quasiconcave: yes}, because concave.
                \emph{Quasiconvex: no}, by (f).
                \emph{Maximum: no}, $u(t,1)\to\infty$ as $t\to\infty$.
                \emph{Minimum: no}, $u(t,1)\to-\infty$ as $t\to0^{+}$; the
                open domain has no boundary to stop at.
                \emph{Increasing: yes}, strictly: $u_{x}>0$, $u_{y}>0$ on a
                convex domain.
                \emph{Decreasing: no}, as it is strictly increasing.
                \emph{Epigraph convex: no}, $u$ is not convex.
        \end{enumerate}
  \item \textbf{Cobb--Douglas with exponent sum above $1$.}
        \begin{enumerate}[label=(\alph*),leftmargin=1.8em,itemsep=0.4ex]
          \item \emph{Name (Vol.~5, Ch.~14).} A Cobb--Douglas function,
                homogeneous of degree $\tfrac34+\tfrac12=\tfrac54>1$:
                $f(tx,ty)=t^{5/4}f(x,y)$, \textbf{increasing returns to
                scale}. Domain $\R^{2}_{+}$: closed, unbounded, convex.
                Range $[0,\infty)$: $f=0$ on the axes and $f(t,t)=t^{5/4}$
                is unbounded.
          \item \emph{Derivative (Vol.~5, Ch.~10).}
                \[
                  f_{x}=\tfrac34x^{-1/4}y^{1/2},\qquad
                  f_{y}=\tfrac12x^{3/4}y^{-1/2},
                \]
                continuous on the open quadrant $x,y>0$, so $f$ is
                differentiable there. On the axis $x=0$ (with $y>0$) the
                factor $x^{-1/4}$ is infinite, and likewise $y^{-1/2}$ on the
                axis $y=0$ (with $x>0$): the partials do not exist there.
                The origin is the exception: $f$ vanishes on both axes, so
                $f_{x}(0,0)=f_{y}(0,0)=0$, and
                $0\le f(h,k)\le\lVert(h,k)\rVert^{5/4}$ shows that $f$ is
                differentiable there (from inside the quadrant) with zero
                derivative. Second partials:
                $f_{xx}=-\tfrac{3}{16}x^{-5/4}y^{1/2}<0$,
                $f_{yy}=-\tfrac14x^{3/4}y^{-3/2}<0$: both marginal products
                \textbf{diminish}.
          \item \emph{Hessian (Vol.~5, Ch.~11).}
                $f_{xy}=\tfrac38x^{-1/4}y^{-1/2}$, and
                \[
                  \det\Hess f=\tfrac{3}{64}x^{-1/2}y^{-1}
                  -\tfrac{9}{64}x^{-1/2}y^{-1}
                  =-\tfrac{3}{32}\,x^{-1/2}y^{-1}<0 :
                \]
                \textbf{indefinite} at every interior point, so $f$ is
                neither concave nor convex. Diminishing marginal products
                (negative diagonal) are \emph{not} enough for concavity;
                the cross effect $f_{xy}$ is too strong. Along the diagonal
                the function is $f(t,t)=t^{5/4}$, strictly convex in $t$,
                which refutes concavity directly.
          \item \emph{Level curves and gradient (Vol.~5, Ch.~9, Ch.~10).}
                $x^{3/4}y^{1/2}=c$ gives $y=c^{2}x^{-3/2}$: decreasing convex
                isoquants. $\grad f(1,1)=(\tfrac34,\tfrac12)$, norm
                $\sqrt{\tfrac{9}{16}+\tfrac{4}{16}}=\tfrac{\sqrt{13}}{4}
                \approx0.90$. The isoquant $y=x^{-3/2}$ through $(1,1)$ has
                slope $-\tfrac32$, tangent $(2,-3)$, and
                $\tfrac34\cdot2-\tfrac12\cdot3=0$ \checkmark. The gradient
                points to the upper right, towards higher output, and
                $\sqrt{13}/4$ is the output gain per unit step in that best
                direction.
          \item \emph{Critical points (Vol.~7, Ch.~5).} None in the
                interior, since $f_{x}>0$. The minimum value $0$ is attained
                on the two axes, boundary points of the domain, which the
                first-order condition does not see. (At the corner $(0,0)$
                the gradient from (b) does vanish, but the other minimizers
                on the axes are points where the partials do not even
                exist.)
          \item \emph{Contour sets.} $\{f\ge1\}=\{y\ge x^{-3/2}\}$, the
                region above a convex curve: \textbf{convex}. $\{f\le1\}$ is
                not: $(4,\tfrac18)$ and $(\tfrac14,8)$ lie on $f=1$, but the
                midpoint $(\tfrac{17}{8},\tfrac{65}{16})$ has
                $f>1$ (both coordinates exceed $1$).
          \item \emph{Checklist (Vol.~7, Ch.~3).}
                \emph{Concave: no}, $f(t,t)=t^{5/4}$ is strictly convex.
                \emph{Convex: no}, $f(t,1)=t^{3/4}$ is strictly concave.
                \emph{Quasiconcave: yes}, all upper sets lie above convex
                isoquants (Cobb--Douglas always is).
                \emph{Quasiconvex: no}, by (f).
                \emph{Maximum: no}, unbounded.
                \emph{Minimum: yes}, $0$ on the axes.
                \emph{Increasing: yes}: if $x\le x'$ and $y\le y'$ then
                $x^{3/4}\le x'^{3/4}$ and $y^{1/2}\le y'^{1/2}$, and the
                product of nonnegative factors grows.
                \emph{Decreasing: no}, $f(0,0)=0<1=f(1,1)$.
                \emph{Epigraph convex: no}, $f$ is not convex.
        \end{enumerate}
  \item \textbf{Cobb--Douglas with constant returns.}
        \begin{enumerate}[label=(\alph*),leftmargin=1.8em,itemsep=0.4ex]
          \item \emph{Name (Vol.~5, Ch.~14).} Cobb--Douglas, homogeneous of
                degree $\tfrac13+\tfrac23=1$: \textbf{constant returns to
                scale}. Domain $\R^{2}_{+}$ (closed, convex, unbounded),
                range $[0,\infty)$, since $f(t,t)=t$.
          \item \emph{Derivative and Hessian (Vol.~5, Ch.~10, Ch.~11).}
                On $x,y>0$:
                $f_{x}=\tfrac13x^{-2/3}y^{2/3}$,
                $f_{y}=\tfrac23x^{1/3}y^{-1/3}$, and
                \[
                  \Hess f=\frac29\begin{pmatrix}
                    -x^{-5/3}y^{2/3}& x^{-2/3}y^{-1/3}\\
                    x^{-2/3}y^{-1/3}& -x^{1/3}y^{-4/3}\end{pmatrix},\qquad
                  \det\Hess f=\tfrac{4}{81}x^{-4/3}y^{-2/3}
                  -\tfrac{4}{81}x^{-4/3}y^{-2/3}=0 .
                \]
                Diagonal negative, determinant zero: \textbf{negative
                semidefinite}, not definite. The null direction is the ray
                itself: the first row times $(x,y)$ is
                $\tfrac29\bigl(-x^{-2/3}y^{2/3}+x^{-2/3}y^{2/3}\bigr)=0$,
                and likewise the second. Along a ray $f(tx,ty)=tf(x,y)$ is
                linear, so it has no curvature there.
          \item \emph{Level curves, gradient, Euler (Vol.~5, Ch.~9,
                Ch.~10, Ch.~14).} $x^{1/3}y^{2/3}=c\iff y=c^{3/2}x^{-1/2}$.
                $f(1,8)=1\cdot4=4$, so the isoquant is $y=8x^{-1/2}$, slope
                $-4x^{-3/2}=-4$ at $x=1$. $\grad f(1,8)=
                (\tfrac13\cdot4,\ \tfrac23\cdot\tfrac12)=(\tfrac43,\tfrac13)$,
                norm $\tfrac{\sqrt{17}}{3}\approx1.37$; with tangent
                $(1,-4)$: $\tfrac43-\tfrac43=0$ \checkmark. Steepest ascent
                points mostly in the $x$ direction, because $x$ is the
                scarce factor at $(1,8)$. Euler: $1\cdot\tfrac43+8\cdot
                \tfrac13=4=f(1,8)$ \checkmark.
          \item \emph{Critical points.} None in the interior ($f_{x}>0$);
                the minimum $0$ sits on the axes.
          \item \emph{Contour sets.} $\{f\ge4\}=\{y\ge8x^{-1/2}\}$:
                \textbf{convex}. $\{f\le4\}$: $(16,2)$ and $(\tfrac14,16)$
                lie on $f=4$, the midpoint $(\tfrac{65}{8},9)$ gives
                $f=(\tfrac{65}{8})^{1/3}\cdot9^{2/3}\approx8.7>4$.
                \textbf{Not convex.}
          \item \emph{Checklist (Vol.~7, Ch.~3).}
                \emph{Concave: yes}, the Hessian is negative semidefinite
                on the interior, and concavity carries over to the axes by
                continuity (Cobb--Douglas with $a+b\le1$).
                \emph{Strictly concave: no}, $f(t,t)=t$ is linear along the
                diagonal.
                \emph{Convex: no}, $f(t,1)=t^{1/3}$ is strictly concave.
                \emph{Quasiconcave: yes}, because concave.
                \emph{Quasiconvex: no}, by (e).
                \emph{Maximum: no}. \emph{Minimum: yes}, $0$ on the axes.
                \emph{Increasing: yes}, as in Problem~3.
                \emph{Decreasing: no}.
                \emph{Epigraph convex: no}; the hypograph is convex.
        \end{enumerate}
  \item \textbf{Logarithm of a positive quadratic.}
        \begin{enumerate}[label=(\alph*),leftmargin=1.8em,itemsep=0.4ex]
          \item \emph{Domain and range (Vol.~2, Ch.~1).}
                $q=4+x^{2}+y^{2}\ge4>0$ always, so $D=\R^{2}$. $q$ takes
                every value in $[4,\infty)$ and $\ln$ is increasing, so
                $f(\R^{2})=[\ln4,\infty)$.
          \item \emph{Gradient (Vol.~5, Ch.~10).}
                $\grad f=\Bigl(\dfrac{2x}{q},\dfrac{2y}{q}\Bigr)$,
                continuous everywhere since $q\ne0$: differentiable on all
                of $\R^{2}$.
          \item \emph{Level curves and gradient (Vol.~5, Ch.~9).}
                $f=c\iff x^{2}+y^{2}=\mathrm{e}^{c}-4$: circles about the
                origin for $c>\ln4$, the origin itself for $c=\ln4$, empty
                below. Through $(1,1)$ runs the circle $x^{2}+y^{2}=2$
                (level $\ln6$). $\grad f(1,1)=(\tfrac26,\tfrac26)
                =(\tfrac13,\tfrac13)$, norm $\tfrac{\sqrt2}{3}$; it is a
                multiple of the radius vector $(1,1)$, and a radius is
                orthogonal to a circle: tangent $(1,-1)$,
                $\tfrac13-\tfrac13=0$ \checkmark. It points outwards,
                towards larger circles and higher values.
          \item \emph{Hessian (Vol.~5, Ch.~11).} Quotient rule:
                \[
                  f_{xx}=\frac{2(4-x^{2}+y^{2})}{q^{2}},\quad
                  f_{yy}=\frac{2(4+x^{2}-y^{2})}{q^{2}},\quad
                  f_{xy}=-\frac{4xy}{q^{2}},
                \]
                \[
                  \det\Hess f=\frac{4\bigl[16-(x^{2}-y^{2})^{2}\bigr]
                  -16x^{2}y^{2}}{q^{4}}
                  =\frac{4\bigl[16-(x^{2}+y^{2})^{2}\bigr]}{q^{4}} .
                \]
                At $(0,0)$: $\Hess f=\mathrm{diag}(\tfrac12,\tfrac12)$,
                positive definite. At $(1,1)$ ($q=6$):
                $\begin{pmatrix}2/9&-1/9\\-1/9&2/9\end{pmatrix}$,
                $\det=\tfrac{1}{27}>0$, $f_{xx}>0$: positive definite. At
                $(2,2)$ ($q=12$):
                $\begin{pmatrix}1/18&-1/9\\-1/9&1/18\end{pmatrix}$,
                $\det=-\tfrac{1}{108}<0$: indefinite.
          \item \emph{Critical points (Vol.~7, Ch.~5).} $2x/q=0$ and
                $2y/q=0$ give only $(0,0)$; positive definite Hessian: strict
                local minimum $\ln4$, and global, since $f\ge\ln4$ by (a).
          \item \emph{Contour sets.} $\{f\le\ln8\}=\{x^{2}+y^{2}\le4\}$, a
                disk: \textbf{convex}. $\{f\ge\ln8\}=\{x^{2}+y^{2}\ge4\}$,
                the outside of the disk: it contains $(2,0)$ and $(-2,0)$
                but not their midpoint $(0,0)$. \textbf{Not convex.}
          \item \emph{Checklist (Vol.~7, Ch.~3).}
                \emph{Concave: no}, the Hessian is positive definite at the
                origin.
                \emph{Convex: no}, it is indefinite at $(2,2)$; along the
                $x$-axis $g(t)=\ln(4+t^{2})$ has
                $g''(t)=\dfrac{2(4-t^{2})}{(4+t^{2})^{2}}<0$ for $|t|>2$.
                \emph{Quasiconvex: yes}, every lower set is a disk (or
                empty, or a point).
                \emph{Quasiconcave: no}, by (f).
                \emph{Maximum: no}, unbounded above.
                \emph{Minimum: yes}, $\ln4$ at $(0,0)$.
                \emph{Increasing: no}, $(-1,0)\le(0,0)$ but
                $f(-1,0)=\ln5>\ln4=f(0,0)$.
                \emph{Decreasing: no}, $(0,0)\le(1,0)$ but
                $f(1,0)=\ln5>f(0,0)$.
                \emph{Epigraph convex: no}, $f$ is not convex.
                \emph{Convex on a disk:} the determinant is positive
                exactly when $x^{2}+y^{2}<2^{2}$, and there
                $f_{xx}=2(4-x^{2}+y^{2})/q^{2}>0$ because $x^{2}<4$. So $f$
                is convex on the disk of radius $2$ (closed, by continuity),
                and on no larger disk, since the determinant turns negative
                just outside it. This is the typical shape of a
                quasiconvex but not convex function: a bowl at the bottom
                that flattens into a trumpet.
        \end{enumerate}
\end{solutions}


\end{document}
"
\ifdefined\KindleEdition
  \documentclass[10pt,oneside]{book}
\else
  \documentclass[11pt,a4paper,twoside,openright]{book}
\fi
% Shared preamble for the Mixed Problems training sheet.
\usepackage[T1]{fontenc}
\usepackage[utf8]{inputenc}
\usepackage{lmodern}
\usepackage[a4paper,margin=2.4cm]{geometry}
\usepackage{amsmath,amssymb,amsthm}
\usepackage{microtype}
\usepackage{enumitem}
\usepackage{xcolor}
\usepackage[most]{tcolorbox}
\usepackage{fancyhdr}
\usepackage[hidelinks]{hyperref}

\definecolor{accent}{HTML}{1F4E79}
\definecolor{soft}{HTML}{F4F7FA}

% --- function graphs ------------------------------------------------------
\usepackage{pgfplots}
\pgfplotsset{compat=1.18}
\usetikzlibrary{arrows.meta}

% One consistent look for every function graph in these books.
\pgfplotsset{
  fnaxis/.style={
    width=5.9cm, height=4.6cm,
    axis lines=middle,
    axis line style={-{Stealth[length=4pt]}, gray!65!black, line width=0.3pt},
    tick style={gray!65!black},
    tick label style={font=\tiny, gray!50!black},
    label style={font=\scriptsize},
    xlabel={$x$}, ylabel={$y$},
    every axis x label/.style={at={(ticklabel* cs:1.0)}, anchor=west},
    every axis y label/.style={at={(ticklabel* cs:1.0)}, anchor=south east,
      xshift=-1pt},
    every axis plot/.append style={accent, very thick, smooth, samples=201},
    clip=true, enlargelimits=false,
    title style={font=\scriptsize, yshift=4pt},
  },
  fnaxis/.append style={/pgfplots/legend style={font=\tiny, draw=none,
    fill=none, at={(0.98,0.03)}, anchor=south east}},
}

% A row of small graphs, evenly spaced with \hfill between them, plus an
% optional caption line: \begin{gallery} <pic>\hfill<pic> \gcap{...} \end{gallery}
\newenvironment{gallery}
  {\par\medskip\noindent\ignorespaces}
  {\par\medskip\noindent}
\newcommand{\gcap}[1]{\par\smallskip{\centering\footnotesize #1\par}}

% --- plane regions, implicit curves, surfaces -----------------------------
\usepgfplotslibrary{fillbetween}
\definecolor{regionfill}{HTML}{9DB4D8}
\pgfplotsset{colormap={house}{color=(soft) color=(accent!45) color=(accent)}}

\pgfplotsset{
  % Squared, gridded axis for plane regions and implicitly given curves.
  geoaxis/.style={
    fnaxis,
    width=5.9cm, height=5.9cm,
    grid=both,
    grid style={gray!22, line width=0.2pt},
    minor tick num=1,
  },
  % Level-curve maps: many thin curves, no single curve emphasised.
  contouraxis/.style={
    geoaxis,
    every axis plot/.append style={accent!75, line width=0.5pt, smooth,
      samples=201},
  },
  % Surfaces: \addplot3[surf, domain=..., y domain=...] {expr};
  surfaxis/.style={
    width=6.0cm, height=5.0cm,
    view={130}{22},
    colormap name=house,
    ticks=none,
    xlabel={$x$}, ylabel={$y$}, zlabel={$f$},
    label style={font=\tiny, gray!50!black},
    title style={font=\scriptsize},
    every axis plot/.append style={samples=25, shader=faceted interp,
      draw=accent!35, line width=0.1pt},
  },
}

% Shorthand for a shaded region between two named paths.
% \region{<path A>}{<path B>}{<soft clip options>}
\newcommand{\region}[3]{%
  \addplot[regionfill, opacity=0.55, forget plot, sharp plot]
    fill between[of=#1 and #2, #3];}

% --- rule boxes -----------------------------------------------------------
\newtcolorbox{rulebox}[1]{%
  breakable, enhanced, colback=soft, colframe=accent,
  boxrule=0.4pt, left=7pt, right=7pt, top=6pt, bottom=6pt,
  fonttitle=\bfseries, coltitle=white, colbacktitle=accent,
  title={#1}}

% --- worked examples ------------------------------------------------------
% One fully narrated solution per chapter, carrying the chapter's central
% technique and the mistakes that technique invites:
%   \begin{worked}{Task} ... \wstep{Step 1.} ... \wtrap{...} ... \end{worked}
\newtcolorbox{worked}[1]{%
  breakable, enhanced, colback=white, colframe=accent!65,
  boxrule=0.4pt, left=7pt, right=7pt, top=6pt, bottom=6pt,
  fonttitle=\bfseries, coltitle=accent, colbacktitle=soft,
  title={#1}}

% Run-in step heading inside a worked example.
\newcommand{\wstep}[1]{\par\smallskip\noindent\textbf{#1}\enspace\ignorespaces}

% A pitfall, set off from the running solution so it can be found again.
\newcommand{\wtrap}[1]{%
  \par\smallskip
  {\leftskip=1.3em\noindent
   \llap{\textcolor{accent}{\rule[0.22ex]{0.9em}{0.6pt}}\hspace{0.4em}}%
   \textbf{Trap.}\enspace\itshape #1\par}%
  \smallskip}

% --- problem / solution lists --------------------------------------------
\newlist{problems}{enumerate}{1}
\setlist[problems]{label=\textbf{\arabic*.}, leftmargin=2.1em,
  itemsep=0.55ex, parsep=0pt, topsep=0.8ex}
\newlist{solutions}{enumerate}{1}
\setlist[solutions]{label=\textbf{\arabic*.}, leftmargin=2.1em,
  itemsep=1.0ex, parsep=0pt, topsep=0.8ex}

% --- running heads --------------------------------------------------------
\pagestyle{fancy}
\fancyhf{}
\fancyhead[LE,RO]{\thepage}
\fancyhead[RE]{\small\itshape\nouppercase{\leftmark}}
\fancyhead[LO]{\small\itshape\nouppercase{\rightmark}}
\setlength{\headheight}{14pt}
\renewcommand{\headrulewidth}{0.3pt}
\renewcommand{\chaptermark}[1]{\markboth{\chaptername\ \thechapter.\ #1}{}}
\renewcommand{\sectionmark}[1]{\markright{\thesection\ #1}}

% --- shorthands used throughout -------------------------------------------
\newcommand{\dd}[2]{\frac{\mathrm{d}#1}{\mathrm{d}#2}}
\newcommand{\pd}[2]{\frac{\partial #1}{\partial #2}}
\newcommand{\pdd}[3]{\frac{\partial^{2} #1}{\partial #2\,\partial #3}}
\newcommand{\R}{\mathbb{R}}
\newcommand{\grad}{\nabla}
\DeclareMathOperator{\sgn}{sgn}
\DeclareMathOperator{\Hess}{H}

% German gloss, used on first mention of a technical term
\newcommand{\de}[1]{\emph{#1}}

\renewcommand*{\thisvolume}{mixed}

\title{\Huge Mixed Problems\\[0.4em]
  \large Cross-Chapter Sets with Worked Solutions}
\author{}
\date{}

\begin{document}
\frontmatter
\maketitle

\chapter*{About this sheet}
\addcontentsline{toc}{chapter}{About this sheet}

Every other volume in this series drills one technique at a time. That is
the right way to learn a method and the wrong way to prepare for an
examination, because a chapter heading tells you what to do before you have
read the question. An exam does not. The single most common way to lose
marks is not botched arithmetic but reaching for the wrong tool: setting up
a Lagrangian where the constraint does not bind, differentiating where a
limit was wanted, solving a system by inverting a matrix that turns out to
be singular.

This sheet exists to remove that crutch. Each set draws on the whole
sequence --- algebra and functions, matrices and elimination, limits and
L'H\^opital, curve discussion, multivariable calculus, constrained
optimization, linear programming --- in no announced order, and no problem
says which chapter it came from. Deciding that \emph{is} the exercise.

Work a set in one sitting, with a time limit and without looking anything
up. Sixty to ninety minutes for twenty problems is a realistic target. Then
read the solutions, and for every problem you got wrong, note not the
arithmetic but the \emph{decision}: which signal in the question should have
pointed at which method? That note is worth more than redoing the
computation.

Solutions name the technique in their first words, so the sets can also be
used the other way round --- read only the first line of each solution as a
drill in classification.

The fourth set works the other way. It announces the methods, all of them,
and takes a single function of two variables through the whole chain:
domain, range, derivative, level curves, gradient, Hessian, critical
points, contour sets and a checklist of properties. There the difficulty is
not choosing a tool but keeping a dozen answers consistent with one
another.

\tableofcontents

\mainmatter
\chapter{Mixed Set A}

Twenty problems, no order, no announced methods. Target time: 75 minutes.

\section{Worked Example}

The problems below announce no method. That is the point: on an
examination, choosing the technique is most of the work, and it is the one
step no drill chapter can train, because a drill chapter has already told
you the answer in its title.

The example here is \emph{not} one of the twenty. It is a typical example of
the kind, a single function examined from five directions, shown so
that you can watch the recognition step happen before you have to do it
yourself.

\begin{worked}{A five-part question on $f(x)=\dfrac{\ln x}{x}$: domain, both edges, extremum, inflection, and $\int_{1}^{\infty}f$}
\wstep{Step 1. Read the parts and name the tools before computing anything.}
\begin{center}\footnotesize
\begin{tabular}{lll}
  \hline
  part & tool & volume\\
  \hline
  domain & obstruction list & 2\\
  edges & limits, growth hierarchy, L'H\^opital & 4, 5\\
  extremum & $f'=0$, then $f''$ & 5\\
  inflection & sign change of $f''$ & 5\\
  convergence & improper integral, substitution & 6\\
  \hline
\end{tabular}
\end{center}
A minute spent on this table keeps the rest of the solution organised.

\wstep{Step 2. Domain.}
$\ln x$ needs $x>0$, and the denominator needs $x\neq0$. The first
condition subsumes the second:
\[
  D_{f}=(0,\infty).
\]

\wstep{Step 3. The left edge --- and it is \emph{not} indeterminate.}
As $x\to0^{+}$: $\ln x\to-\infty$ and $1/x\to+\infty$. The form is
$(-\infty)\cdot(+\infty)$, which is determinate:
\[
  \lim_{x\to0^{+}}\frac{\ln x}{x}=-\infty .
\]
\wtrap{Reaching for L'H\^opital here is the standard reflex and it is wrong
--- the rule needs $0/0$ or $\infty/\infty$, and this is neither. Always
name the form before choosing a tool. Here both factors push the same way,
so there is nothing to resolve.}

\wstep{Step 4. The right edge --- and here it \emph{is} indeterminate.}
As $x\to\infty$ the form is $\infty/\infty$, so either L'H\^opital or the
growth hierarchy applies:
\[
  \lim_{x\to\infty}\frac{\ln x}{x}
  \overset{\infty/\infty}{=}\lim_{x\to\infty}\frac{1/x}{1}=0 ,
\]
consistent with $\ln x\ll x$. So $y=0$ is a horizontal asymptote.

\wstep{Step 5. The extremum.}
\[
  f'(x)=\frac{\tfrac1x\cdot x-\ln x}{x^{2}}=\frac{1-\ln x}{x^{2}} .
\]
$f'=0$ when $\ln x=1$, i.e.\ $x=\mathrm{e}$, and
$f(\mathrm{e})=1/\mathrm{e}\approx0.368$.

The denominator $x^{2}$ is positive throughout the domain, so the sign of
$f'$ is the sign of $1-\ln x$: positive for $x<\mathrm{e}$, negative
after. A \textbf{maximum}, and since it is the only critical point and the
function rises to it from $-\infty$ and falls to $0$, it is the
\emph{global} maximum.
\wtrap{With a single critical point, the sign of $f'$ on either side settles
the classification faster than computing $f''$ --- and it settles the global
question too, which $f''$ never does.}

\wstep{Step 6. The inflection.}
\[
  f''(x)=\frac{2\ln x-3}{x^{3}} .
\]
This vanishes at $\ln x=\tfrac32$, i.e.\ $x=\mathrm{e}^{3/2}\approx4.482$,
and changes sign there (the denominator is positive), so it is a genuine
inflection at
\[
  \Bigl(\mathrm{e}^{3/2},\ \tfrac{3}{2}\mathrm{e}^{-3/2}\Bigr)\approx(4.48,\ 0.335).
\]
Check the ordering: $\mathrm{e}\approx2.72<4.48$, so the curve peaks and
\emph{then} changes curvature --- which is what a graph tending to zero from
above must do.

\wstep{Step 7. The integral.}
Substituting $u=\ln x$, $\mathrm{d}u=\mathrm{d}x/x$:
\[
  \int_{1}^{b}\frac{\ln x}{x}\,\mathrm{d}x
  =\int_{0}^{\ln b}u\,\mathrm{d}u
  =\frac{(\ln b)^{2}}{2}\xrightarrow[b\to\infty]{}\infty .
\]
\[
  \boxed{\int_{1}^{\infty}\frac{\ln x}{x}\,\mathrm{d}x\ \text{diverges.}}
\]
\wtrap{The integrand tends to $0$ --- Step~4 proved it --- and the integral
still diverges. ``The function goes to zero'' is never sufficient for
convergence, exactly as ``the terms go to zero'' was never sufficient for a
series. Here $f$ behaves like $1/x$ up to a logarithm, and $\int_{1}^{\infty}
\mathrm{d}x/x$ diverges.}

\wstep{Step 8. Assemble the picture, and check it against every part.}
Rising from $-\infty$ at the origin, crossing zero at $x=1$, peaking at
$(\mathrm{e},1/\mathrm{e})$, inflecting at $x\approx4.48$, and decaying to
$0$ --- with a tail thick enough to have infinite area. Every part of the
answer must be visible in that description, and if one is not, something is
wrong.
\begin{center}
\begin{tikzpicture}
\begin{axis}[fnaxis, width=8.0cm, height=4.4cm, domain=0.13:12,
             samples=201, xmin=-0.6, xmax=12.6, ymin=-2.2, ymax=0.9,
             xtick={1,2.718,4.482,8}, xticklabels={$1$,$\mathrm{e}$,$\mathrm{e}^{3/2}$,$8$},
             ytick={0.368}, yticklabels={$1/\mathrm{e}$}]
  \addplot {ln(x)/x};
  \addplot[only marks, mark=*, mark size=1.2pt, accent!70, forget plot]
    coordinates {(2.71828,0.36788) (4.48169,0.33470)};
\end{axis}
\end{tikzpicture}
\end{center}
\end{worked}

\section{Problems}

\begin{problems}
  \item Give the maximal domain of
        $f(x)=\dfrac{\sqrt{x^{2}-9}}{\ln(x-2)}$.
  \item Solve $3^{2x}-4\cdot3^{x}+3=0$.
  \item Find the inverse of $f(x)=\dfrac{2x-1}{x+3}$ and state its domain.
  \item Describe the set $\{(x,y): x^{2}+y^{2}\le25,\ y\ge x\}$ and decide
        whether it is convex.
  \item Give the level curves of $f(x,y)=x^{2}-4y$ at $c=0$ and $c=4$, and
        name the family.
  \item Evaluate $\displaystyle\lim_{x\to0}
        \frac{\mathrm{e}^{3x}-1-3x}{x^{2}}$.
  \item Evaluate $\displaystyle\lim_{x\to\infty}\frac{(\ln x)^{3}}{x}$.
  \item For $f(x)=\dfrac{x}{x^{2}-1}$: give the domain, the symmetry, all
        asymptotes, and the intervals of monotonicity.
  \item Solve $x+2y-z=3$, \;$2x-y+z=5$, \;$3x+y+2z=14$.
  \item Give the rank of
        $A=\begin{pmatrix}1&2&3\\2&4&6\\1&0&1\end{pmatrix}$ and say whether
        $A$ is invertible.
  \item Find the eigenvalues of
        $\begin{pmatrix}2&1\\1&2\end{pmatrix}$.
  \item Determine the definiteness of the Hessian of
        $f(x,y)=x^{2}+4xy+y^{2}$ and what it implies about the critical
        point at the origin.
  \item Find and classify all critical points of
        $f(x,y)=x^{3}-3xy+y^{3}$.
  \item Maximise $U=xy$ subject to $3x+4y=24$, and give the multiplier.
  \item Using the previous answer only, estimate the change in maximal
        utility when the income $24$ rises to $25$, then compute the exact
        change.
  \item Maximise $f(x,y)=x+y$ subject to $x^{2}+y^{2}\le2$, $x,y\ge0$.
  \item Maximise $3x+2y$ subject to $x+y\le4$, $x\le3$, $x,y\ge0$.
  \item Evaluate $\displaystyle\int x\,\mathrm{e}^{x}\,\mathrm{d}x$.
  \item Compute the price elasticity of demand for $q(p)=100/p^{2}$.
  \item At a continuously compounded rate of $5\%$, how long does a
        capital take to double?
\end{problems}

\section{Solutions}

\begin{solutions}
  \item \textbf{Maximal domain.} Three conditions: $x^{2}-9\ge0$ gives
        $|x|\ge3$; $\ln(x-2)$ needs $x>2$; and the denominator must not
        vanish, so $x-2\neq1$, i.e. $x\neq3$. Intersecting,
        $x\ge3$ and $x>2$ and $x\neq3$: \;$D=(3,\infty)$.
  \item \textbf{Substitution in an exponential equation.} With
        $u=3^{x}>0$: $u^{2}-4u+3=0$, so $u=1$ or $u=3$, giving
        $x=0$ and $x=1$.
  \item \textbf{Inverse function.} $y(x+3)=2x-1\Rightarrow x(y-2)=-1-3y$,
        so $x=\dfrac{3y+1}{2-y}$ and
        $f^{-1}(x)=\dfrac{3x+1}{2-x}$ with
        $D(f^{-1})=\R\setminus\{2\}$ --- which is exactly the range of
        $f$, as it must be.
  \item \textbf{Region in the plane.} The half of the disc of radius $5$
        lying on or above the line $y=x$. It is the intersection of two
        convex sets (a disc and a half-plane), hence \textbf{convex}.
  \item \textbf{Level curves.} $x^{2}-4y=c\Rightarrow y=\dfrac{x^{2}-c}{4}$:
        for $c=0$, $y=x^{2}/4$; for $c=4$, $y=\dfrac{x^{2}-4}{4}$. All
        level curves are vertical translates of the same parabola.
  \item \textbf{L'H\^opital, twice.} $\dfrac{3\mathrm{e}^{3x}-3}{2x}$ is
        still $0/0$, then
        $\dfrac{9\mathrm{e}^{3x}}{2}\to\dfrac92$.
  \item \textbf{Growth hierarchy.} $(\ln x)^{3}\ll x$, so the limit is
        $0$. (Or three applications of L'H\^opital, each reducing the power
        of the logarithm by one.)
  \item \textbf{Curve discussion.} Domain $\R\setminus\{-1,1\}$; the
        function is \emph{odd}; single root $x=0$; poles at $x=\pm1$;
        horizontal asymptote $y=0$. And
        $f'=\dfrac{(x^{2}-1)-x\cdot2x}{(x^{2}-1)^{2}}
        =\dfrac{-(x^{2}+1)}{(x^{2}-1)^{2}}<0$ everywhere, so $f$ is
        strictly decreasing on each of the three branches, with \emph{no}
        extrema.
  \item \textbf{Elimination.} Adding the first two equations gives
        $3x+y=8$; substituting into the third, $8+2z=14$ so $z=3$. Then
        $x+2y=6$ and $y=8-3x$ give $x=2$, $y=2$. Solution $(2,2,3)$.
        (Check the second: $4-2+3=5$ \checkmark.)
  \item \textbf{Rank by row reduction.} $R_{2}=2R_{1}$, so that row
        vanishes; $R_{3}-R_{1}=(0,-2,-2)$ survives.
        $\mathrm{rk}(A)=2<3$, so $A$ is \textbf{singular} --- equivalently
        $\det A=0$.
  \item \textbf{Characteristic polynomial.}
        $(2-\lambda)^{2}-1=0\Rightarrow\lambda-2=\pm1$, so
        $\lambda_{1}=3$ and $\lambda_{2}=1$. Both positive, so the matrix
        is positive definite.
  \item \textbf{Definiteness.}
        $\Hess f=\begin{pmatrix}2&4\\4&2\end{pmatrix}$ with
        $\det=4-16=-12<0$: \textbf{indefinite}. So the origin is a
        \textbf{saddle point}, not an extremum.
  \item \textbf{Critical points in two variables.}
        $f_{x}=3x^{2}-3y=0$ gives $y=x^{2}$; $f_{y}=-3x+3y^{2}=0$ gives
        $x=y^{2}$. Hence $x=x^{4}$, so $x(x^{3}-1)=0$ and $x=0$ or $x=1$:
        the points $(0,0)$ and $(1,1)$. With $f_{xx}=6x$, $f_{yy}=6y$,
        $f_{xy}=-3$: at $(0,0)$ the determinant is $0-9=-9<0$, a
        \textbf{saddle}; at $(1,1)$ it is $36-9=27>0$ with $f_{xx}=6>0$, a
        \textbf{local minimum} with $f=-1$.
  \item \textbf{Lagrange.} $\mathcal{L}=xy+\lambda(24-3x-4y)$ gives
        $y=3\lambda$ and $x=4\lambda$, so $4y=3x$. With the constraint,
        $3x+3x=24$, hence $x^{*}=4$, $y^{*}=3$, $U^{*}=12$ and
        $\lambda^{*}=1$.
  \item \textbf{Envelope theorem.} $\mathrm{d}U^{*}/\mathrm{d}M
        =\lambda^{*}=1$, so the estimate is $+1$. Exactly:
        $U^{*}(M)=M^{2}/48$, so
        $U^{*}(25)-U^{*}(24)=\dfrac{625-576}{48}=\dfrac{49}{48}
        \approx1.021$. The multiplier is the first-order term.
  \item \textbf{Kuhn--Tucker (or symmetry).} The objective increases in
        both variables, so the circle constraint binds; by symmetry
        $x=y$, and $2x^{2}=2$ gives $x=y=1$. The non-negativity
        constraints are slack, and $f^{*}=2$. (Formally
        $\grad f=(1,1)^{T}=\lambda(2x,2y)^{T}$ forces $x=y$ with
        $\lambda=\tfrac12>0$.)
  \item \textbf{Linear programming, graphically.} The feasible vertices
        are $(0,0)$, $(3,0)$, $(3,1)$ and $(0,4)$, where the objective
        takes the values $0$, $9$, $11$ and $8$. The maximum is
        $\mathbf{11}$ at $(3,1)$, where both $x\le3$ and $x+y\le4$ bind.
  \item \textbf{Integration by parts.} With $u=x$,
        $v'=\mathrm{e}^{x}$:
        $\int x\mathrm{e}^{x}\mathrm{d}x
        =x\mathrm{e}^{x}-\int\mathrm{e}^{x}\mathrm{d}x
        =(x-1)\mathrm{e}^{x}+C$.
  \item \textbf{Elasticity.}
        $\varepsilon=\dfrac{q'(p)\,p}{q(p)}
        =\dfrac{-200p^{-3}\cdot p}{100p^{-2}}=-2$: constant, independent of
        the price --- the signature of a power demand function.
  \item \textbf{Continuous compounding.}
        $\mathrm{e}^{0.05t}=2\Rightarrow t=\dfrac{\ln2}{0.05}
        \approx13.86$ years. (The ``rule of $70$'': $70/5=14$.)
\end{solutions}

\chapter{Mixed Set B}

Twenty problems, no order, no announced methods. Target time: 75 minutes.

\section{Worked Example}

Again the example is not one of the twenty. This one is a constrained
optimisation, one of the most frequently examined topics, worked
with every check that the marker expects to see.

\begin{worked}{Maximise $f(x,y)=\ln x+2\ln y$ subject to $x+y=12$, with $x,y>0$: solve, certify, and interpret the multiplier}
\wstep{Step 1. Recognise the type before writing anything.}
Two variables, one equality constraint, a differentiable objective: this is
Lagrange. The constraint is linear and could be substituted instead --- keep
that in reserve as the check.

\wstep{Step 2. Set up and differentiate.}
\[
  \mathcal L=\ln x+2\ln y+\lambda(12-x-y),
\]
\[
  \mathcal L_{x}=\frac1x-\lambda=0,
  \qquad
  \mathcal L_{y}=\frac2y-\lambda=0,
  \qquad
  x+y=12 .
\]

\wstep{Step 3. Eliminate $\lambda$.}
\[
  \frac1x=\frac2y
  \quad\Longrightarrow\quad
  y=2x .
\]
The optimal bundle puts twice as much into $y$ as into $x$ --- the ratio of
the coefficients $2$ and $1$ in the objective, exactly as Cobb--Douglas
shares behave, and $\ln$-linear objectives are Cobb--Douglas in disguise.

\wstep{Step 4. Impose the constraint.}
\[
  x+2x=12
  \quad\Longrightarrow\quad
  \boxed{x^{*}=4,\quad y^{*}=8,\quad \lambda^{*}=\tfrac14 }
\]
\[
  f^{*}=\ln4+2\ln8=\ln4+\ln64=\ln256\approx5.545 .
\]

\wstep{Step 5. Certify it with the bordered Hessian.}
\[
  \mathcal L_{xx}=-\frac{1}{x^{2}}=-\tfrac{1}{16},
  \qquad
  \mathcal L_{yy}=-\frac{2}{y^{2}}=-\tfrac{1}{32},
  \qquad
  \mathcal L_{xy}=0,
  \qquad
  g_{x}=g_{y}=1 .
\]
\[
  \overline H=\begin{pmatrix}0&1&1\\1&-\tfrac{1}{16}&0\\1&0&-\tfrac{1}{32}\end{pmatrix},
  \qquad
  \det\overline H=\tfrac{1}{32}+\tfrac{1}{16}=\tfrac{3}{32}>0 ,
\]
which for $n=2$, $m=1$ means a constrained \textbf{maximum}. \checkmark

\wstep{Step 6. Check again by substitution, which is available here.}
$y=12-x$ gives $h(x)=\ln x+2\ln(12-x)$, so
\[
  h'(x)=\frac1x-\frac{2}{12-x}=0
  \;\Longrightarrow\;
  12-x=2x
  \;\Longrightarrow\;
  x=4 ,
\]
\[
  h''(x)=-\frac{1}{x^{2}}-\frac{2}{(12-x)^{2}}<0\ \text{everywhere:\ concave, so the maximum is global.}
\]
\wtrap{Two independent routes to the same answer is the best use of spare
examination time --- far better than re-reading the first one. And here the
substitution route delivers something the bordered Hessian does not: a
\emph{global} conclusion, because $h''<0$ on the whole feasible interval,
not merely at the candidate point.}

\wstep{Step 7. Verify the multiplier's meaning.}
With a general budget $M$ in place of $12$: $x^{*}=M/3$, $y^{*}=2M/3$, so
\[
  V(M)=\ln\frac M3+2\ln\frac{2M}{3}=3\ln M+\text{constant},
  \qquad
  \dd{V}{M}=\frac3M .
\]
At $M=12$ this is $\tfrac14=\lambda^{*}$. \checkmark\ One more unit of budget
raises the maximised objective by about $0.25$.

\wstep{Step 8. The checks a marker looks for.}
\begin{center}\footnotesize
\begin{tabular}{ll}
  \hline
  check & here\\
  \hline
  constraint satisfied & $4+8=12$ \checkmark\\
  domain respected & $x,y>0$ \checkmark\\
  second-order condition & $\det\overline H>0$ \checkmark\\
  multiplier sign & $\lambda>0$: relaxing helps \checkmark\\
  interpretation stated & $\lambda=\partial V/\partial M$ \checkmark\\
  \hline
\end{tabular}
\end{center}
Every one of these is a line, and together they are usually worth more marks
than the algebra that produced the numbers.
\end{worked}

\section{Problems}

\begin{problems}
  \item Simplify $\dfrac{x^{2}-4}{x^{2}-x-6}$ and state where the
        simplification is valid.
  \item Solve $\log_{2}x+\log_{2}(x-2)=3$.
  \item Give the range of $f(x)=\dfrac{x^{2}+1}{x}$ on $x>0$.
  \item Write a system of inequalities for the part of the elliptical disc
        $\dfrac{x^{2}}{4}+\dfrac{y^{2}}{9}\le1$ lying in the second
        quadrant.
  \item Two functions have concentric circles as level curves. In one the
        circles are equally spaced, in the other they crowd outwards. Which
        is $\sqrt{x^{2}+y^{2}}$ and which is $x^{2}+y^{2}$?
  \item Evaluate $\displaystyle\lim_{x\to0^{+}}x^{\sin x}$.
  \item Evaluate $\displaystyle\lim_{x\to2}\frac{x^{3}-8}{x^{2}-4}$.
  \item For $f(x)=\dfrac{\mathrm{e}^{x}}{x}$: give the domain, the
        asymptotes, the monotonicity and any extremum.
  \item Solve $2x+y=7$, \;$4x+2y=14$, and describe the solution set.
  \item Compute
        $\det\begin{pmatrix}1&2&0\\0&3&1\\2&0&4\end{pmatrix}$.
  \item Solve $\begin{pmatrix}2&1\\1&3\end{pmatrix}
        \mathbf{x}=\binom{5}{10}$.
  \item Give $\grad f(1,2)$ for $f(x,y)=x^{2}y^{3}$.
  \item Give the directional derivative of the same $f$ at $(1,2)$ in the
        direction of $(3,4)^{T}$.
  \item Use the total differential to estimate $f(1.1,1.9)$ for the same
        $f$, and compare with the exact value.
  \item Find $\dd{y}{x}$ at $(1,2)$ on the curve $x^{2}+xy+y^{2}=7$.
  \item Minimise $2x+3y$ subject to $xy=54$, with $x,y>0$.
  \item Maximise $\ln x+y$ subject to $x+y\le4$, $x>0$, $y\ge0$.
  \item Maximise $4x+3y$ subject to $x+y\le10$, $2x+y\le16$, $x,y\ge0$.
  \item What is $1000$ received in five years worth today at $6\%$ annual
        interest?
  \item Sum the series $3+\tfrac32+\tfrac34+\tfrac38+\cdots$.
\end{problems}

\section{Solutions}

\begin{solutions}
  \item \textbf{Factor and cancel.}
        $\dfrac{(x-2)(x+2)}{(x-3)(x+2)}=\dfrac{x-2}{x-3}$, valid for
        $x\neq-2$ and $x\neq3$. The point $x=-2$ remains excluded from the
        domain even though the factor cancelled --- a hole, not a pole.
  \item \textbf{Logarithm law, then check the domain.}
        $\log_{2}\bigl(x(x-2)\bigr)=3\Rightarrow x^{2}-2x=8$, so
        $x^{2}-2x-8=0$ and $x=4$ or $x=-2$. The domain requires $x>2$, so
        $x=-2$ is rejected: \;$x=4$.
  \item \textbf{Range via AM--GM (or calculus).}
        $f(x)=x+\dfrac1x$, and for $x>0$, $x+1/x\ge2$ with equality at
        $x=1$. As $x\to0^{+}$ and $x\to\infty$ the function grows without
        bound, so $W=[2,\infty)$.
  \item \textbf{Region.} $\dfrac{x^{2}}{4}+\dfrac{y^{2}}{9}\le1$,
        \;$x\le0$, \;$y\ge0$.
  \item \textbf{Contour spacing.} Equally spaced circles belong to the
        \textbf{cone} $\sqrt{x^{2}+y^{2}}$, whose level $c$ has radius $c$.
        Circles crowding outwards belong to the \textbf{paraboloid}
        $x^{2}+y^{2}$, whose level $c$ has radius $\sqrt{c}$: the surface
        steepens, so equal steps in height need ever smaller steps in
        radius.
  \item \textbf{Exponential form, take logs.}
        $\ln y=\sin x\ln x$, and since $\sin x\sim x$ near $0$ this behaves
        like $x\ln x\to0$. Hence the limit is $\mathrm{e}^{0}=1$.
  \item \textbf{L'H\^opital (or factor).}
        $\dfrac{3x^{2}}{2x}\to\dfrac{12}{4}=3$. Factoring gives the same:
        $\dfrac{(x-2)(x^{2}+2x+4)}{(x-2)(x+2)}\to\dfrac{12}{4}$.
  \item \textbf{Curve discussion.} Domain $\R\setminus\{0\}$. As
        $x\to0^{+}$, $f\to+\infty$; as $x\to0^{-}$, $f\to-\infty$: a pole
        at $x=0$. As $x\to-\infty$, $f\to0^{-}$, so $y=0$ is a horizontal
        asymptote on the left only; as $x\to+\infty$, $f\to+\infty$. And
        $f'=\dfrac{\mathrm{e}^{x}(x-1)}{x^{2}}$, negative for $x<1$
        (excluding $0$) and positive for $x>1$: decreasing, then
        increasing, with a local (and global, on $x>0$) minimum at
        $(1,\mathrm{e})$.
  \item \textbf{Elimination.} The second equation is exactly twice the
        first, so it carries no information: one pivot, two unknowns.
        \textbf{Infinitely many} solutions, $y=7-2x$ with $x$ free ---
        the line itself.
  \item \textbf{Cofactor expansion along the first row.}
        $1\cdot(3\cdot4-1\cdot0)-2\cdot(0\cdot4-1\cdot2)+0
        =12+4=\mathbf{16}$.
  \item \textbf{Inverse or elimination.} $\det=6-1=5$, so
        $\mathbf{x}=\dfrac15\begin{pmatrix}3&-1\\-1&2\end{pmatrix}
        \binom{5}{10}=\dfrac15\binom{5}{15}=\binom{1}{3}$.
        (Check: $2+3=5$ \checkmark, $1+9=10$ \checkmark.)
  \item \textbf{Partial derivatives.} $f_{x}=2xy^{3}$ and
        $f_{y}=3x^{2}y^{2}$, so at $(1,2)$:
        $\grad f=(2\cdot8,\;3\cdot4)^{T}=(16,12)^{T}$.
  \item \textbf{Directional derivative --- normalise first.}
        $\|(3,4)\|=5$, so $\mathbf{u}=(0.6,0.8)^{T}$ and
        $D_{\mathbf{u}}f=(16,12)\cdot(0.6,0.8)=9.6+9.6=\mathbf{19.2}$.
  \item \textbf{Total differential.} $f(1,2)=1\cdot8=8$ and
        $\mathrm{d}f=16\,\mathrm{d}x+12\,\mathrm{d}y
        =16(0.1)+12(-0.1)=0.4$, so $f(1.1,1.9)\approx8.4$. Exactly
        $1.1^{2}\cdot1.9^{3}=1.21\cdot6.859=8.299$; the linear estimate is
        high by about $0.1$, the second-order error.
  \item \textbf{Implicit differentiation.} $2x+y+xy'+2yy'=0$, so
        $y'=-\dfrac{2x+y}{x+2y}$. At $(1,2)$ --- which does lie on the
        curve, since $1+2+4=7$ --- this is
        $-\dfrac{4}{5}$.
  \item \textbf{Lagrange (cost minimisation).}
        $2=\lambda y$ and $3=\lambda x$ give $\dfrac{x}{y}=\dfrac32$, so
        $x=1.5y$. Then $1.5y^{2}=54$, $y^{2}=36$, $y=6$ and $x=9$. Minimal
        cost $18+18=36$ --- the two expenditures are equal, as always for
        a product constraint with linear costs.
  \item \textbf{Kuhn--Tucker.} With
        $\mathcal{L}=\ln x+y+\lambda(4-x-y)$: $1/x=\lambda$ and
        $1=\lambda$, so $\lambda=1$ and $x=1$. The budget binds, so $y=3$,
        which satisfies $y\ge0$. Optimum $(1,3)$ with value
        $\ln1+3=3$.
  \item \textbf{Linear programming.} The two constraints meet where
        $x+y=10$ and $2x+y=16$, i.e. $(6,4)$. Vertices $(0,0)$, $(8,0)$,
        $(6,4)$, $(0,10)$ give objective values $0$, $32$, $36$, $30$: the
        maximum is $\mathbf{36}$ at $(6,4)$, where both constraints bind.
  \item \textbf{Discounting.}
        $\dfrac{1000}{1.06^{5}}=\dfrac{1000}{1.33823}\approx747.26$.
  \item \textbf{Geometric series.} First term $a=3$, ratio
        $r=\tfrac12$, so the sum is $\dfrac{3}{1-\tfrac12}=6$.
\end{solutions}

\chapter{Mixed Set C}

Twenty problems, no order, no announced methods. Target time: 75 minutes.

\section{Worked Example}

The last example is an applied one, of the kind that hides three techniques
inside a sentence of prose. As before it is not among the twenty.

\begin{worked}{A machine yields a net flow of $D(t)=8000\,\mathrm{e}^{-0.1t}$ euros per year, discounted continuously at $r=5\%$. Find (a) the present value over ten years, (b) over an infinite horizon, (c) the time at which half the infinite value has accrued}
\wstep{Step 1. Translate the sentence into an integral.}
``Continuous flow'' plus ``continuous discounting'' means
\[
  PV(T)=\int_{0}^{T}D(t)\,\mathrm{e}^{-rt}\,\mathrm{d}t .
\]
\wtrap{Two exponentials appear and they are different things: $\mathrm{e}^{-0.1t}$
is the machine \emph{decaying}, $\mathrm{e}^{-0.05t}$ is money being
\emph{discounted}. Both must be present, and using only one is the standard
misreading of the sentence.}

\wstep{Step 2. Combine the exponents --- this is what makes it easy.}
\[
  PV(T)=\int_{0}^{T}8000\,\mathrm{e}^{-0.1t}\mathrm{e}^{-0.05t}\,\mathrm{d}t
  =8000\int_{0}^{T}\mathrm{e}^{-0.15t}\,\mathrm{d}t .
\]
The two rates \emph{add}, because $\mathrm{e}^{a}\mathrm{e}^{b}
=\mathrm{e}^{a+b}$: the effective rate of decline is $15\%$.

\wstep{Step 3. Integrate.}
\[
  8000\Bigl[\frac{\mathrm{e}^{-0.15t}}{-0.15}\Bigr]_{0}^{T}
  =\frac{8000}{0.15}\bigl(1-\mathrm{e}^{-0.15T}\bigr)
  =53\,333.33\,\bigl(1-\mathrm{e}^{-0.15T}\bigr).
\]
\wtrap{The factor $1/(-0.15)$ turns $\mathrm{e}^{-0.15T}-1$ into
$1-\mathrm{e}^{-0.15T}$. Get this backwards and the present value comes out
negative, which signals an error: a positive flow cannot have a
negative value.}

\wstep{Step 4. Part (a): ten years.}
\[
  PV(10)=53\,333.33\,\bigl(1-\mathrm{e}^{-1.5}\bigr)
        =53\,333.33\,(1-0.22313)
        =\boxed{41\,433\ \text{euros}}
\]

\wstep{Step 5. Part (b): the infinite horizon.}
\[
  PV(\infty)=\lim_{T\to\infty}53\,333.33\bigl(1-\mathrm{e}^{-0.15T}\bigr)
  =\boxed{53\,333\ \text{euros}}
\]
The limit exists because $0.15>0$. Had the machine's output grown faster
than the discount rate --- $D(t)=8000\mathrm{e}^{+0.06t}$ against $r=0.05$
--- the integral would diverge, and the asset would have no finite value.
That convergence condition is the same $g<r$ that governs a growing
perpetuity in \chref{foundations}{2}.

\wstep{Step 6. Part (c): solve for the time, which is a logarithm problem.}
\[
  53\,333.33\bigl(1-\mathrm{e}^{-0.15T}\bigr)=\tfrac12\cdot53\,333.33
  \;\Longrightarrow\;
  1-\mathrm{e}^{-0.15T}=\tfrac12
  \;\Longrightarrow\;
  \mathrm{e}^{-0.15T}=\tfrac12 .
\]
\[
  -0.15T=\ln\tfrac12=-\ln2
  \quad\Longrightarrow\quad
  T=\frac{\ln2}{0.15}=\boxed{4.62\ \text{years}}
\]
The constant $53\,333.33$ cancels before any arithmetic is done --- notice
that before computing it.
\wtrap{This is a half-life, and it depends only on the \emph{combined} rate
$0.15$, not on the size of the flow. Recognising the shape
$\mathrm{e}^{-kT}=\tfrac12\Rightarrow T=\ln2/k$ saves the whole
manipulation.}

\wstep{Step 7. Sanity-check all three answers together.}
$41\,433<53\,333$ \checkmark: ten years cannot be worth more than the infinite horizon. The ten-year
figure is $77.7\%$ of the total, and $4.62$ years already delivers $50\%$
\checkmark --- consistent with a flow that decays. And $53\,333=8000/0.15$
\checkmark, the perpetuity formula with the combined rate.

\wstep{Step 8. What made this an applied problem rather than a calculus exercise.}
Three separate recognitions: that ``continuous flow'' means an integral,
that two exponential rates add, and that ``half the value'' is an equation
to be solved rather than a quantity to be computed. The calculus itself was
one substitution. That ratio --- most of the difficulty in the translation,
little in the technique --- is characteristic of the applied problems in
these sets, and it is why reading the sentence twice is time well spent.
\end{worked}

\section{Problems}

\begin{problems}
  \item Solve $|2x-5|\le3$.
  \item Give the maximal domain of $f(x)=\ln(x^{2}-x-6)$.
  \item Describe $y=-2\sqrt{x+1}+3$ as a chain of transformations, and give
        its domain and range.
  \item Describe the region $\{y\ge x^{2},\ y\le 2x+3\}$ and give the range
        of $x$ it occupies.
  \item Give the level curve of $f(x,y)=xy$ at $c=-4$ and say which
        quadrants it lies in.
  \item Evaluate $\displaystyle\lim_{x\to\infty}
        x\bigl(\sqrt{x^{2}+1}-x\bigr)$.
  \item Evaluate $\displaystyle\lim_{x\to1}\frac{x^{n}-1}{x-1}$ for a
        positive integer $n$.
  \item Carry out a complete discussion of $f(x)=x^{3}-12x$.
  \item For which $t$ does the system $x+y+tz=1$, \;$x+ty+z=1$,
        \;$tx+y+z=1$ fail to have a unique solution?
  \item Give the eigenvalues of
        $\begin{pmatrix}4&2\\2&1\end{pmatrix}$ and its definiteness.
  \item Determine the definiteness of the Hessian of
        $f(x,y)=x^{2}+y^{2}-2xy$ and describe the set of minimisers.
  \item At the point $(3,4)$ on $x^{2}+y^{2}=25$, may $y$ be written
        locally as a function of $x$? Give $\dd{y}{x}$ there.
  \item Show that $f(K,L)=K^{0.3}L^{0.7}$ is homogeneous of degree $1$ and
        verify Euler's theorem $Kf_{K}+Lf_{L}=f$.
  \item Maximise $-(x-2)^{2}-(y-3)^{2}$ subject to $x+y\le10$. Does the
        constraint bind?
  \item For $\max xy$ subject to $px+qy=M$, show $x^{*}=M/(2p)$ and give
        $\partial x^{*}/\partial p$.
  \item Minimise $2x+3y$ subject to $x+y\ge4$, $x\ge1$, $y\ge0$.
  \item Evaluate $\displaystyle\int_{0}^{1}x^{2}\mathrm{e}^{x^{3}}
        \,\mathrm{d}x$.
  \item Evaluate $\displaystyle\int_{1}^{\infty}\frac{1}{x^{2}}
        \,\mathrm{d}x$.
  \item Demand is $p=20-2q$. At the equilibrium quantity $q=5$, compute the
        consumer surplus.
  \item What is an annuity of $500$ per year for ten years worth today at
        $4\%$?
\end{problems}

\section{Solutions}

\begin{solutions}
  \item \textbf{Absolute-value inequality.}
        $-3\le2x-5\le3\Rightarrow2\le2x\le8$, so $1\le x\le4$.
  \item \textbf{Logarithm domain, hence a quadratic inequality.}
        $x^{2}-x-6=(x-3)(x+2)>0$ holds outside the roots:
        $D=(-\infty,-2)\cup(3,\infty)$.
  \item \textbf{Transformations.} From $\sqrt{x}$: shift $1$ left, stretch
        vertically by $2$, reflect in the $x$-axis, shift $3$ up. Domain
        $[-1,\infty)$; since $-2\sqrt{x+1}\in(-\infty,0]$, the range is
        $(-\infty,3]$, with the curve starting at $(-1,3)$ and falling.
  \item \textbf{Region between two curves.} They meet where
        $x^{2}=2x+3$, i.e. $x=-1$ and $x=3$, so the region is the lens
        above the parabola and below the line, occupying
        $-1\le x\le3$. Convex.
  \item \textbf{Level curve.} $xy=-4\Rightarrow y=-4/x$: a hyperbola whose
        branches lie in the \textbf{second and fourth} quadrants, since the
        product is negative.
  \item \textbf{Conjugate, then a limit at infinity.}
        $x\bigl(\sqrt{x^{2}+1}-x\bigr)
        =\dfrac{x}{\sqrt{x^{2}+1}+x}
        =\dfrac{1}{\sqrt{1+1/x^{2}}+1}\to\dfrac12$.
  \item \textbf{L'H\^opital (or the factorisation of $x^{n}-1$).}
        $\dfrac{nx^{n-1}}{1}\to n$.
  \item \textbf{Complete discussion.} Domain $\R$; \emph{odd}, so
        symmetric about the origin; $f=x(x^{2}-12)$ gives roots
        $0,\pm2\sqrt3$; ends $\mp\infty$. $f'=3x^{2}-12=3(x-2)(x+2)$:
        increasing on $(-\infty,-2)$, decreasing on $(-2,2)$, increasing on
        $(2,\infty)$, so a local maximum $(-2,16)$ and a local minimum
        $(2,-16)$. $f''=6x$: concave for $x<0$, convex for $x>0$, with an
        inflection at the origin $(0,0)$.
  \item \textbf{Determinant with a parameter.} Expanding,
        \[
          \det\begin{pmatrix}1&1&t\\1&t&1\\t&1&1\end{pmatrix}
          =(t-1)\bigl[2-t(1+t)\bigr]
          =-(t-1)^{2}(t+2).
        \]
        This vanishes for $t=1$ and $t=-2$. At $t=1$ all three equations
        become $x+y+z=1$: infinitely many solutions. At $t=-2$ the system
        is inconsistent (adding all three gives $0=3$).
  \item \textbf{Eigenvalues.} Trace $5$ and determinant $4-4=0$, so
        $\lambda^{2}-5\lambda=0$ and $\lambda=0$ or $\lambda=5$. Both
        $\ge0$ with one zero: \textbf{positive semidefinite}, not definite.
  \item \textbf{Hessian and the minimiser set.}
        $\Hess f=\begin{pmatrix}2&-2\\-2&2\end{pmatrix}$ with
        $\det=0$ and trace $4>0$: positive semidefinite, so $f$ is convex
        but not strictly so. Indeed $f=(x-y)^{2}\ge0$, and the minimum
        value $0$ is attained on the whole \textbf{line} $y=x$ --- the
        direction in which the Hessian is singular.
  \item \textbf{Implicit function theorem.} With
        $F=x^{2}+y^{2}-25$: $F_{y}=2y=8\neq0$ at $(3,4)$, so yes, a local
        $y(x)$ exists, and
        $\dd{y}{x}=-\dfrac{F_{x}}{F_{y}}=-\dfrac{6}{8}=-\dfrac34$.
  \item \textbf{Homogeneity and Euler.}
        $f(tK,tL)=(tK)^{0.3}(tL)^{0.7}=t^{0.3+0.7}K^{0.3}L^{0.7}=tf(K,L)$,
        so the degree is $1$. And
        $f_{K}=0.3K^{-0.7}L^{0.7}$, $f_{L}=0.7K^{0.3}L^{-0.3}$, so
        \[
          Kf_{K}+Lf_{L}=0.3K^{0.3}L^{0.7}+0.7K^{0.3}L^{0.7}
          =K^{0.3}L^{0.7}=f .
        \]
        Economically: with constant returns to scale, paying each factor
        its marginal product exactly exhausts the product.
  \item \textbf{Check the unconstrained optimum first.} The objective is
        maximal at $(2,3)$, where $x+y=5\le10$: the point is already
        feasible, so the constraint does \textbf{not} bind, its multiplier
        is $0$, and the maximum is $f=0$ at $(2,3)$. Setting up a
        Lagrangian with an equality would have produced the wrong answer.
  \item \textbf{Lagrange with general parameters.} $y=\lambda p$ and
        $x=\lambda q$ give $px=qy$, so the budget splits evenly:
        $px=M/2$, hence $x^{*}=\dfrac{M}{2p}$ and $y^{*}=\dfrac{M}{2q}$.
        Then $\dfrac{\partial x^{*}}{\partial p}=-\dfrac{M}{2p^{2}}<0$:
        the own-price effect is negative, and the cross-price effect
        $\partial x^{*}/\partial q=0$ --- Cobb--Douglas demands do not
        respond to the other price.
  \item \textbf{Linear programming, minimisation.} The vertices of the
        feasible set are $(4,0)$ (where $x+y=4$ meets $y=0$) and $(1,3)$
        (where $x=1$ meets $x+y=4$). The objective takes the values
        $8$ and $11$, so the minimum is $\mathbf{8}$ at $(4,0)$. The region
        is unbounded, but the objective increases in both variables, so a
        minimum exists.
  \item \textbf{Substitution.} With $u=x^{3}$, $\mathrm{d}u=3x^{2}
        \mathrm{d}x$:
        $\dfrac13\displaystyle\int_{0}^{1}\mathrm{e}^{u}\mathrm{d}u
        =\dfrac{\mathrm{e}-1}{3}\approx0.573$.
  \item \textbf{Improper integral.}
        $\displaystyle\int_{1}^{b}x^{-2}\mathrm{d}x
        =1-\dfrac1b\to1$ as $b\to\infty$: the integral \textbf{converges}
        to $1$.
  \item \textbf{Consumer surplus.} At $q=5$ the price is
        $p=20-10=10$, and
        \[
          CS=\int_{0}^{5}(20-2q)\,\mathrm{d}q-10\cdot5
            =\bigl[20q-q^{2}\bigr]_{0}^{5}-50
            =75-50=25 .
        \]
  \item \textbf{Annuity present value.}
        $500\cdot\dfrac{1-1.04^{-10}}{0.04}
        =500\cdot8.1109\approx4055.4$.
\end{solutions}

\chapter{Mixed Set D: Full Investigations}

The first three sets drill one skill per problem. This set does the
opposite. Each problem hands you a single function of two variables and
runs it through the whole chain of questions that a function of two
variables invites: where it lives, what values it takes, where it can be
differentiated, what its level curves look like, where its gradient
points, what its Hessian says, whether it has critical points, which of
its contour sets are convex, and finally a checklist of properties, each
to be decided \emph{with a reason}.

None of the steps is new. What is new is that they feed each other: the
domain decides where the derivative may be taken, the Hessian decides
concavity, concavity decides quasiconcavity, the domain and the range
decide whether a maximum exists. A slip early in the chain propagates,
and a good investigation catches it, because the later answers stop
fitting together. Target time: twenty to twenty-five minutes per
function. The solutions tag every step with the volume and chapter whose
rule it uses, so that a weak link can be traced back and repaired.

\section{The Investigation Chain}

\begin{rulebox}{The chain, in order}
\begin{enumerate}[label=\arabic*.,leftmargin=1.6em,itemsep=0.3ex]
  \item \textbf{Domain.} Collect the obstructions (roots need a
        nonnegative radicand, logarithms a positive argument,
        denominators must not vanish), sketch the set, and say whether it
        is closed, bounded and convex. (Vol.~2, Ch.~1 and Ch.~4;
        Vol.~5, Ch.~10; Vol.~7, Ch.~2.)
  \item \textbf{Range.} Find the smallest and largest values, or show that
        there are none, and use continuity on a connected domain to fill
        in everything between. (Vol.~2, Ch.~1; Vol.~5, Ch.~9.)
  \item \textbf{Derivative.} Compute $f_{x}$, $f_{y}$ and say where they
        exist and are continuous: on that open set $f$ is differentiable
        (continuous partials are sufficient, not necessary). Where a
        partial fails to exist, $f$ is not differentiable. Roots are the
        usual offenders on the boundary. (Vol.~5, Ch.~10.)
  \item \textbf{Level curves.} Solve $f(x,y)=c$ for one variable and
        sketch several values of $c$. (Vol.~5, Ch.~9.)
  \item \textbf{Gradient at a point.} Evaluate $\grad f$, draw it, check
        that it is orthogonal to the level curve through the point, and
        read it: it points in the direction of steepest ascent, and its
        norm $\lVert\grad f\rVert$ is the slope in that direction.
        (Vol.~5, Ch.~10 and Ch.~13.)
  \item \textbf{Hessian.} Second partials, Young's theorem, and the
        definiteness by Sylvester's criterion. (Vol.~5, Ch.~11.)
  \item \textbf{Critical points.} Solve $\grad f=\mathbf{0}$ \emph{inside
        the domain}, or explain why there is no solution; classify with the
        Hessian. (Vol.~7, Ch.~5.)
  \item \textbf{Contour sets.} Mark an upper set $\{f\ge c\}$ and a lower
        set $\{f\le c\}$ and decide whether each is convex.
        (Vol.~5, Ch.~9; Vol.~7, Ch.~2.)
  \item \textbf{Checklist.} Concave, convex, quasiconcave, quasiconvex,
        has a maximum, has a minimum, increasing, decreasing, epigraph
        convex: each answered with a proof or a counterexample.
        (Vol.~7, Ch.~1 and Ch.~3.)
\end{enumerate}
\end{rulebox}

\begin{rulebox}{Three definitions the checklist needs}
\begin{itemize}[leftmargin=1.6em,itemsep=0.3ex]
  \item \textbf{Monotone in several variables.} Order points componentwise:
        $(x,y)\le(x',y')$ means $x\le x'$ and $y\le y'$. Then $f$ is
        \emph{increasing} \de{monoton wachsend} if
        $(x,y)\le(x',y')\Rightarrow f(x,y)\le f(x',y')$, and
        \emph{decreasing} if the inequality between the values is reversed.
        On a convex domain with interior points, a continuous $f$ with
        $f_{x}\ge0$ and $f_{y}\ge0$ at every interior point is increasing.
        To refute it, exhibit two ordered points whose values are in the
        wrong order.
  \item \textbf{Epigraph and hypograph.} The \emph{epigraph}
        \de{Epigraph} of $f$ is the set of points on or above its graph,
        $\{(x,y,z): (x,y)\in D,\ z\ge f(x,y)\}$; the \emph{hypograph} is
        the set on or below it. On a convex domain, the epigraph is a convex
        set exactly when $f$ is convex, and the hypograph exactly when $f$
        is concave. So the epigraph question is the convexity question
        asked again, and the two answers must agree.
  \item \textbf{Existence of extrema.} A continuous function on a closed
        and bounded domain attains a maximum and a minimum (Weierstrass,
        Vol.~7, Ch.~1). On an unbounded or open domain, look at the range
        first: if $f$ grows without bound there is no maximum, and if the
        infimum is approached but never reached there is no minimum.
        Extrema on the boundary of the domain are not found by
        $\grad f=\mathbf{0}$, and are often points where $f$ is not even
        differentiable.
\end{itemize}
\end{rulebox}

\begin{rulebox}{Shortcuts that keep the checklist consistent}
\begin{itemize}[leftmargin=1.6em,itemsep=0.3ex]
  \item Concave $\Rightarrow$ quasiconcave; convex $\Rightarrow$
        quasiconvex. A function that is both concave and convex is affine.
  \item A strictly increasing transform $F$ ($\ln$ of a positive $f$,
        $\sqrt{\cdot}$ or a positive power of a nonnegative $f$) only
        relabels the contour sets, $\{F(f)\ge F(c)\}=\{f\ge c\}$, so it
        preserves quasiconcavity and quasiconvexity but not necessarily
        concavity.
  \item If $g$ is concave and $F$ is increasing and concave, $F\circ g$ is
        concave: this is how $\sqrt{\text{concave}}$ and
        $\ln(\text{concave})$ are handled without a Hessian.
  \item One straight line in the domain along which $f$ is strictly convex
        refutes concavity; one along which it is strictly concave refutes
        convexity. Two points in a contour set whose midpoint is not in it
        refute quasiconcavity or quasiconvexity.
  \item Cobb--Douglas $x^{a}y^{b}$ ($a,b>0$) on $\R^{2}_{+}$: always
        quasiconcave; concave iff $a+b\le1$; never convex.
        (Vol.~7, Ch.~3.)
\end{itemize}
\end{rulebox}

\section{Worked Example}

\begin{worked}{Investigate $f(x,y)=\sqrt{8-2x^{2}-y^{2}}$ completely: domain, range, derivative, level curves, gradient at $(1,2)$, Hessian, critical points, contour sets, and the full property checklist}
\wstep{Step 1. Domain (Vol.~2, Ch.~1; Vol.~2, Ch.~4).}
The only obstruction is the root: $8-2x^{2}-y^{2}\ge0$, that is
\[
  D=\{(x,y): 2x^{2}+y^{2}\le8\}
   =\Bigl\{\tfrac{x^{2}}{4}+\tfrac{y^{2}}{8}\le1\Bigr\},
\]
the filled ellipse with semi-axes $2$ along $x$ and $2\sqrt2$ along $y$,
boundary included. It is \textbf{closed} (the inequality is weak),
\textbf{bounded}, and \textbf{convex}: it is a lower contour set of the
convex function $2x^{2}+y^{2}$ (Vol.~7, Ch.~2).
\wtrap{Writing $2x^{2}+y^{2}<8$ loses the boundary. The root of zero is
zero, so the boundary ellipse belongs to the domain, and it will carry the
minimum.}

\wstep{Step 2. Range (Vol.~2, Ch.~1).}
The radicand is largest at the origin, where it equals $8$, and smallest
($0$) on the boundary. Since $\sqrt{\cdot}$ is increasing,
$0\le f\le\sqrt8=2\sqrt2$. $f$ is continuous on the connected set $D$, so
every value in between is taken:
\[
  f(D)=[\,0,\,2\sqrt2\,].
\]

\wstep{Step 3. Derivative, and where it exists (Vol.~5, Ch.~10).}
By the chain rule,
\[
  f_{x}=\frac{-4x}{2\sqrt{8-2x^{2}-y^{2}}}=-\frac{2x}{f},\qquad
  f_{y}=\frac{-2y}{2\sqrt{8-2x^{2}-y^{2}}}=-\frac{y}{f}.
\]
Both are continuous wherever $f>0$, so $f$ is differentiable on the
\textbf{interior} $2x^{2}+y^{2}<8$. On the boundary the formulas divide by
zero, and the difference quotient confirms the failure: at $(2,0)$,
$\bigl(f(2-h,0)-f(2,0)\bigr)/h=\sqrt{8h-2h^{2}}/h\to\infty$ as
$h\to0^{+}$. The surface meets the plane $z=0$ vertically, like the edge
of a dome.

\wstep{Step 4. Level curves (Vol.~5, Ch.~9).}
For $0\le c\le2\sqrt2$, $f=c$ means $2x^{2}+y^{2}=8-c^{2}$: concentric
ellipses with the same shape as the boundary, shrinking to the origin as
$c\to2\sqrt2$. The values $c=0,1,\sqrt2,2,\sqrt7$ give
$2x^{2}+y^{2}=8,7,6,4,1$ (left panel below). The level $c=0$ is the
boundary of the domain; the level $c=2\sqrt2$ is the single point
$(0,0)$.

\wstep{Step 5. Gradient at $(1,2)$ (Vol.~5, Ch.~10).}
Here $f(1,2)=\sqrt{8-2-4}=\sqrt2$, so
\[
  \grad f(1,2)=\begin{pmatrix}-2/\sqrt2\\-2/\sqrt2\end{pmatrix}
  =\begin{pmatrix}-\sqrt2\\-\sqrt2\end{pmatrix},
  \qquad
  \lVert\grad f(1,2)\rVert=\sqrt{2+2}=2 .
\]
\emph{Orthogonality.} The level curve through $(1,2)$ is
$2x^{2}+y^{2}=6$; implicit differentiation gives
$y'=-4x/(2y)=-1$ there, so a tangent direction is $(1,-1)$, and
$(-\sqrt2)(1)+(-\sqrt2)(-1)=0$ \checkmark.
\emph{Direction.} The gradient points into the ellipse, towards higher
levels: the direction of steepest ascent. \emph{Norm.} $2$ is the slope of
the surface in that direction, the largest directional derivative at
$(1,2)$ (Vol.~5, Ch.~13). In every direction along the level curve the
slope is $0$.

\begin{gallery}
\begin{tikzpicture}
  \begin{axis}[contouraxis, xmin=-3.2, xmax=3.2, ymin=-3.2, ymax=3.2,
      xtick={-2,-1,1,2}, ytick={-2,-1,1,2},
      title={domain, levels, $\grad f(1,2)$}]
    \addplot[name path=top, draw=none, domain=-2:2]
      {sqrt(max(8-2*x^2,0))};
    \addplot[name path=bot, draw=none, domain=-2:2]
      {-sqrt(max(8-2*x^2,0))};
    \region{top}{bot}{}
    \addplot[accent, line width=1pt, domain=0:360, samples=121]
      ({2*cos(x)},{sqrt(8)*sin(x)});
    \addplot[domain=0:360, samples=121]({sqrt(3.5)*cos(x)},{sqrt(7)*sin(x)});
    \addplot[domain=0:360, samples=121]({sqrt(3)*cos(x)},{sqrt(6)*sin(x)});
    \addplot[domain=0:360, samples=121]({sqrt(2)*cos(x)},{2*sin(x)});
    \addplot[domain=0:360, samples=121]({sqrt(0.5)*cos(x)},{sin(x)});
    \draw[dashed, gray!70!black, line width=0.5pt]
      (axis cs:0.1,2.9) -- (axis cs:1.9,1.1);
    \draw[-{Stealth[length=6pt,width=4pt]}, accent, line width=1pt]
      (axis cs:1,2) -- (axis cs:-0.414,0.586);
    \addplot[only marks, mark=*, mark size=1.2pt, black, forget plot]
      coordinates {(1,2) (0,0)};
  \end{axis}
\end{tikzpicture}\hfill
\begin{tikzpicture}
  \begin{axis}[contouraxis, xmin=-3.2, xmax=3.2, ymin=-3.2, ymax=3.2,
      xtick={-2,-1,1,2}, ytick={-2,-1,1,2},
      title={upper set $\{f\ge2\}$ (shaded)}]
    \addplot[name path=itop, draw=none, domain=-1.4142:1.4142]
      {sqrt(max(4-2*x^2,0))};
    \addplot[name path=ibot, draw=none, domain=-1.4142:1.4142]
      {-sqrt(max(4-2*x^2,0))};
    \region{itop}{ibot}{}
    \addplot[accent, line width=1pt, domain=0:360, samples=121]
      ({2*cos(x)},{sqrt(8)*sin(x)});
    \addplot[domain=0:360, samples=121]({sqrt(2)*cos(x)},{2*sin(x)});
    \draw[gray!70!black, line width=0.5pt] (axis cs:0,2) -- (axis cs:0,-2);
    \addplot[only marks, mark=*, mark size=1.2pt, black, forget plot]
      coordinates {(0,2) (0,-2)};
    \addplot[only marks, mark=o, mark size=1.6pt, black, forget plot]
      coordinates {(0,0)};
  \end{axis}
\end{tikzpicture}
\gcap{Left: the domain (shaded), its boundary $f=0$ (bold), the levels
$f=1,\sqrt2,2,\sqrt7$, and the gradient at $(1,2)$, orthogonal to the
dashed tangent of its level curve. Right: the upper set $\{f\ge2\}$ is a
filled ellipse; the lower set $\{f\le2\}$ is the ring between the two
curves, which contains $(0,\pm2)$ but not their midpoint.}
\end{gallery}

\wstep{Step 6. Hessian (Vol.~5, Ch.~11).}
Differentiate $f_{x}=-2xf^{-1}$ and $f_{y}=-yf^{-1}$ once more, using
$f_{x}$, $f_{y}$ from Step 3:
\[
  f_{xx}=-\frac{2}{f}+\frac{2x f_{x}}{f^{2}}
        =-\frac{2f^{2}+4x^{2}}{f^{3}}
        =-\frac{2(8-y^{2})}{f^{3}},\qquad
  f_{yy}=-\frac{1}{f}+\frac{y f_{y}}{f^{2}}
        =-\frac{8-2x^{2}}{f^{3}},
\]
\[
  f_{xy}=\frac{2x f_{y}}{f^{2}}=-\frac{2xy}{f^{3}}
  \quad(\text{equal to } f_{yx}\text{ by Young}).
\]
The determinant simplifies remarkably:
\[
  \det\Hess f=\frac{2(8-y^{2})(8-2x^{2})-4x^{2}y^{2}}{f^{6}}
  =\frac{16(8-2x^{2}-y^{2})}{f^{6}}=\frac{16}{f^{4}}>0 .
\]
In the interior $y^{2}<8$, so $f_{xx}<0$: the Hessian is
\textbf{negative definite} at every interior point. At $(1,2)$, with
$f^{3}=2\sqrt2$:
\[
  \Hess f(1,2)=\begin{pmatrix}-2\sqrt2&-\sqrt2\\[2pt]-\sqrt2&-\tfrac32\sqrt2
  \end{pmatrix},\qquad \det=6-2=4=\frac{16}{(\sqrt2)^{4}}\ \checkmark
\]
\wtrap{Do not expand $\det\Hess f$ with the roots still in it. Writing
every entry over the common denominator $f^{3}$ and replacing $f^{2}$ by
$8-2x^{2}-y^{2}$ is what makes the sign visible.}

\wstep{Step 7. Critical points (Vol.~7, Ch.~5).}
$f_{x}=0\iff x=0$ and $f_{y}=0\iff y=0$, so the only candidate is
$(0,0)$, which lies in the interior. There
$\Hess f(0,0)=\begin{pmatrix}-1/\sqrt2&0\\0&-1/(2\sqrt2)\end{pmatrix}$ is
negative definite: a strict local maximum, with value $2\sqrt2$. The
boundary points are not critical points; $f$ is not differentiable there.

\wstep{Step 8. Contour sets (Vol.~5, Ch.~9; Vol.~7, Ch.~2).}
Take $c=2$. The upper set $\{f\ge2\}=\{2x^{2}+y^{2}\le4\}$ is a filled
ellipse: \textbf{convex}. The lower set
$\{f\le2\}=\{4\le2x^{2}+y^{2}\le8\}$ is an elliptical ring: it contains
$(0,2)$ and $(0,-2)$ but not their midpoint $(0,0)$, so it is
\textbf{not convex} (right panel).

\wstep{Step 9. The checklist, each with its reason.}
\begin{itemize}[leftmargin=1.6em,itemsep=0.3ex]
  \item \emph{Concave: yes.} The radicand $8-2x^{2}-y^{2}$ is concave and
        $\sqrt{\cdot}$ is increasing and concave, so the composition is
        concave on $D$ (Vol.~7, Ch.~3). Step 6 confirms it, even strictly,
        on the interior.
  \item \emph{Convex: no.} The Hessian is negative definite; concretely,
        $f(0,0)=2\sqrt2$ lies above the chord value
        $\tfrac12\bigl(f(-2,0)+f(2,0)\bigr)=0$.
  \item \emph{Quasiconcave: yes}, because concave; every upper set is a
        filled ellipse (Step 8).
  \item \emph{Quasiconvex: no.} The lower set $\{f\le2\}$ is not convex
        (Step 8).
  \item \emph{Maximum: yes}, $2\sqrt2$ at $(0,0)$, unique; global because
        $f$ is concave and $(0,0)$ is a critical point (Vol.~7, Ch.~5).
  \item \emph{Minimum: yes}, $0$, attained at \emph{every} point of the
        boundary ellipse. Weierstrass guarantees it (continuous on a closed
        bounded set); the first-order condition could never find it, since
        $f$ is not differentiable there.
  \item \emph{Increasing: no.} $(0,0)\le(1,0)$ but
        $f(1,0)=\sqrt6<2\sqrt2=f(0,0)$.
  \item \emph{Decreasing: no.} $(-1,0)\le(0,0)$ but
        $f(-1,0)=\sqrt6<f(0,0)$.
  \item \emph{Epigraph convex: no}, since $f$ is not convex. Its
        hypograph, the solid dome, is convex, matching concavity.
\end{itemize}

\wstep{Step 10. Do the answers fit together?}
Concave with an interior critical point gives a global maximum
\checkmark. Not quasiconvex, yet concave: a dome has convex upper sets and
ring-shaped lower sets \checkmark. The range $[0,2\sqrt2]$ matches the
maximum and minimum found \checkmark. Neither monotone, as one expects of
a function that peaks in the middle of its domain \checkmark.
\end{worked}

\section{Problems}

Each problem is one investigation. Work every part, and give a reason for
every item of the checklist.

\begin{problems}
  \item Investigate $f(x,y)=\sqrt{6-x-y^{2}}$.
        \begin{enumerate}[label=(\alph*),leftmargin=1.8em,itemsep=0.2ex]
          \item Give and sketch the domain. Is it closed, bounded, convex?
                What values does $f$ take?
          \item Compute $f_{x}$ and $f_{y}$. Is $f$ differentiable on its
                whole domain?
          \item Describe and sketch the level curves.
          \item Compute $\grad f(1,1)$, draw it, show that it is orthogonal
                to the level curve through $(1,1)$, and state what its
                direction and its norm mean.
          \item Compute the Hessian and its definiteness.
          \item Does $f$ have critical points?
          \item Mark an upper and a lower contour set at level $2$. Which is
                convex?
          \item Decide, with reasons: concave, convex, quasiconcave,
                quasiconvex, has a maximum, has a minimum, increasing,
                decreasing, epigraph convex.
        \end{enumerate}
  \item Investigate the utility function
        $u:\R^{2}_{++}\to\R$, $u(x,y)=\ln\bigl(x^{1/3}y^{1/2}\bigr)$, where
        $\R^{2}_{++}=\{(x,y):x>0,\ y>0\}$.
        \begin{enumerate}[label=(\alph*),leftmargin=1.8em,itemsep=0.2ex]
          \item Rewrite $u$ as a sum. Describe the domain (closed? bounded?
                convex?) and the range.
          \item Describe and sketch the level curves. How do they compare
                with those of $x^{1/3}y^{1/2}$?
          \item Compute $\grad u(1,1)$, draw it, check orthogonality to the
                level curve, and interpret direction and norm.
          \item Compute the Hessian. Show that both marginal utilities are
                diminishing, and give the definiteness.
          \item Does $u$ have critical points?
          \item Mark an upper and a lower contour set at level $0$. Which is
                convex?
          \item Decide, with reasons: concave, convex, quasiconcave,
                quasiconvex, has a maximum, has a minimum, increasing,
                decreasing, epigraph convex.
        \end{enumerate}
  \item Investigate the production function $f:\R^{2}_{+}\to\R$,
        $f(x,y)=x^{3/4}y^{1/2}$.
        \begin{enumerate}[label=(\alph*),leftmargin=1.8em,itemsep=0.2ex]
          \item Name this type of function and its returns to scale. Give
                the domain (closed? convex?) and the range.
          \item Compute the derivative and say where it exists. Show that
                both marginal products are diminishing.
          \item Compute the Hessian and its definiteness. Does diminishing
                marginal productivity make $f$ concave?
          \item Sketch the level curves, and draw $\grad f(1,1)$; check
                orthogonality and interpret direction and norm.
          \item Does $f$ have critical points? Where does its minimum come
                from?
          \item Mark an upper and a lower contour set at level $1$. Which is
                convex?
          \item Decide, with reasons: concave, convex, quasiconcave,
                quasiconvex, has a maximum, has a minimum, increasing,
                decreasing, epigraph convex.
        \end{enumerate}
  \item Investigate $f:\R^{2}_{+}\to\R$, $f(x,y)=x^{1/3}y^{2/3}$.
        \begin{enumerate}[label=(\alph*),leftmargin=1.8em,itemsep=0.2ex]
          \item Name the function and its returns to scale; give domain and
                range.
          \item Compute the derivative and the Hessian, and the
                definiteness. Which direction does the Hessian annihilate?
          \item Sketch the level curves and draw $\grad f(1,8)$; check
                orthogonality, interpret direction and norm, and verify
                Euler's theorem at $(1,8)$.
          \item Does $f$ have critical points?
          \item Mark an upper and a lower contour set at level $4$. Which is
                convex?
          \item Decide, with reasons: concave, strictly concave, convex,
                quasiconcave, quasiconvex, has a maximum, has a minimum,
                increasing, decreasing, epigraph convex.
        \end{enumerate}
  \item Investigate $f(x,y)=\ln\bigl(4+x^{2}+y^{2}\bigr)$.
        \begin{enumerate}[label=(\alph*),leftmargin=1.8em,itemsep=0.2ex]
          \item Give the domain and the range.
          \item Compute the gradient; where is $f$ differentiable?
          \item Describe the level curves, compute $\grad f(1,1)$, and check
                that it is orthogonal to the level curve through $(1,1)$.
          \item Compute the Hessian, its determinant in general, and its
                definiteness at $(0,0)$, $(1,1)$ and $(2,2)$.
          \item Find and classify the critical points.
          \item Mark the upper and lower contour sets at level $\ln8$.
                Which is convex?
          \item Decide, with reasons: concave, convex, quasiconcave,
                quasiconvex, has a maximum, has a minimum, increasing,
                decreasing, epigraph convex. On which disk centred at the
                origin is $f$ convex?
        \end{enumerate}
\end{problems}

\section{Solutions}

\begin{solutions}
  \item \textbf{Root of a concave expression on an unbounded closed domain.}
        \begin{enumerate}[label=(\alph*),leftmargin=1.8em,itemsep=0.4ex]
          \item \emph{Domain (Vol.~2, Ch.~1 and Ch.~4).} Radicand
                $\ge0$: $D=\{x\le6-y^{2}\}$, the region on and to the left of
                the parabola $x=6-y^{2}$ (vertex $(6,0)$, opening to the
                left). Closed (weak inequality), unbounded (it contains the
                whole ray $y=0$, $x\le6$), and convex: it is
                the upper set $\{g\ge0\}$ of the concave function
                $g=6-x-y^{2}$ (Vol.~7, Ch.~2). \emph{Range
                (Vol.~2, Ch.~1):} $f=0$ on the parabola and
                $f(x,0)=\sqrt{6-x}\to\infty$ as $x\to-\infty$; by
                continuity $f(D)=[0,\infty)$.
          \item \emph{Derivative (Vol.~5, Ch.~10).}
                $f_{x}=\dfrac{-1}{2\sqrt{6-x-y^{2}}}=-\dfrac{1}{2f}$,
                \;$f_{y}=\dfrac{-2y}{2f}=-\dfrac{y}{f}$, continuous where
                $f>0$. So $f$ is differentiable on the interior
                $x<6-y^{2}$ only. On the parabola the partials blow up:
                at $(6,0)$, $\bigl(f(6-h,0)-0\bigr)/h=\sqrt h/h\to\infty$.
          \item \emph{Level curves (Vol.~5, Ch.~9).} $f=c\ge0$ gives
                $x=6-c^{2}-y^{2}$: copies of the boundary parabola, shifted
                left by $c^{2}$. For $c=0,1,2,3$ the vertices sit at
                $x=6,5,2,-3$; higher levels lie further left.
          \item \emph{Gradient (Vol.~5, Ch.~10).} $f(1,1)=\sqrt4=2$, so
                \[
                  \grad f(1,1)=\begin{pmatrix}-1/4\\-1/2\end{pmatrix},\qquad
                  \lVert\grad f(1,1)\rVert=\sqrt{\tfrac1{16}+\tfrac14}
                  =\frac{\sqrt5}{4}\approx0.56 .
                \]
                The level curve through $(1,1)$ is $x=2-y^{2}$; there
                $\mathrm{d}x/\mathrm{d}y=-2y=-2$, so
                $\mathrm{d}y/\mathrm{d}x=-\tfrac12$ and $(2,-1)$ is a tangent
                direction. $(-\tfrac14)(2)+(-\tfrac12)(-1)=0$ \checkmark.
                The gradient points left and down, towards the inside of the
                parabolas where the levels rise: steepest ascent, with slope
                $\sqrt5/4$ in that direction (Vol.~5, Ch.~13).
          \item \emph{Hessian (Vol.~5, Ch.~11).} From $f_{x}=-\tfrac12f^{-1}$
                and $f_{y}=-yf^{-1}$:
                \[
                  f_{xx}=\tfrac12f^{-2}f_{x}=-\frac{1}{4f^{3}},\qquad
                  f_{xy}=\tfrac12f^{-2}f_{y}=-\frac{y}{2f^{3}},
                \]
                \[
                  f_{yy}=-\frac1f+\frac{yf_{y}}{f^{2}}
                        =-\frac{f^{2}+y^{2}}{f^{3}}=-\frac{6-x}{f^{3}} .
                \]
                $\det\Hess f=\dfrac{(6-x)-y^{2}}{4f^{6}}
                =\dfrac{f^{2}}{4f^{6}}=\dfrac{1}{4f^{4}}>0$ and
                $f_{xx}<0$: \textbf{negative definite} on the interior. At
                $(1,1)$: $\Hess f=\begin{pmatrix}-1/32&-1/16\\-1/16&-5/8
                \end{pmatrix}$, $\det=\tfrac{5}{256}-\tfrac{1}{256}
                =\tfrac1{64}=\tfrac{1}{4\cdot2^{4}}$ \checkmark.
          \item \emph{Critical points (Vol.~7, Ch.~5).} None:
                $f_{x}=-1/(2f)$ is never $0$. Moving left always raises $f$.
          \item \emph{Contour sets (Vol.~5, Ch.~9; Vol.~7, Ch.~2).}
                $\{f\ge2\}=\{x\le2-y^{2}\}$, the inside of a parabola:
                \textbf{convex}. $\{f\le2\}=\{2-y^{2}\le x\le6-y^{2}\}$, a
                curved band between two parabolas: it contains $(-2,2)$ and
                $(-2,-2)$ (check: $-2\le-2\le2$) but not their midpoint
                $(-2,0)$ (since $2\le-2$ fails). \textbf{Not convex.}
          \item \emph{Checklist (Vol.~7, Ch.~1 and Ch.~3).}
                \emph{Concave: yes}, $\sqrt{\cdot}$ increasing and concave
                applied to the concave $6-x-y^{2}$ (and (e) confirms it).
                \emph{Convex: no}, negative definite Hessian.
                \emph{Quasiconcave: yes}, because concave.
                \emph{Quasiconvex: no}, by the band in (g).
                \emph{Maximum: no}, $f$ is unbounded above.
                \emph{Minimum: yes}, $0$ on the whole parabola $x=6-y^{2}$,
                boundary points where $f$ is not differentiable.
                \emph{Increasing: no}, $f$ falls in $x$: $(0,0)\le(1,0)$
                with $f(1,0)=\sqrt5<\sqrt6=f(0,0)$.
                \emph{Decreasing: no}, $(0,-1)\le(0,0)$ with
                $f(0,-1)=\sqrt5<\sqrt6=f(0,0)$.
                \emph{Epigraph convex: no}, since $f$ is not convex (the
                hypograph is).
        \end{enumerate}
  \item \textbf{Log-Cobb--Douglas.}
        \begin{enumerate}[label=(\alph*),leftmargin=1.8em,itemsep=0.4ex]
          \item \emph{Logarithm rules (Vol.~1, Ch.~2).}
                $u=\tfrac13\ln x+\tfrac12\ln y$. \emph{Domain
                (Vol.~5, Ch.~10):} the open quadrant $\R^{2}_{++}$, open
                (not closed), unbounded, convex. \emph{Range:}
                $u(t,1)=\tfrac13\ln t$ runs through all of $\R$ as $t$
                runs through $(0,\infty)$, so $u(\R^{2}_{++})=\R$.
          \item \emph{Level curves (Vol.~5, Ch.~9).}
                $\tfrac13\ln x+\tfrac12\ln y=c\iff\ln y=2c-\tfrac23\ln x
                \iff y=\mathrm{e}^{2c}x^{-2/3}$: decreasing, convex curves
                asymptotic to both axes, moving outwards as $c$ grows. They
                are exactly the level curves of $x^{1/3}y^{1/2}$, relabelled
                ($u=\ln v$ with $v$ the Cobb--Douglas): an increasing
                transform changes the labels, not the curves.
          \item \emph{Gradient (Vol.~5, Ch.~10).} $u_{x}=\dfrac{1}{3x}$,
                $u_{y}=\dfrac{1}{2y}$, so $\grad u(1,1)=(\tfrac13,\tfrac12)$
                with norm $\sqrt{\tfrac19+\tfrac14}=\tfrac{\sqrt{13}}{6}
                \approx0.60$. The level curve through $(1,1)$ is $u=0$, i.e.
                $y=x^{-2/3}$, with slope $-\tfrac23$ at $x=1$ and tangent
                direction $(3,-2)$; $\tfrac13\cdot3-\tfrac12\cdot2=0$
                \checkmark. The gradient points up and to the right, the
                direction in which utility rises fastest, at rate
                $\sqrt{13}/6$.
          \item \emph{Hessian (Vol.~5, Ch.~11).}
                $\Hess u=\begin{pmatrix}-\frac{1}{3x^{2}}&0\\0&-\frac{1}{2y^{2}}
                \end{pmatrix}$. $u_{xx}<0$ and $u_{yy}<0$: each marginal
                utility $u_{x}=1/(3x)$, $u_{y}=1/(2y)$ falls as its own good
                increases. Diagonal with negative entries: \textbf{negative
                definite} everywhere.
          \item \emph{Critical points (Vol.~7, Ch.~5).} None: $u_{x}=1/(3x)>0$.
          \item \emph{Contour sets.} $u\le0\iff x^{1/3}y^{1/2}\le1\iff
                x^{2}y^{3}\le1$. The upper set $\{u\ge0\}=\{y\ge x^{-2/3}\}$
                lies above a convex curve: \textbf{convex}. The lower set is
                not: $(8,\tfrac14)$ and $(\tfrac18,4)$ both satisfy
                $x^{2}y^{3}=1$, but at the midpoint
                $(\tfrac{65}{16},\tfrac{17}{8})$ the product is far above
                $1$ (already $x^{2}>16$ and $y^{3}>8$).
          \item \emph{Checklist (Vol.~7, Ch.~3).}
                \emph{Concave: yes}, strictly, by the negative definite
                Hessian (or as a positive sum of the concave $\ln x$,
                $\ln y$).
                \emph{Convex: no}, the Hessian is negative definite.
                \emph{Quasiconcave: yes}, because concave.
                \emph{Quasiconvex: no}, by (f).
                \emph{Maximum: no}, $u(t,1)\to\infty$ as $t\to\infty$.
                \emph{Minimum: no}, $u(t,1)\to-\infty$ as $t\to0^{+}$; the
                open domain has no boundary to stop at.
                \emph{Increasing: yes}, strictly: $u_{x}>0$, $u_{y}>0$ on a
                convex domain.
                \emph{Decreasing: no}, as it is strictly increasing.
                \emph{Epigraph convex: no}, $u$ is not convex.
        \end{enumerate}
  \item \textbf{Cobb--Douglas with exponent sum above $1$.}
        \begin{enumerate}[label=(\alph*),leftmargin=1.8em,itemsep=0.4ex]
          \item \emph{Name (Vol.~5, Ch.~14).} A Cobb--Douglas function,
                homogeneous of degree $\tfrac34+\tfrac12=\tfrac54>1$:
                $f(tx,ty)=t^{5/4}f(x,y)$, \textbf{increasing returns to
                scale}. Domain $\R^{2}_{+}$: closed, unbounded, convex.
                Range $[0,\infty)$: $f=0$ on the axes and $f(t,t)=t^{5/4}$
                is unbounded.
          \item \emph{Derivative (Vol.~5, Ch.~10).}
                \[
                  f_{x}=\tfrac34x^{-1/4}y^{1/2},\qquad
                  f_{y}=\tfrac12x^{3/4}y^{-1/2},
                \]
                continuous on the open quadrant $x,y>0$, so $f$ is
                differentiable there. On the axis $x=0$ (with $y>0$) the
                factor $x^{-1/4}$ is infinite, and likewise $y^{-1/2}$ on the
                axis $y=0$ (with $x>0$): the partials do not exist there.
                The origin is the exception: $f$ vanishes on both axes, so
                $f_{x}(0,0)=f_{y}(0,0)=0$, and
                $0\le f(h,k)\le\lVert(h,k)\rVert^{5/4}$ shows that $f$ is
                differentiable there (from inside the quadrant) with zero
                derivative. Second partials:
                $f_{xx}=-\tfrac{3}{16}x^{-5/4}y^{1/2}<0$,
                $f_{yy}=-\tfrac14x^{3/4}y^{-3/2}<0$: both marginal products
                \textbf{diminish}.
          \item \emph{Hessian (Vol.~5, Ch.~11).}
                $f_{xy}=\tfrac38x^{-1/4}y^{-1/2}$, and
                \[
                  \det\Hess f=\tfrac{3}{64}x^{-1/2}y^{-1}
                  -\tfrac{9}{64}x^{-1/2}y^{-1}
                  =-\tfrac{3}{32}\,x^{-1/2}y^{-1}<0 :
                \]
                \textbf{indefinite} at every interior point, so $f$ is
                neither concave nor convex. Diminishing marginal products
                (negative diagonal) are \emph{not} enough for concavity;
                the cross effect $f_{xy}$ is too strong. Along the diagonal
                the function is $f(t,t)=t^{5/4}$, strictly convex in $t$,
                which refutes concavity directly.
          \item \emph{Level curves and gradient (Vol.~5, Ch.~9, Ch.~10).}
                $x^{3/4}y^{1/2}=c$ gives $y=c^{2}x^{-3/2}$: decreasing convex
                isoquants. $\grad f(1,1)=(\tfrac34,\tfrac12)$, norm
                $\sqrt{\tfrac{9}{16}+\tfrac{4}{16}}=\tfrac{\sqrt{13}}{4}
                \approx0.90$. The isoquant $y=x^{-3/2}$ through $(1,1)$ has
                slope $-\tfrac32$, tangent $(2,-3)$, and
                $\tfrac34\cdot2-\tfrac12\cdot3=0$ \checkmark. The gradient
                points to the upper right, towards higher output, and
                $\sqrt{13}/4$ is the output gain per unit step in that best
                direction.
          \item \emph{Critical points (Vol.~7, Ch.~5).} None in the
                interior, since $f_{x}>0$. The minimum value $0$ is attained
                on the two axes, boundary points of the domain, which the
                first-order condition does not see. (At the corner $(0,0)$
                the gradient from (b) does vanish, but the other minimizers
                on the axes are points where the partials do not even
                exist.)
          \item \emph{Contour sets.} $\{f\ge1\}=\{y\ge x^{-3/2}\}$, the
                region above a convex curve: \textbf{convex}. $\{f\le1\}$ is
                not: $(4,\tfrac18)$ and $(\tfrac14,8)$ lie on $f=1$, but the
                midpoint $(\tfrac{17}{8},\tfrac{65}{16})$ has
                $f>1$ (both coordinates exceed $1$).
          \item \emph{Checklist (Vol.~7, Ch.~3).}
                \emph{Concave: no}, $f(t,t)=t^{5/4}$ is strictly convex.
                \emph{Convex: no}, $f(t,1)=t^{3/4}$ is strictly concave.
                \emph{Quasiconcave: yes}, all upper sets lie above convex
                isoquants (Cobb--Douglas always is).
                \emph{Quasiconvex: no}, by (f).
                \emph{Maximum: no}, unbounded.
                \emph{Minimum: yes}, $0$ on the axes.
                \emph{Increasing: yes}: if $x\le x'$ and $y\le y'$ then
                $x^{3/4}\le x'^{3/4}$ and $y^{1/2}\le y'^{1/2}$, and the
                product of nonnegative factors grows.
                \emph{Decreasing: no}, $f(0,0)=0<1=f(1,1)$.
                \emph{Epigraph convex: no}, $f$ is not convex.
        \end{enumerate}
  \item \textbf{Cobb--Douglas with constant returns.}
        \begin{enumerate}[label=(\alph*),leftmargin=1.8em,itemsep=0.4ex]
          \item \emph{Name (Vol.~5, Ch.~14).} Cobb--Douglas, homogeneous of
                degree $\tfrac13+\tfrac23=1$: \textbf{constant returns to
                scale}. Domain $\R^{2}_{+}$ (closed, convex, unbounded),
                range $[0,\infty)$, since $f(t,t)=t$.
          \item \emph{Derivative and Hessian (Vol.~5, Ch.~10, Ch.~11).}
                On $x,y>0$:
                $f_{x}=\tfrac13x^{-2/3}y^{2/3}$,
                $f_{y}=\tfrac23x^{1/3}y^{-1/3}$, and
                \[
                  \Hess f=\frac29\begin{pmatrix}
                    -x^{-5/3}y^{2/3}& x^{-2/3}y^{-1/3}\\
                    x^{-2/3}y^{-1/3}& -x^{1/3}y^{-4/3}\end{pmatrix},\qquad
                  \det\Hess f=\tfrac{4}{81}x^{-4/3}y^{-2/3}
                  -\tfrac{4}{81}x^{-4/3}y^{-2/3}=0 .
                \]
                Diagonal negative, determinant zero: \textbf{negative
                semidefinite}, not definite. The null direction is the ray
                itself: the first row times $(x,y)$ is
                $\tfrac29\bigl(-x^{-2/3}y^{2/3}+x^{-2/3}y^{2/3}\bigr)=0$,
                and likewise the second. Along a ray $f(tx,ty)=tf(x,y)$ is
                linear, so it has no curvature there.
          \item \emph{Level curves, gradient, Euler (Vol.~5, Ch.~9,
                Ch.~10, Ch.~14).} $x^{1/3}y^{2/3}=c\iff y=c^{3/2}x^{-1/2}$.
                $f(1,8)=1\cdot4=4$, so the isoquant is $y=8x^{-1/2}$, slope
                $-4x^{-3/2}=-4$ at $x=1$. $\grad f(1,8)=
                (\tfrac13\cdot4,\ \tfrac23\cdot\tfrac12)=(\tfrac43,\tfrac13)$,
                norm $\tfrac{\sqrt{17}}{3}\approx1.37$; with tangent
                $(1,-4)$: $\tfrac43-\tfrac43=0$ \checkmark. Steepest ascent
                points mostly in the $x$ direction, because $x$ is the
                scarce factor at $(1,8)$. Euler: $1\cdot\tfrac43+8\cdot
                \tfrac13=4=f(1,8)$ \checkmark.
          \item \emph{Critical points.} None in the interior ($f_{x}>0$);
                the minimum $0$ sits on the axes.
          \item \emph{Contour sets.} $\{f\ge4\}=\{y\ge8x^{-1/2}\}$:
                \textbf{convex}. $\{f\le4\}$: $(16,2)$ and $(\tfrac14,16)$
                lie on $f=4$, the midpoint $(\tfrac{65}{8},9)$ gives
                $f=(\tfrac{65}{8})^{1/3}\cdot9^{2/3}\approx8.7>4$.
                \textbf{Not convex.}
          \item \emph{Checklist (Vol.~7, Ch.~3).}
                \emph{Concave: yes}, the Hessian is negative semidefinite
                on the interior, and concavity carries over to the axes by
                continuity (Cobb--Douglas with $a+b\le1$).
                \emph{Strictly concave: no}, $f(t,t)=t$ is linear along the
                diagonal.
                \emph{Convex: no}, $f(t,1)=t^{1/3}$ is strictly concave.
                \emph{Quasiconcave: yes}, because concave.
                \emph{Quasiconvex: no}, by (e).
                \emph{Maximum: no}. \emph{Minimum: yes}, $0$ on the axes.
                \emph{Increasing: yes}, as in Problem~3.
                \emph{Decreasing: no}.
                \emph{Epigraph convex: no}; the hypograph is convex.
        \end{enumerate}
  \item \textbf{Logarithm of a positive quadratic.}
        \begin{enumerate}[label=(\alph*),leftmargin=1.8em,itemsep=0.4ex]
          \item \emph{Domain and range (Vol.~2, Ch.~1).}
                $q=4+x^{2}+y^{2}\ge4>0$ always, so $D=\R^{2}$. $q$ takes
                every value in $[4,\infty)$ and $\ln$ is increasing, so
                $f(\R^{2})=[\ln4,\infty)$.
          \item \emph{Gradient (Vol.~5, Ch.~10).}
                $\grad f=\Bigl(\dfrac{2x}{q},\dfrac{2y}{q}\Bigr)$,
                continuous everywhere since $q\ne0$: differentiable on all
                of $\R^{2}$.
          \item \emph{Level curves and gradient (Vol.~5, Ch.~9).}
                $f=c\iff x^{2}+y^{2}=\mathrm{e}^{c}-4$: circles about the
                origin for $c>\ln4$, the origin itself for $c=\ln4$, empty
                below. Through $(1,1)$ runs the circle $x^{2}+y^{2}=2$
                (level $\ln6$). $\grad f(1,1)=(\tfrac26,\tfrac26)
                =(\tfrac13,\tfrac13)$, norm $\tfrac{\sqrt2}{3}$; it is a
                multiple of the radius vector $(1,1)$, and a radius is
                orthogonal to a circle: tangent $(1,-1)$,
                $\tfrac13-\tfrac13=0$ \checkmark. It points outwards,
                towards larger circles and higher values.
          \item \emph{Hessian (Vol.~5, Ch.~11).} Quotient rule:
                \[
                  f_{xx}=\frac{2(4-x^{2}+y^{2})}{q^{2}},\quad
                  f_{yy}=\frac{2(4+x^{2}-y^{2})}{q^{2}},\quad
                  f_{xy}=-\frac{4xy}{q^{2}},
                \]
                \[
                  \det\Hess f=\frac{4\bigl[16-(x^{2}-y^{2})^{2}\bigr]
                  -16x^{2}y^{2}}{q^{4}}
                  =\frac{4\bigl[16-(x^{2}+y^{2})^{2}\bigr]}{q^{4}} .
                \]
                At $(0,0)$: $\Hess f=\mathrm{diag}(\tfrac12,\tfrac12)$,
                positive definite. At $(1,1)$ ($q=6$):
                $\begin{pmatrix}2/9&-1/9\\-1/9&2/9\end{pmatrix}$,
                $\det=\tfrac{1}{27}>0$, $f_{xx}>0$: positive definite. At
                $(2,2)$ ($q=12$):
                $\begin{pmatrix}1/18&-1/9\\-1/9&1/18\end{pmatrix}$,
                $\det=-\tfrac{1}{108}<0$: indefinite.
          \item \emph{Critical points (Vol.~7, Ch.~5).} $2x/q=0$ and
                $2y/q=0$ give only $(0,0)$; positive definite Hessian: strict
                local minimum $\ln4$, and global, since $f\ge\ln4$ by (a).
          \item \emph{Contour sets.} $\{f\le\ln8\}=\{x^{2}+y^{2}\le4\}$, a
                disk: \textbf{convex}. $\{f\ge\ln8\}=\{x^{2}+y^{2}\ge4\}$,
                the outside of the disk: it contains $(2,0)$ and $(-2,0)$
                but not their midpoint $(0,0)$. \textbf{Not convex.}
          \item \emph{Checklist (Vol.~7, Ch.~3).}
                \emph{Concave: no}, the Hessian is positive definite at the
                origin.
                \emph{Convex: no}, it is indefinite at $(2,2)$; along the
                $x$-axis $g(t)=\ln(4+t^{2})$ has
                $g''(t)=\dfrac{2(4-t^{2})}{(4+t^{2})^{2}}<0$ for $|t|>2$.
                \emph{Quasiconvex: yes}, every lower set is a disk (or
                empty, or a point).
                \emph{Quasiconcave: no}, by (f).
                \emph{Maximum: no}, unbounded above.
                \emph{Minimum: yes}, $\ln4$ at $(0,0)$.
                \emph{Increasing: no}, $(-1,0)\le(0,0)$ but
                $f(-1,0)=\ln5>\ln4=f(0,0)$.
                \emph{Decreasing: no}, $(0,0)\le(1,0)$ but
                $f(1,0)=\ln5>f(0,0)$.
                \emph{Epigraph convex: no}, $f$ is not convex.
                \emph{Convex on a disk:} the determinant is positive
                exactly when $x^{2}+y^{2}<2^{2}$, and there
                $f_{xx}=2(4-x^{2}+y^{2})/q^{2}>0$ because $x^{2}<4$. So $f$
                is convex on the disk of radius $2$ (closed, by continuity),
                and on no larger disk, since the determinant turns negative
                just outside it. This is the typical shape of a
                quasiconvex but not convex function: a bowl at the bottom
                that flattens into a trumpet.
        \end{enumerate}
\end{solutions}


\end{document}
